%% GENERATED by generate_papers.py -- edit the bank_*.py files, not this file.
\documentclass[11pt,a4paper]{article}
\usepackage[a4paper, top=20mm, bottom=20mm, left=18mm, right=18mm, headheight=14pt]{geometry}
\usepackage[T1]{fontenc}
\usepackage{lmodern}
\usepackage{amsmath, amssymb, mathtools}
\usepackage{braket}
\usepackage{array, tabularx, booktabs, multicol}
\usepackage{enumitem}
\usepackage[dvipsnames]{xcolor}
\usepackage[most]{tcolorbox}
\usepackage{needspace}
\usepackage{fancyhdr, lastpage}
\usepackage{pifont}
\usepackage{tikz}
\usepackage{graphicx}
\usepackage{eso-pic}
% SILICO FELLER watermark: official logo, centred, ~8 % opacity, behind all content, on every page
\AddToShipoutPictureBG{\AtPageCenter{\makebox(0,0){\tikz\node[opacity=0.08, inner sep=0pt]{\includegraphics[width=0.72\paperwidth]{../../../assets/branding/silico_feller_logo.png}};}}}
\usepackage[hidelinks]{hyperref}
\newcolumntype{L}{>{\raggedright\arraybackslash}X}
\definecolor{brand}{HTML}{1F4E79}
\definecolor{okgreen}{HTML}{1B7F3B}
\setlength{\parindent}{0pt}
\setlength{\parskip}{3pt}
\newcommand{\optlabel}[1]{\textbf{(#1)}}
\newcommand{\correctlabel}[1]{\textcolor{okgreen}{\textbf{(#1)}\,\ding{51}}}
\newenvironment{qitem}[1]{\begin{list}{}{\setlength{\leftmargin}{3.3em}\setlength{\labelwidth}{3.1em}%
  \setlength{\labelsep}{0.2em}\setlength{\topsep}{3pt}\setlength{\parsep}{3pt}}\item[\textbf{Q.#1}\hfill]}{\end{list}}
\newcommand{\qsep}{\par\vspace{2pt}{\color{black!15}\hrule}\vspace{4pt}}
\newcommand{\qmeta}[2]{{\footnotesize\color{black!60}\textsc{#1} \;$\cdot$\; #2}\par}
\newtcolorbox{passagebox}[1]{enhanced, breakable, colback=brand!4, colframe=brand!70, fonttitle=\bfseries,
  title={#1}, left=6pt, right=6pt, top=4pt, bottom=4pt}
\newtcolorbox{sectionbanner}{enhanced, colback=brand, colframe=brand, coltext=white, fontupper=\large\bfseries,
  left=8pt, right=8pt, top=6pt, bottom=6pt, sharp corners}
\newtcolorbox{infobox}[1]{enhanced, breakable, colback=black!3, colframe=black!45, fonttitle=\bfseries,
  title={#1}, left=6pt, right=6pt, top=4pt, bottom=4pt}
\newtcolorbox{solbox}[1]{enhanced, breakable, colback=okgreen!3, colframe=okgreen!60!black, fonttitle=\bfseries,
  title={Answer: (#1) \quad Detailed solution}, left=6pt, right=6pt, top=3pt, bottom=4pt, before skip=4pt, after skip=6pt}
\newcommand{\VQEop}{\mathrm{VQE}}
\allowdisplaybreaks
\setlength{\emergencystretch}{3em}

\begin{document}
\pagestyle{fancy}\fancyhf{}
\fancyhead[L]{\small\color{black!70}Variational Quantum Eigensolver --- Assessment}
\fancyhead[R]{\small\color{black!70}\textbf{Answer Book}}
\fancyfoot[C]{\small Page \thepage\ of \pageref*{LastPage}}
\renewcommand{\headrulewidth}{0.4pt}

\thispagestyle{empty}
\begin{center}
{\color{brand}\rule{\linewidth}{1.2pt}}\\[6pt]
{\Huge\bfseries\color{brand} Variational Quantum Eigensolver}\\[4pt]
{\Large Assessment of the VQE algorithm}\\[4pt]
{\large\bfseries ANSWER BOOK}\\[2pt]
{\color{brand}\rule{\linewidth}{1.2pt}}
\end{center}
This book accompanies the separate \emph{Question Paper} (80 questions, 3 hours, 240 marks). It contains
\begin{enumerate}[leftmargin=1.4em, itemsep=0pt, topsep=2pt]
\item the \textbf{answer key}, section by section, with the topic tested by every question (page~\pageref{key}),
\item a \textbf{topic-wise map} grouping the questions by theme and a \textbf{quick-marking grid} (page~\pageref{topicmap}), and
\item the \textbf{Answer Bank}: every question restated with the correct option marked \textcolor{okgreen}{\ding{51}} and a
      complete solution with derivations (from page~\pageref{bank}).
\end{enumerate}
\begin{infobox}{Marking scheme (every question)}
Each question has \textbf{four} options (A), (B), (C) and (D). \textbf{ONLY ONE} of these four options is correct.\\[2pt]
\begin{tabular}{@{}ll@{}}
\emph{Full marks}: & $+3$ if ONLY the correct option is chosen;\\
\emph{Zero marks}: & $\phantom{+}0$ if none of the options is chosen (question unanswered);\\
\emph{Negative marks}: & $-1$ in all other cases (wrong option, or more than one option chosen).
\end{tabular}
\end{infobox}

\begin{infobox}{How to evaluate}
For each section count $N_{\mathrm c}$ = number of correct answers and $N_{\mathrm w}$ = number of questions answered wrongly (or with more than one option marked). Section marks $= 3N_{\mathrm c} - N_{\mathrm w}$; unanswered questions score $0$. The total is out of $240$.
Every numerical answer was re-computed independently (NumPy linear algebra, PySCF/OpenFermion for the chemistry, and CUDA-Q for the notebook's circuits) by
\texttt{verify\_answers.py} in the source folder.
\end{infobox}
\begin{infobox}{Conventions (identical to the Question Paper)}

\begin{itemize}[leftmargin=1.2em, itemsep=1pt, topsep=2pt]
\item $n$ qubits; a Pauli string is a word over $\{I, X, Y, Z\}$ whose character $k$ acts on qubit $k$; basis states $\ket{q_0q_1\dots q_{n-1}}$ are written
      \textbf{qubit 0 first}. $H = \sum_kc_kP_k$ with real $c_k$; $E_0$ = smallest eigenvalue; $E(\boldsymbol\theta) = \braket{\psi(\boldsymbol\theta)|H|\psi(\boldsymbol\theta)}$.
\item $R_Y(\theta) = e^{-i\theta Y/2}$, $R_X(\theta) = e^{-i\theta X/2}$, $R_Z(\theta) = e^{-i\theta Z/2}$, $S = \operatorname{diag}(1, i)$; CUDA-Q's \texttt{exp\_pauli(theta, q, word)} applies $e^{+i\theta P}$.
      Two-qubit matrices use the order $\ket{00}, \ket{01}, \ket{10}, \ket{11}$.
\item \textbf{Jordan--Wigner} (qubit $\ket1$ = occupied): $a_j^\dagger = Z_0\cdots Z_{j-1}\,\tfrac12(X_j - iY_j)$, $n_j = \tfrac12(I - Z_j)$.
\item \textbf{The notebook} (\texttt{VQE.ipynb}): default Ising chain $H = -J\sum_{i}Z_iZ_{i+1} - h\sum_iX_i$ (open, $N = 15$, $J = 1$, $h = 0.5$, 29 terms);
      hardware-efficient ansatz with $L = 2$ layers of [$R_Y$ on all qubits, CNOT ladder] plus a final $R_Y$ layer; COBYLA, 300 evaluations; exact
      energies via \texttt{cudaq.observe}. H$_2$ example: 4 qubits, 15 Pauli terms (STO-3G, Jordan--Wigner, spin-orbitals ordered $\sigma_g\alpha, \sigma_g\beta, \sigma_u\alpha, \sigma_u\beta$;
      coefficients correspond to $R = 0.740$ \AA), Hartree--Fock $\ket{1100}$, UCC state $e^{i\theta\,X_0X_1X_2Y_3}\ket{1100}$, FCI energy $-1.137284$ Ha.
      A \emph{previously recorded run} means the outputs stored in the notebook (CUDA-Q 0.16.0, CPU target \texttt{qpp-cpu}).
\item Energies of molecules in Hartree (Ha); 1 Ha $= 627.51$ kcal/mol; chemical accuracy $1.6$ mHa. Numerical answers are rounded.
\end{itemize}
\end{infobox}
\newpage
\section*{Answer key}\label{key}
{\small Answer distribution: (A) 20, (B) 20, (C) 20, (D) 20.}

\needspace{14\baselineskip}
\subsection*{Section A --- Mathematical and Quantum Chemistry Foundations}
{\small\renewcommand{\arraystretch}{1.15}
\begin{tabularx}{\linewidth}{@{}r c L @{\hspace{1em}} r c L@{}}
\toprule
\textbf{Q.} & \textbf{Ans.} & \textbf{Topic tested} & \textbf{Q.} & \textbf{Ans.} & \textbf{Topic tested}\\ \midrule
1 & \textbf{C} & $2\times2$ eigenproblem & 11 & \textbf{A} & Fermionic anticommutation \\
2 & \textbf{C} & Rayleigh quotient & 12 & \textbf{B} & $O(N^4)$ Hamiltonian terms \\
3 & \textbf{B} & Trace, determinant, spectrum & 13 & \textbf{C} & $e^{i\theta P}$ identity \\
4 & \textbf{C} & Stationary states & 14 & \textbf{D} & One-qubit $E(\theta)$ and $E_0$ \\
5 & \textbf{A} & Energy units, chemical accuracy & 15 & \textbf{D} & Optimal angle, starting error \\
6 & \textbf{D} & Born--Oppenheimer; constant term & 16 & \textbf{A} & Pauli decomposition ($2\times2$) \\
7 & \textbf{A} & Variational principle & 17 & \textbf{D} & Pauli decomposition (diagonal) \\
8 & \textbf{D} & Hermitian spectra & 18 & \textbf{B} & RHF at dissociation \\
9 & \textbf{B} & Counting determinants & 19 & \textbf{A} & Correlation energy \\
10 & \textbf{C} & LiH qubit count & 20 & \textbf{B} & Excited-state contamination \\
\bottomrule
\end{tabularx}}

\needspace{14\baselineskip}
\subsection*{Section B --- Quantum Computing Foundations}
{\small\renewcommand{\arraystretch}{1.15}
\begin{tabularx}{\linewidth}{@{}r c L @{\hspace{1em}} r c L@{}}
\toprule
\textbf{Q.} & \textbf{Ans.} & \textbf{Topic tested} & \textbf{Q.} & \textbf{Ans.} & \textbf{Topic tested}\\ \midrule
21 & \textbf{D} & $R_Y$ and the Born rule & 31 & \textbf{A} & $\texttt{XXXY}\ket{1100}$ \\
22 & \textbf{A} & Pauli products & 32 & \textbf{C} & Real expectation values \\
23 & \textbf{D} & Bloch vectors & 33 & \textbf{D} & State-vector / density-matrix memory \\
24 & \textbf{C} & Measurement basis changes & 34 & \textbf{B} & $\braket{Z_0Z_1}$ through a CNOT \\
25 & \textbf{D} & Bell-state correlations & 35 & \textbf{B} & Measuring $Y$ \\
26 & \textbf{B} & Qubit ordering in CUDA-Q & 36 & \textbf{A} & Two-qubit Ising: trial energy \\
27 & \textbf{A} & Entangling with CNOT & 37 & \textbf{B} & Two-qubit Ising: exact $E_0$ \\
28 & \textbf{C} & Normalisation; measuring $X$ & 38 & \textbf{D} & Pauli basis size \\
29 & \textbf{B} & Commuting vs.\ qubit-wise commuting & 39 & \textbf{A} & Simultaneous measurement \\
30 & \textbf{C} & \texttt{exp\_pauli} on eigenstates & 40 & \textbf{C} & Energy variance \\
\bottomrule
\end{tabularx}}

\needspace{14\baselineskip}
\subsection*{Section C --- VQE Fundamentals}
{\small\renewcommand{\arraystretch}{1.15}
\begin{tabularx}{\linewidth}{@{}r c L @{\hspace{1em}} r c L@{}}
\toprule
\textbf{Q.} & \textbf{Ans.} & \textbf{Topic tested} & \textbf{Q.} & \textbf{Ans.} & \textbf{Topic tested}\\ \midrule
41 & \textbf{C} & Energy below $E_0$ = bug & 51 & \textbf{B} & Gradient cost \\
42 & \textbf{C} & Rayleigh--Ritz bounds & 52 & \textbf{C} & COBYLA; best-of-history \\
43 & \textbf{D} & Hardware-efficient ansatz counts & 53 & \textbf{B} & Chemical accuracy check \\
44 & \textbf{C} & Hartree--Fock preparation & 54 & \textbf{A} & Measurement grouping \\
45 & \textbf{A} & Two-state energy formula & 55 & \textbf{B} & Recorded relative error \\
46 & \textbf{A} & VQE $\leftrightarrow$ code & 56 & \textbf{B} & Real ansatz for real $H$ \\
47 & \textbf{C} & Energy from Pauli averages & 57 & \textbf{D} & Symmetry breaking \\
48 & \textbf{D} & HF energy from $Z$ terms & 58 & \textbf{A} & Optimizers \\
49 & \textbf{A} & Exactness of the UCC ansatz & 59 & \textbf{D} & H$_2$ two-level reduction: $E_0$ \\
50 & \textbf{D} & Energy of $\ket{0\dots0}$ & 60 & \textbf{B} & H$_2$: optimal angle, weights \\
\bottomrule
\end{tabularx}}

\needspace{14\baselineskip}
\subsection*{Section D --- Advanced VQE and Implementation}
{\small\renewcommand{\arraystretch}{1.15}
\begin{tabularx}{\linewidth}{@{}r c L @{\hspace{1em}} r c L@{}}
\toprule
\textbf{Q.} & \textbf{Ans.} & \textbf{Topic tested} & \textbf{Q.} & \textbf{Ans.} & \textbf{Topic tested}\\ \midrule
61 & \textbf{D} & Jordan--Wigner $a_j^\dagger$ & 71 & \textbf{D} & Variational bound vs.\ shot noise \\
62 & \textbf{B} & JW hopping term & 72 & \textbf{B} & Zero-noise extrapolation \\
63 & \textbf{A} & JW vs.\ BK Pauli weight & 73 & \textbf{C} & GPU request without GPU \\
64 & \textbf{B} & Fermion encodings & 74 & \textbf{D} & Exact references \\
65 & \textbf{D} & Bravyi--Kitaev encoding & 75 & \textbf{B} & Free-fermion matrix size \\
66 & \textbf{A} & UCCSD excitation count & 76 & \textbf{D} & Recorded timing \\
67 & \textbf{C} & JW double excitation & 77 & \textbf{A} & Frozen core \\
68 & \textbf{B} & Obstacles and remedies & 78 & \textbf{C} & LiH Pauli terms \\
69 & \textbf{C} & Shots for chemical accuracy & 79 & \textbf{A} & Recorded run: recovered fraction \\
70 & \textbf{A} & Fully depolarised energy & 80 & \textbf{C} & Recorded run: next step \\
\bottomrule
\end{tabularx}}

\needspace{18\baselineskip}
\subsection*{Topic-wise map of the questions}\label{topicmap}
{\small Every question appears under at least one theme; several appear under two.\par}
{\small\renewcommand{\arraystretch}{1.2}
\begin{tabularx}{\linewidth}{@{}>{\raggedright\arraybackslash}p{0.42\linewidth}L@{}}
\toprule
\textbf{Theme} & \textbf{Question numbers}\\ \midrule
Eigenproblems and linear algebra & 1, 2, 3, 8, 13, 14, 15, 37, 42, 45, 59 \\
Variational principle and its use as a check & 7, 20, 41, 53, 71 \\
Quantum chemistry: Hamiltonians, determinants, correlation & 4, 5, 6, 9, 10, 12, 18, 19, 77, 78 \\
Second quantisation and fermion-to-qubit mappings & 11, 61, 62, 63, 64, 65, 67 \\
Pauli decomposition, Pauli algebra and measurement & 16, 17, 22, 24, 25, 29, 31, 35, 38, 39, 47, 54 \\
Qubits, gates and circuits & 21, 23, 26, 27, 28, 30, 32, 34, 36 \\
Ansatz design: hardware-efficient and UCC & 43, 44, 49, 56, 60, 66, 67 \\
Classical optimisation and convergence & 51, 52, 58, 68, 80 \\
Shots, variance and noise & 40, 69, 70, 71, 72 \\
Simulation backends, memory and exact references & 33, 73, 74, 75, 76 \\
The notebook's code and recorded results & 46, 48, 50, 55, 57, 76, 79, 80 \\
\bottomrule
\end{tabularx}}

\subsection*{Quick-marking grid}
\begin{center}\renewcommand{\arraystretch}{1.3}
\begin{tabular}{|c|c|c|c|c|c|c|c|c|c|}\hline
1\,\textbf{C} & 2\,\textbf{C} & 3\,\textbf{B} & 4\,\textbf{C} & 5\,\textbf{A} & 6\,\textbf{D} & 7\,\textbf{A} & 8\,\textbf{D} & 9\,\textbf{B} & 10\,\textbf{C} \\
\hline
11\,\textbf{A} & 12\,\textbf{B} & 13\,\textbf{C} & 14\,\textbf{D} & 15\,\textbf{D} & 16\,\textbf{A} & 17\,\textbf{D} & 18\,\textbf{B} & 19\,\textbf{A} & 20\,\textbf{B} \\
\hline
21\,\textbf{D} & 22\,\textbf{A} & 23\,\textbf{D} & 24\,\textbf{C} & 25\,\textbf{D} & 26\,\textbf{B} & 27\,\textbf{A} & 28\,\textbf{C} & 29\,\textbf{B} & 30\,\textbf{C} \\
\hline
31\,\textbf{A} & 32\,\textbf{C} & 33\,\textbf{D} & 34\,\textbf{B} & 35\,\textbf{B} & 36\,\textbf{A} & 37\,\textbf{B} & 38\,\textbf{D} & 39\,\textbf{A} & 40\,\textbf{C} \\
\hline
41\,\textbf{C} & 42\,\textbf{C} & 43\,\textbf{D} & 44\,\textbf{C} & 45\,\textbf{A} & 46\,\textbf{A} & 47\,\textbf{C} & 48\,\textbf{D} & 49\,\textbf{A} & 50\,\textbf{D} \\
\hline
51\,\textbf{B} & 52\,\textbf{C} & 53\,\textbf{B} & 54\,\textbf{A} & 55\,\textbf{B} & 56\,\textbf{B} & 57\,\textbf{D} & 58\,\textbf{A} & 59\,\textbf{D} & 60\,\textbf{B} \\
\hline
61\,\textbf{D} & 62\,\textbf{B} & 63\,\textbf{A} & 64\,\textbf{B} & 65\,\textbf{D} & 66\,\textbf{A} & 67\,\textbf{C} & 68\,\textbf{B} & 69\,\textbf{C} & 70\,\textbf{A} \\
\hline
71\,\textbf{D} & 72\,\textbf{B} & 73\,\textbf{C} & 74\,\textbf{D} & 75\,\textbf{B} & 76\,\textbf{D} & 77\,\textbf{A} & 78\,\textbf{C} & 79\,\textbf{A} & 80\,\textbf{C} \\
\hline\end{tabular}\end{center}
\newpage
\section*{Answer Bank --- questions with detailed solutions}\label{bank}

\begin{sectionbanner}SECTION A \enspace$\cdot$\enspace Mathematical and Quantum Chemistry Foundations \hfill Q.1--20\end{sectionbanner}

\needspace{13\baselineskip}
\begin{qitem}{1}
\qmeta{Single-correct MCQ}{$2\times2$ eigenproblem}
The ground-state energy (smallest eigenvalue) of $H = \begin{pmatrix}1 & 2\\ 2 & -2\end{pmatrix}$ and a normalised eigenvector for it are
\par\smallskip\noindent\begin{tabularx}{\linewidth}{@{}*{2}{L}@{}}
\optlabel{A}\,$-2$ and $(0, 1)^{\mathsf T}$ & \optlabel{B}\,$2$ and $\tfrac{1}{\sqrt5}(2, 1)^{\mathsf T}$ \\[3pt]
\correctlabel{C}\,$-3$ and $\tfrac{1}{\sqrt5}(1, -2)^{\mathsf T}$ & \optlabel{D}\,$-3$ and $\tfrac{1}{\sqrt5}(2, 1)^{\mathsf T}$
\end{tabularx}
\end{qitem}
\begin{solbox}{C}
\textbf{\color{brand}Correct answer.} \textcolor{okgreen}{\textbf{(C)}} $-3$ and $\tfrac{1}{\sqrt5}(1, -2)^{\mathsf T}$

\textbf{\color{brand}Concepts required.} Eigenvalues from the characteristic polynomial; eigenvectors from the null space of $H - \lambda I$.

\textbf{\color{brand}Relevant formulae.} $\det(H - \lambda I) = 0$; $(H - \lambda I)v = 0$; normalise $v/\|v\|$.

\textbf{\color{brand}Step-by-step explanation.} Characteristic polynomial: $\det(H - \lambda I) = (1 - \lambda)(-2 - \lambda) - 4 = \lambda^2 + \lambda - 6 = (\lambda + 3)(\lambda - 2)$,
so $\lambda \in \{-3, 2\}$ and $E_0 = -3$.
Eigenvector for $-3$: $(H + 3I)v = 0$ gives $4v_1 + 2v_2 = 0$, i.e.\ $v_2 = -2v_1$: $v = (1, -2)/\sqrt5$.
Check: $Hv = (1 - 4, 2 + 4)/\sqrt5 = (-3, 6)/\sqrt5 = -3v$. The vector $(2, 1)/\sqrt5$ is orthogonal to it and belongs to $\lambda = 2$.

\textbf{\color{brand}Why the other options are incorrect.}
\begin{itemize}[leftmargin=2.4em, itemsep=1pt, topsep=2pt]
\item[\textbf{(A)}] $-2$ is the smallest \emph{diagonal} entry; off-diagonal coupling pushes the true minimum lower, to $-3$.
\item[\textbf{(B)}] $2$ is the largest eigenvalue, not the ground state.
\item[\textbf{(D)}] $(2, 1)/\sqrt5$ is the eigenvector of the \emph{other} eigenvalue $2$: $H(2,1) = (4, 2) = 2(2, 1)$.
\end{itemize}
\textbf{\color{brand}Key takeaway.} The ground state lies below the smallest diagonal element whenever off-diagonal couplings exist.
\end{solbox}
\qsep

\needspace{13\baselineskip}
\begin{qitem}{2}
\qmeta{Single-correct MCQ}{Rayleigh quotient}
For the same matrix $H = \begin{pmatrix}1 & 2\\ 2 & -2\end{pmatrix}$, the Rayleigh quotient $R(v) = \dfrac{v^{\mathsf T}Hv}{v^{\mathsf T}v}$ of the trial
vector $v = (1, -1)^{\mathsf T}$ is
\par\smallskip\noindent\begin{tabularx}{\linewidth}{@{}*{2}{L}@{}}
\optlabel{A}\,$-5$, which lies below $E_0$ & \optlabel{B}\,$-3$, exactly $E_0$ \\[3pt]
\correctlabel{C}\,$-2.5$, which lies above $E_0 = -3$ & \optlabel{D}\,$0.5$
\end{tabularx}
\end{qitem}
\begin{solbox}{C}
\textbf{\color{brand}Correct answer.} \textcolor{okgreen}{\textbf{(C)}} $-2.5$, which lies above $E_0 = -3$

\textbf{\color{brand}Concepts required.} Rayleigh quotient and its bounds.

\textbf{\color{brand}Relevant formulae.} $R(v) = \frac{v^{\mathsf T}Hv}{v^{\mathsf T}v}$, $\ \lambda_{\min} \le R(v) \le \lambda_{\max}$.

\textbf{\color{brand}Step-by-step explanation.} $Hv = (1 - 2,\ 2 + 2) = (-1, 4)$, so $v^{\mathsf T}Hv = 1\cdot(-1) + (-1)\cdot4 = -5$, and $v^{\mathsf T}v = 2$: $R(v) = -5/2 = -2.5$.
The Rayleigh quotient of any non-zero vector lies between the smallest and largest eigenvalue: $-3 \le -2.5 \le 2$.
Eigenbasis check: with $u_0 = (1, -2)/\sqrt5$ (eigenvalue $-3$), the weight of $u_0$ in $v/\|v\|$ is $\langle v, u_0\rangle^2/\|v\|^2 = (3/\sqrt5)^2/2 = 0.9$;
the remaining $0.1$ belongs to eigenvalue $2$: $R = 0.9(-3) + 0.1(2) = -2.5$. This is the variational principle in its simplest form.

\textbf{\color{brand}Why the other options are incorrect.}
\begin{itemize}[leftmargin=2.4em, itemsep=1pt, topsep=2pt]
\item[\textbf{(A)}] $-5$ is $v^{\mathsf T}Hv$ without dividing by $v^{\mathsf T}v = 2$; no Rayleigh quotient can lie below $\lambda_{\min} = -3$.
\item[\textbf{(B)}] $-3$ would require $v \propto (1, -2)$; $(1, -1)$ is not an eigenvector.
\item[\textbf{(D)}] $0.5$ comes from sign errors in $Hv$; the correct $Hv = (-1, 4)$.
\end{itemize}
\textbf{\color{brand}Key takeaway.} Any trial vector's energy is an upper bound to the ground energy --- VQE in one line.
\end{solbox}
\qsep

\needspace{13\baselineskip}
\begin{qitem}{3}
\qmeta{Matching lists}{Trace, determinant, spectrum}
For the real symmetric matrix $M = \begin{pmatrix}2 & 1 & 0\\ 1 & 2 & 0\\ 0 & 0 & 5\end{pmatrix}$, match each quantity (List-I) with its value (List-II).
\par\smallskip\noindent{\renewcommand{\arraystretch}{1.35}%
\begin{tabularx}{\linewidth}{@{}c>{\raggedright\arraybackslash}p{0.33\linewidth}@{\hspace{1.2em}}cL@{}}
\toprule
\multicolumn{2}{@{}l}{\textbf{List-I}} & \multicolumn{2}{l@{}}{\textbf{List-II}} \\ \midrule
\textbf{(P)} & $\operatorname{tr} M$ & \textbf{(1)} & $15$ \\
\textbf{(Q)} & $\det M$ & \textbf{(2)} & $1$ \\
\textbf{(R)} & smallest eigenvalue of $M$ & \textbf{(3)} & $9$ \\
\textbf{(S)} & largest eigenvalue of $M$ & \textbf{(4)} & $5$ \\
\bottomrule
\end{tabularx}}
\par\smallskip\noindent\begin{tabularx}{\linewidth}{@{}*{2}{L}@{}}
\optlabel{A}\,P$\to$3, Q$\to$1, R$\to$4, S$\to$2 & \correctlabel{B}\,P$\to$3, Q$\to$1, R$\to$2, S$\to$4 \\[3pt]
\optlabel{C}\,P$\to$1, Q$\to$3, R$\to$2, S$\to$4 & \optlabel{D}\,P$\to$3, Q$\to$4, R$\to$2, S$\to$1
\end{tabularx}
\end{qitem}
\begin{solbox}{B}
\textbf{\color{brand}Correct answer.} \textcolor{okgreen}{\textbf{(B)}} P$\to$3, Q$\to$1, R$\to$2, S$\to$4

\textbf{\color{brand}Concepts required.} Trace = sum, determinant = product of eigenvalues; block-diagonal matrices.

\textbf{\color{brand}Relevant formulae.} $\operatorname{tr}M = \sum\lambda_i$, $\det M = \prod\lambda_i$; $\begin{psmallmatrix}a&b\\b&a\end{psmallmatrix}$ has eigenvalues $a \pm b$.

\textbf{\color{brand}Step-by-step explanation.} $M$ is block diagonal: the block $\begin{psmallmatrix}2&1\\1&2\end{psmallmatrix}$ has eigenvalues $2 \pm 1 = 1, 3$ (eigenvectors $(1, \mp1)/\sqrt2$), and the
block $(5)$ has eigenvalue $5$. Spectrum $\{1, 3, 5\}$.
\begin{itemize}\itemsep0pt
\item P: trace = sum of diagonal entries $= 2 + 2 + 5 = 9$ = sum of eigenvalues (3).
\item Q: determinant = product of eigenvalues $= 1\cdot3\cdot5 = 15$ (also $(4 - 1)\cdot5$) (1).
\item R: smallest eigenvalue $1$ (2) --- the ``ground state'' of $M$.
\item S: largest eigenvalue $5$ (4).
\end{itemize}

\textbf{\color{brand}Why the other options are incorrect.}
\begin{itemize}[leftmargin=2.4em, itemsep=1pt, topsep=2pt]
\item[\textbf{(A)}] R$\to$4, S$\to$2 swaps smallest and largest eigenvalues.
\item[\textbf{(C)}] P$\to$1, Q$\to$3 swaps trace and determinant: $9 = 2+2+5$, $15 = 1\cdot3\cdot5$.
\item[\textbf{(D)}] Q$\to$4, S$\to$1 gives the determinant 5 and the largest eigenvalue 15; the eigenvalues are only $1, 3, 5$.
\end{itemize}
\textbf{\color{brand}Key takeaway.} Check eigenvalues against trace and determinant.
\end{solbox}
\qsep

\needspace{13\baselineskip}
\begin{qitem}{4}
\qmeta{Single-correct MCQ}{Stationary states}
An energy eigenstate $\ket{\psi_E}$ of a time-independent Hamiltonian ($H\ket{\psi_E} = E\ket{\psi_E}$) evolves under the time-dependent
Schr\"odinger equation $i\hbar\,\partial_t\ket{\psi(t)} = H\ket{\psi(t)}$. Which statement is correct?
\begin{itemize}[leftmargin=2.6em, labelsep=0.4em, itemsep=2pt, topsep=3pt]
\item[\optlabel{A}] $\ket{\psi(t)} = e^{-Et/\hbar}\ket{\psi_E}$: the state decays in time
\item[\optlabel{B}] The state is constant only if $E = 0$; otherwise it moves to other eigenstates
\item[\correctlabel{C}] $\ket{\psi(t)} = e^{-iEt/\hbar}\ket{\psi_E}$: only a global phase changes, so all measurement probabilities are constant in time
\item[\optlabel{D}] The probabilities oscillate with angular frequency $E/\hbar$
\end{itemize}
\end{qitem}
\begin{solbox}{C}
\textbf{\color{brand}Correct answer.} \textcolor{okgreen}{\textbf{(C)}} $\ket{\psi(t)} = e^{-iEt/\hbar}\ket{\psi_E}$: only a global phase changes, so all measurement probabilities are constant in time

\textbf{\color{brand}Concepts required.} Stationary states; global versus relative phase.

\textbf{\color{brand}Relevant formulae.} $\ket{\psi(t)} = e^{-iHt/\hbar}\ket{\psi(0)}$; on an eigenvector $e^{-iHt/\hbar}\ket{\psi_E} = e^{-iEt/\hbar}\ket{\psi_E}$.

\textbf{\color{brand}Step-by-step explanation.} The formal solution is $\ket{\psi(t)} = e^{-iHt/\hbar}\ket{\psi(0)}$. On an eigenvector $H$ acts as the number $E$, so
$e^{-iHt/\hbar}\ket{\psi_E} = e^{-iEt/\hbar}\ket{\psi_E}$ (expand the exponential as a power series). The factor has modulus 1: a
\emph{global} phase, invisible in every measurement. Hence eigenstates are called \emph{stationary states}, and finding them
reduces to the time-independent equation $H\ket\psi = E\ket\psi$ --- the eigenvalue problem VQE solves. A superposition of
eigenstates with different energies, by contrast, acquires \emph{relative} phases and its probabilities can oscillate.

\textbf{\color{brand}Why the other options are incorrect.}
\begin{itemize}[leftmargin=2.4em, itemsep=1pt, topsep=2pt]
\item[\textbf{(A)}] $e^{-Et/\hbar}$ lacks the $i$: Schr\"odinger evolution is unitary, norms are conserved.
\item[\textbf{(B)}] Eigenstates stay eigenstates for any $E$; only superpositions of different energies evolve non-trivially.
\item[\textbf{(D)}] A global phase has no observable effect; probabilities of an eigenstate do not oscillate.
\end{itemize}
\textbf{\color{brand}Key takeaway.} Energy eigenstates are stationary: finding them is the time-independent problem $H\ket\psi = E\ket\psi$.
\end{solbox}
\qsep

\needspace{13\baselineskip}
\begin{qitem}{5}
\qmeta{Single-correct MCQ}{Energy units, chemical accuracy}
Using $1\ \text{Hartree} = 627.51$ kcal/mol, an energy error of $0.005$ Hartree corresponds to about ... and is ... chemically
accurate (chemical accuracy: 1 kcal/mol $\approx$ 1.6 mHa).
\par\smallskip\noindent\begin{tabularx}{\linewidth}{@{}*{2}{L}@{}}
\correctlabel{A}\,$3.14$ kcal/mol; not & \optlabel{B}\,$3.14$ kcal/mol; still \\[3pt]
\optlabel{C}\,$0.31$ kcal/mol; still & \optlabel{D}\,$125.5$ kcal/mol; not
\end{tabularx}
\end{qitem}
\begin{solbox}{A}
\textbf{\color{brand}Correct answer.} \textcolor{okgreen}{\textbf{(A)}} $3.14$ kcal/mol; not

\textbf{\color{brand}Concepts required.} Energy units; chemical accuracy.

\textbf{\color{brand}Relevant formulae.} $1$ Ha $= 627.51$ kcal/mol $= 27.211$ eV; chemical accuracy $1$ kcal/mol $= 1.594$ mHa.

\textbf{\color{brand}Step-by-step explanation.} $0.005 \times 627.51 = 3.138$ kcal/mol. Chemical accuracy means an error below $1$ kcal/mol $= 1/627.51$ Ha $= 1.594$ mHa
($\approx 1.6$ mHa, the threshold printed by the notebook). $5$ mHa is about three times too large, so the result is \emph{not}
chemically accurate. The threshold matters because reaction energies and barriers decide chemistry at the kcal/mol scale.

\textbf{\color{brand}Why the other options are incorrect.}
\begin{itemize}[leftmargin=2.4em, itemsep=1pt, topsep=2pt]
\item[\textbf{(B)}] $3.14$ kcal/mol exceeds the 1 kcal/mol threshold.
\item[\textbf{(C)}] $0.31$ multiplies by $62.751$ instead of $627.51$.
\item[\textbf{(D)}] $125.5$ multiplies by $25\,102$ (mixing units); $0.005\times627.51 = 3.14$.
\end{itemize}
\textbf{\color{brand}Key takeaway.} $1.6$ mHa $\approx 1$ kcal/mol is the target for useful chemistry.
\end{solbox}
\qsep

\needspace{13\baselineskip}
\begin{qitem}{6}
\qmeta{Single-correct MCQ}{Born--Oppenheimer; constant term}
In the Born--Oppenheimer approximation the nuclei are held fixed and only the electrons are treated quantum mechanically. In the
notebook's 4-qubit H$_2$ Hamiltonian, the nuclear--nuclear repulsion energy appears
\begin{itemize}[leftmargin=2.6em, labelsep=0.4em, itemsep=2pt, topsep=3pt]
\item[\optlabel{A}] nowhere, because nuclear repulsion is neglected in quantum chemistry
\item[\optlabel{B}] spread over the four $\texttt{XXYY}$-type exchange terms
\item[\optlabel{C}] as the coefficient of the $\texttt{ZZII}$ term
\item[\correctlabel{D}] inside the coefficient of the identity string $\texttt{IIII}$, as a constant energy shift
\end{itemize}
\end{qitem}
\begin{solbox}{D}
\textbf{\color{brand}Correct answer.} \textcolor{okgreen}{\textbf{(D)}} inside the coefficient of the identity string $\texttt{IIII}$, as a constant energy shift

\textbf{\color{brand}Concepts required.} Born--Oppenheimer approximation; constants in a qubit Hamiltonian.

\textbf{\color{brand}Relevant formulae.} $V_{NN} = \sum_{A<B}\frac{Z_AZ_B}{R_{AB}}$ (atomic units); a constant $c$ maps to $c\,I^{\otimes n}$.

\textbf{\color{brand}Step-by-step explanation.} The molecular Hamiltonian is $H = T_e + V_{eN} + V_{ee} + V_{NN}$. With fixed nuclei (Born--Oppenheimer), $V_{NN} = \frac{Z_AZ_B}{R}$ is
just a number for a given bond length $R$. A number multiplies the identity operator, so after the fermion-to-qubit mapping it
ends up in the coefficient of $\texttt{IIII}$ (together with constant pieces of the electronic terms). The notebook's comment
says so: ``the constant term includes the nuclear repulsion energy''. It shifts every eigenvalue equally and needs no
measurement, but it must be included to compare with total energies such as the FCI value $-1.137284$ Ha.

\textbf{\color{brand}Why the other options are incorrect.}
\begin{itemize}[leftmargin=2.4em, itemsep=1pt, topsep=2pt]
\item[\textbf{(A)}] $V_{NN}$ is included; otherwise total energies like $-1.137$ Ha could not be compared with experiment or FCI.
\item[\textbf{(B)}] The exchange terms describe the double excitation $\ket{1100} \leftrightarrow \ket{0011}$, not a constant.
\item[\textbf{(C)}] $\texttt{ZZII}$ is an electron--electron interaction between spin-orbitals 0 and 1, not the nuclear repulsion.
\end{itemize}
\textbf{\color{brand}Key takeaway.} Constants go into the identity coefficient: no measurement needed, but never forget them.
\end{solbox}
\qsep

\needspace{13\baselineskip}
\begin{qitem}{7}
\qmeta{Assertion--Reason}{Variational principle}
\textbf{Assertion (A):} For every normalised state $\ket\psi$, $\braket{\psi|H|\psi} \ge E_0$, where $E_0$ is the smallest eigenvalue of the
Hermitian operator $H$.\\[1mm]
\textbf{Reason (R):} Expanding $\ket\psi = \sum_k c_k\ket{E_k}$ in the orthonormal eigenbasis of $H$ gives $\braket{\psi|H|\psi} = \sum_k |c_k|^2E_k$ with
$\sum_k|c_k|^2 = 1$.
\begin{itemize}[leftmargin=2.6em, labelsep=0.4em, itemsep=2pt, topsep=3pt]
\item[\correctlabel{A}] Both Assertion (A) and Reason (R) are true, and R is the correct explanation of A
\item[\optlabel{B}] Both A and R are true, but R is \emph{not} the correct explanation of A
\item[\optlabel{C}] A is true, but R is false
\item[\optlabel{D}] A is false, but R is true
\end{itemize}
\end{qitem}
\begin{solbox}{A}
\textbf{\color{brand}Correct answer.} \textcolor{okgreen}{\textbf{(A)}} Both Assertion (A) and Reason (R) are true, and R is the correct explanation of A

\textbf{\color{brand}Concepts required.} Spectral theorem; variational principle.

\textbf{\color{brand}Relevant formulae.} $\braket{\psi|H|\psi} = \sum_k|c_k|^2E_k \ge E_0$, $\sum|c_k|^2 = 1$.

\textbf{\color{brand}Step-by-step explanation.} \textbf{R is true:} a Hermitian $H$ has an orthonormal eigenbasis (spectral theorem). Then
$\braket{\psi|H|\psi} = \sum_{k,l} c_k^*c_l\braket{E_k|H|E_l} = \sum_{k,l}c_k^*c_lE_l\delta_{kl} = \sum_k|c_k|^2E_k$, and normalisation gives $\sum|c_k|^2 = 1$.
\textbf{A follows from R:} a weighted average with non-negative weights summing to 1 cannot be below the smallest value:
$\sum_k|c_k|^2E_k \ge \sum_k|c_k|^2E_0 = E_0$; equality iff all weight sits on ground states. R is the correct explanation. Answer: (A).

\textbf{\color{brand}Why the other options are incorrect.}
\begin{itemize}[leftmargin=2.4em, itemsep=1pt, topsep=2pt]
\item[\textbf{(B)}] needs A and R both true but R not explaining A; R is exactly the reason.
\item[\textbf{(C)}] needs A true and R false; R is true.
\item[\textbf{(D)}] needs A false and R true; A is true.
\end{itemize}
\textbf{\color{brand}Key takeaway.} Energy expectation = weighted average of eigenvalues $\Rightarrow$ never below $E_0$.
\end{solbox}
\qsep

\needspace{13\baselineskip}
\begin{qitem}{8}
\qmeta{Statement I / II}{Hermitian spectra}
\textbf{Statement I:} A real symmetric matrix can have complex (non-real) eigenvalues.\\[1mm]
\textbf{Statement II:} Eigenvectors of a Hermitian matrix that belong to different eigenvalues are orthogonal.
\begin{itemize}[leftmargin=2.6em, labelsep=0.4em, itemsep=2pt, topsep=3pt]
\item[\optlabel{A}] Both Statement I and Statement II are true
\item[\optlabel{B}] Both Statement I and Statement II are false
\item[\optlabel{C}] Statement I is true, but Statement II is false
\item[\correctlabel{D}] Statement I is false, but Statement II is true
\end{itemize}
\end{qitem}
\begin{solbox}{D}
\textbf{\color{brand}Correct answer.} \textcolor{okgreen}{\textbf{(D)}} Statement I is false, but Statement II is true

\textbf{\color{brand}Concepts required.} Spectral properties of Hermitian matrices.

\textbf{\color{brand}Relevant formulae.} $\lambda = v^\dagger Hv = \lambda^*$; $(\mu - \lambda)u^\dagger v = 0$.

\textbf{\color{brand}Step-by-step explanation.} \textbf{I is false:} real symmetric matrices are Hermitian; if $Hv = \lambda v$ with $\|v\| = 1$ then
$\lambda = v^\dagger Hv = (v^\dagger Hv)^* = \lambda^*$, so $\lambda$ is real. (Complex eigenvalues occur for non-symmetric real matrices such as rotations.)
\textbf{II is true:} from $Hu = \lambda u$, $Hv = \mu v$: $\mu\,u^\dagger v = u^\dagger Hv = (Hu)^\dagger v = \lambda u^\dagger v$, so $(\mu - \lambda)u^\dagger v = 0$ and
$u^\dagger v = 0$ when $\lambda \neq \mu$. Real energies and orthogonal eigenstates are exactly what the variational argument uses. Answer: (D).

\textbf{\color{brand}Why the other options are incorrect.}
\begin{itemize}[leftmargin=2.4em, itemsep=1pt, topsep=2pt]
\item[\textbf{(A)}] needs both statements true; I is false.
\item[\textbf{(B)}] needs both statements false; II is true.
\item[\textbf{(C)}] needs I true and II false; I is false and II is true.
\end{itemize}
\textbf{\color{brand}Key takeaway.} Hermitian $\Rightarrow$ real eigenvalues and orthogonal eigenvectors.
\end{solbox}
\qsep

\needspace{13\baselineskip}
\begin{qitem}{9}
\qmeta{Single-correct MCQ}{Counting determinants}
H$_2$ in the STO-3G basis has 2 spatial orbitals, i.e.\ 4 spin-orbitals, and 2 electrons. The number of 2-electron Slater determinants,
and the number of them with one spin-up and one spin-down electron ($S_z = 0$), are
\par\smallskip\noindent\begin{tabularx}{\linewidth}{@{}*{4}{L}@{}}
\optlabel{A}\,$6$ and $2$ & \correctlabel{B}\,$6$ and $4$ & \optlabel{C}\,$16$ and $4$ & \optlabel{D}\,$4$ and $2$
\end{tabularx}
\end{qitem}
\begin{solbox}{B}
\textbf{\color{brand}Correct answer.} \textcolor{okgreen}{\textbf{(B)}} $6$ and $4$

\textbf{\color{brand}Concepts required.} Slater determinants as occupation patterns; spin sectors.

\textbf{\color{brand}Relevant formulae.} $\#\text{det} = \binom{N_{\mathrm{so}}}{N_e}$; $S_z = 0$: $\binom{N_{\mathrm{orb}}}{N_\alpha}\binom{N_{\mathrm{orb}}}{N_\beta}$.

\textbf{\color{brand}Step-by-step explanation.} A Slater determinant is fixed by choosing which spin-orbitals are occupied: $\binom{4}{2} = 6$ ways (the 4-qubit Fock space has $2^4 = 16$
basis states for \emph{all} electron numbers). With one $\alpha$ electron (2 choices of spatial orbital) and one $\beta$ electron (2 choices):
$2\times2 = 4$ determinants with $S_z = 0$. Exact diagonalisation within such a determinant space is \emph{full configuration
interaction} (FCI). The notebook's reference $-1.137284$ Ha is the FCI energy; its ground state uses only two of these
determinants, $\ket{1100}$ and $\ket{0011}$ (see Q.59).

\textbf{\color{brand}Why the other options are incorrect.}
\begin{itemize}[leftmargin=2.4em, itemsep=1pt, topsep=2pt]
\item[\textbf{(A)}] $2$ counts only the closed-shell determinants ($\ket{1100}$, $\ket{0011}$).
\item[\textbf{(C)}] $16 = 2^4$ counts all occupation patterns for any electron number.
\item[\textbf{(D)}] $4$ would be $\binom41$ (one electron).
\end{itemize}
\textbf{\color{brand}Key takeaway.} FCI = exact diagonalisation within the fixed-electron-number determinant space.
\end{solbox}
\qsep

\needspace{13\baselineskip}
\begin{qitem}{10}
\qmeta{Single-correct MCQ}{LiH qubit count}
LiH in the minimal STO-3G basis has one $1s$ orbital on H and the orbitals $1s, 2s, 2p_x, 2p_y, 2p_z$ on Li. With the Jordan--Wigner
mapping (one qubit per spin-orbital, no reductions), the number of qubits and the number of electrons are
\par\smallskip\noindent\begin{tabularx}{\linewidth}{@{}*{2}{L}@{}}
\optlabel{A}\,$24$ qubits, $4$ electrons & \optlabel{B}\,$12$ qubits, $3$ electrons \\[3pt]
\correctlabel{C}\,$12$ qubits, $4$ electrons & \optlabel{D}\,$6$ qubits, $4$ electrons
\end{tabularx}
\end{qitem}
\begin{solbox}{C}
\textbf{\color{brand}Correct answer.} \textcolor{okgreen}{\textbf{(C)}} $12$ qubits, $4$ electrons

\textbf{\color{brand}Concepts required.} Minimal basis; spin-orbitals; Jordan--Wigner qubit count.

\textbf{\color{brand}Relevant formulae.} qubits $= 2\times$(spatial orbitals); electrons $=\sum Z_{\mathrm{atom}}$ (neutral).

\textbf{\color{brand}Step-by-step explanation.} Spatial orbitals: $5$ (Li) $+ 1$ (H) $= 6$. Each holds a spin-up and a spin-down spin-orbital: $12$ spin-orbitals $\Rightarrow$ $12$
qubits under Jordan--Wigner. Electrons: Li has 3, H has 1: $4$. (Checked with PySCF in \texttt{verify\_answers.py}.) Freezing the Li $1s$
core and using symmetries reduces the qubit count further (Q.77). LiH is not part of the notebook; it is a standard
illustration of how molecular problems grow.

\textbf{\color{brand}Why the other options are incorrect.}
\begin{itemize}[leftmargin=2.4em, itemsep=1pt, topsep=2pt]
\item[\textbf{(A)}] $24$ double-counts spin.
\item[\textbf{(B)}] Li has 3 electrons and H one: 4, not 3.
\item[\textbf{(D)}] $6$ counts spatial orbitals; each needs two qubits (spin up and down).
\end{itemize}
\textbf{\color{brand}Key takeaway.} Count spatial orbitals, double for spin: LiH/STO-3G $\to$ 12 qubits.
\end{solbox}
\qsep

\needspace{13\baselineskip}
\begin{qitem}{11}
\qmeta{Matching lists}{Fermionic anticommutation}
Fermionic creation and annihilation operators $a_p^\dagger$, $a_p$ satisfy the canonical anticommutation relations. Match each expression
(List-I) with its value (List-II); $\ket{\mathrm{vac}}$ is the empty (vacuum) state.
\par\smallskip\noindent{\renewcommand{\arraystretch}{1.35}%
\begin{tabularx}{\linewidth}{@{}c>{\raggedright\arraybackslash}p{0.33\linewidth}@{\hspace{1.2em}}cL@{}}
\toprule
\multicolumn{2}{@{}l}{\textbf{List-I}} & \multicolumn{2}{l@{}}{\textbf{List-II}} \\ \midrule
\textbf{(P)} & $\{a_p, a_q^\dagger\} = a_pa_q^\dagger + a_q^\dagger a_p$ & \textbf{(1)} & $\{0, 1\}$ \\
\textbf{(Q)} & $a_p^\dagger a_p^\dagger$ & \textbf{(2)} & $\delta_{pq}\,I$ \\
\textbf{(R)} & spectrum of $n_p = a_p^\dagger a_p$ & \textbf{(3)} & $0$ (the zero vector) \\
\textbf{(S)} & $a_p\ket{\mathrm{vac}}$ & \textbf{(4)} & $0$ (the zero operator) \\
\bottomrule
\end{tabularx}}
\par\smallskip\noindent\begin{tabularx}{\linewidth}{@{}*{2}{L}@{}}
\correctlabel{A}\,P$\to$2, Q$\to$4, R$\to$1, S$\to$3 & \optlabel{B}\,P$\to$2, Q$\to$3, R$\to$1, S$\to$4 \\[3pt]
\optlabel{C}\,P$\to$2, Q$\to$4, R$\to$3, S$\to$1 & \optlabel{D}\,P$\to$4, Q$\to$2, R$\to$1, S$\to$3
\end{tabularx}
\end{qitem}
\begin{solbox}{A}
\textbf{\color{brand}Correct answer.} \textcolor{okgreen}{\textbf{(A)}} P$\to$2, Q$\to$4, R$\to$1, S$\to$3

\textbf{\color{brand}Concepts required.} Canonical anticommutation relations and their consequences.

\textbf{\color{brand}Relevant formulae.} $\{a_p, a_q^\dagger\} = \delta_{pq}$, $\{a_p, a_q\} = 0$, $n_p^2 = n_p$.

\textbf{\color{brand}Step-by-step explanation.} \begin{itemize}\itemsep0pt
\item P: canonical anticommutation relation $\{a_p, a_q^\dagger\} = \delta_{pq}I$ (2).
\item Q: $\{a_p^\dagger, a_p^\dagger\} = 2(a_p^\dagger)^2 = 0$, so $(a_p^\dagger)^2 = 0$ as an \emph{operator}: no orbital can be filled twice --- Pauli exclusion (4).
\item R: $n_p^2 = a_p^\dagger(a_pa_p^\dagger)a_p = a_p^\dagger(1 - a_p^\dagger a_p)a_p = n_p$, so $n_p$ is a projector with eigenvalues $0$ and $1$ (empty/occupied) (1).
\item S: there is nothing to remove from the vacuum: $a_p\ket{\mathrm{vac}}$ is the zero \emph{vector} (3).
\end{itemize}
These relations are what the Jordan--Wigner $Z$-strings must reproduce on qubits.

\textbf{\color{brand}Why the other options are incorrect.}
\begin{itemize}[leftmargin=2.4em, itemsep=1pt, topsep=2pt]
\item[\textbf{(B)}] Q$\to$3, S$\to$4 confuses the zero operator $(a_p^\dagger)^2$ with the zero vector $a_p\ket{\mathrm{vac}}$.
\item[\textbf{(C)}] R$\to$3, S$\to$1: the number operator has eigenvalues $\{0, 1\}$; annihilating the vacuum gives the zero vector.
\item[\textbf{(D)}] P$\to$4, Q$\to$2 swaps the CAR value with Pauli exclusion.
\end{itemize}
\textbf{\color{brand}Key takeaway.} Fermions anticommute; qubit mappings must reproduce these signs.
\end{solbox}
\qsep

\needspace{13\baselineskip}
\begin{qitem}{12}
\qmeta{Single-correct MCQ}{$O(N^4)$ Hamiltonian terms}
In second quantisation $H = \sum_{pq}h_{pq}a_p^\dagger a_q + \tfrac12\sum_{pqrs}h_{pqrs}a_p^\dagger a_q^\dagger a_ra_s$ (+ constant). For $N$ spin-orbitals the number of
terms in this Hamiltonian grows as
\begin{itemize}[leftmargin=2.6em, labelsep=0.4em, itemsep=2pt, topsep=3pt]
\item[\optlabel{A}] $O(N)$, one term per orbital
\item[\correctlabel{B}] $O(N^4)$, from the two-electron integrals $h_{pqrs}$
\item[\optlabel{C}] $O(N^2)$, from the one-electron integrals only
\item[\optlabel{D}] $O(2^N)$, like the dimension of the Fock space
\end{itemize}
\end{qitem}
\begin{solbox}{B}
\textbf{\color{brand}Correct answer.} \textcolor{okgreen}{\textbf{(B)}} $O(N^4)$, from the two-electron integrals $h_{pqrs}$

\textbf{\color{brand}Concepts required.} Second-quantised Hamiltonian; number of integrals.

\textbf{\color{brand}Relevant formulae.} one-body $h_{pq}$: $N^2$; two-body $h_{pqrs}$: $N^4$ (reduced $\approx 8\times$ by symmetry).

\textbf{\color{brand}Step-by-step explanation.} The one-body part has $N^2$ coefficients $h_{pq}$, the two-body part $N^4$ coefficients $h_{pqrs}$ (four free indices). Permutation symmetries
reduce the count by a constant factor (about 8) but not the power, so the total is $O(N^4)$. After a fermion-to-qubit mapping each
term becomes a few Pauli strings, so the number of Pauli terms --- and the measurement cost of VQE --- also grows like $N^4$.
The Hilbert space is $2^N$-dimensional, but the Hamiltonian itself has only polynomially many terms; that is what makes
$\braket{H} = \sum_kc_k\braket{P_k}$ practical. H$_2$ (STO-3G): 15 Pauli terms.

\textbf{\color{brand}Why the other options are incorrect.}
\begin{itemize}[leftmargin=2.4em, itemsep=1pt, topsep=2pt]
\item[\textbf{(A)}] $O(N)$ would mean no interactions between orbitals.
\item[\textbf{(C)}] $O(N^2)$ ignores the two-electron integrals, which dominate.
\item[\textbf{(D)}] $2^N$ is the state-space dimension, not the number of terms.
\end{itemize}
\textbf{\color{brand}Key takeaway.} Polynomially many terms make $\braket H = \sum c_k\braket{P_k}$ feasible; $N^4$ still makes measurement costly.
\end{solbox}
\qsep

\needspace{13\baselineskip}
\begin{qitem}{13}
\qmeta{Single-correct MCQ}{$e^{i\theta P}$ identity}
For a Pauli operator $P$ ($P^2 = I$), $e^{i\theta P} = \cos\theta\,I + i\sin\theta\,P$. Hence $e^{i\pi X/4}$ equals
\par\smallskip\noindent\begin{tabularx}{\linewidth}{@{}*{4}{L}@{}}
\optlabel{A}\,$\begin{pmatrix}0 & i\\ i & 0\end{pmatrix}$ & \optlabel{B}\,$\dfrac{1}{\sqrt2}\begin{pmatrix}1 & 1\\ 1 & -1\end{pmatrix}$ & \correctlabel{C}\,$\dfrac{1}{\sqrt2}\begin{pmatrix}1 & i\\ i & 1\end{pmatrix}$ & \optlabel{D}\,$\dfrac{1}{\sqrt2}\begin{pmatrix}1 & -i\\ -i & 1\end{pmatrix}$
\end{tabularx}
\end{qitem}
\begin{solbox}{C}
\textbf{\color{brand}Correct answer.} \textcolor{okgreen}{\textbf{(C)}} $\dfrac{1}{\sqrt2}\begin{pmatrix}1 & i\\ i & 1\end{pmatrix}$

\textbf{\color{brand}Concepts required.} Exponential of an involutory matrix.

\textbf{\color{brand}Relevant formulae.} $e^{i\theta P} = \cos\theta\,I + i\sin\theta\,P$ for $P^2 = I$.

\textbf{\color{brand}Step-by-step explanation.} Power series: even powers of $P$ give $I$, odd powers give $P$, so $e^{i\theta P} = \sum_k\frac{(i\theta)^k}{k!}P^k = \cos\theta\,I + i\sin\theta\,P$.
With $\theta = \pi/4$: $\cos\frac\pi4 = \sin\frac\pi4 = \frac{1}{\sqrt2}$, so $e^{i\pi X/4} = \frac{1}{\sqrt2}(I + iX) = \frac{1}{\sqrt2}\begin{psmallmatrix}1 & i\\ i & 1\end{psmallmatrix}$.
The same identity, with $P = X_0X_1X_2Y_3$, is how CUDA-Q's \texttt{exp\_pauli(theta, q, "XXXY")} acts in the notebook's UCC ansatz.
$\frac{1}{\sqrt2}(I - iX)$ would be $e^{-i\pi X/4} = R_X(\pi/2)$.

\textbf{\color{brand}Why the other options are incorrect.}
\begin{itemize}[leftmargin=2.4em, itemsep=1pt, topsep=2pt]
\item[\textbf{(A)}] $iX = e^{i\pi X/2}$, a quarter period further.
\item[\textbf{(B)}] That is the Hadamard gate, not an $X$ rotation.
\item[\textbf{(D)}] $\frac{1}{\sqrt2}(I - iX)$ is $e^{-i\pi X/4}$ (opposite sign).
\end{itemize}
\textbf{\color{brand}Key takeaway.} Same identity powers \texttt{exp\_pauli}: $e^{i\theta P} = \cos\theta + i\sin\theta P$.
\end{solbox}
\qsep

\needspace{20\baselineskip}
\begin{passagebox}{Paragraph for Questions 14 and 15}
A single qubit has the Hamiltonian $H = -Z + \tfrac12 X = \begin{pmatrix}-1 & \tfrac12\\ \tfrac12 & 1\end{pmatrix}$.
A one-parameter trial state is prepared as $\ket{\psi(\theta)} = R_Y(\theta)\ket0 = \cos\tfrac\theta2\ket0 + \sin\tfrac\theta2\ket1$, and its
energy $E(\theta) = \braket{\psi(\theta)|H|\psi(\theta)}$ is minimised over $\theta \in (-\pi, \pi]$.
\end{passagebox}
\needspace{13\baselineskip}
\begin{qitem}{14}
\qmeta{Paragraph based}{One-qubit $E(\theta)$ and $E_0$}
The energy function $E(\theta)$ and the exact ground-state energy $E_0$ are
\begin{itemize}[leftmargin=2.6em, labelsep=0.4em, itemsep=2pt, topsep=3pt]
\item[\optlabel{A}] $E(\theta) = -\cos\theta + \tfrac12\sin\theta$ and $E_0 = -1$
\item[\optlabel{B}] $E(\theta) = -\cos\tfrac\theta2 + \tfrac12\sin\tfrac\theta2$ and $E_0 = -\tfrac{\sqrt5}{2}$
\item[\optlabel{C}] $E(\theta) = -\cos\theta + \sin\theta$ and $E_0 = -\sqrt2$
\item[\correctlabel{D}] $E(\theta) = -\cos\theta + \tfrac12\sin\theta$ and $E_0 = -\tfrac{\sqrt5}{2} \approx -1.118$
\end{itemize}
\end{qitem}
\begin{solbox}{D}
\textbf{\color{brand}Correct answer.} \textcolor{okgreen}{\textbf{(D)}} $E(\theta) = -\cos\theta + \tfrac12\sin\theta$ and $E_0 = -\tfrac{\sqrt5}{2} \approx -1.118$

\textbf{\color{brand}Concepts required.} Expectation values of a parameterised qubit; diagonalising a $2\times2$ Hamiltonian.

\textbf{\color{brand}Relevant formulae.} $R_Y(\theta)\ket0$: $\braket Z = \cos\theta$, $\braket X = \sin\theta$; $\lambda = \pm\sqrt{a^2 + b^2}$ for $aZ + bX$.

\textbf{\color{brand}Step-by-step explanation.} For $\ket\psi = \cos\frac\theta2\ket0 + \sin\frac\theta2\ket1$: $\braket Z = \cos^2\frac\theta2 - \sin^2\frac\theta2 = \cos\theta$ and
$\braket X = 2\cos\frac\theta2\sin\frac\theta2 = \sin\theta$. By linearity $E(\theta) = -\braket Z + \tfrac12\braket X = -\cos\theta + \tfrac12\sin\theta$.
Eigenvalues of $\begin{psmallmatrix}-1&1/2\\1/2&1\end{psmallmatrix}$: $\lambda^2 - 1 - \tfrac14 = 0$, $\lambda = \pm\sqrt{5/4} = \pm\frac{\sqrt5}{2}$, so $E_0 = -1.1180$.
Consistency: $E(\theta) = -\sqrt{1 + \tfrac14}\cos(\theta + \varphi)$ has minimum $-\frac{\sqrt5}{2}$ --- the one-parameter $R_Y$ ansatz reaches $E_0$ exactly,
because the ground state of a real $2\times2$ Hamiltonian is a real vector.

\textbf{\color{brand}Why the other options are incorrect.}
\begin{itemize}[leftmargin=2.4em, itemsep=1pt, topsep=2pt]
\item[\textbf{(A)}] $-1$ is $E(0) = \braket{0|H|0}$, the diagonal element, not the minimum.
\item[\textbf{(B)}] Half angles appear in the amplitudes, not in the expectation values.
\item[\textbf{(C)}] The $X$ coefficient is $\frac12$, not $1$; $-\sqrt2$ belongs to $-Z + X$.
\end{itemize}
\textbf{\color{brand}Key takeaway.} For $aZ + bX$: $E_0 = -\sqrt{a^2 + b^2}$, reachable with one $R_Y$ angle.
\end{solbox}
\qsep

\needspace{13\baselineskip}
\begin{qitem}{15}
\qmeta{Paragraph based}{Optimal angle, starting error}
The optimal angle $\theta^*$, and the energy error $E(0) - E_0$ of the starting point $\theta = 0$ (the state $\ket0$), are
\par\smallskip\noindent\begin{tabularx}{\linewidth}{@{}*{2}{L}@{}}
\optlabel{A}\,$\theta^* = -\arctan\tfrac12 \approx -0.464$ rad; error $0$ & \optlabel{B}\,$\theta^* = -\arctan 2 \approx -1.107$ rad; error $\approx 0.118$ \\[3pt]
\optlabel{C}\,$\theta^* = +\arctan\tfrac12 \approx 0.464$ rad; error $\approx 0.118$ & \correctlabel{D}\,$\theta^* = -\arctan\tfrac12 \approx -0.464$ rad; error $\approx 0.118$
\end{tabularx}
\end{qitem}
\begin{solbox}{D}
\textbf{\color{brand}Correct answer.} \textcolor{okgreen}{\textbf{(D)}} $\theta^* = -\arctan\tfrac12 \approx -0.464$ rad; error $\approx 0.118$

\textbf{\color{brand}Concepts required.} Stationary points; second-derivative test.

\textbf{\color{brand}Relevant formulae.} $E'(\theta) = \sin\theta + \frac12\cos\theta$; $\tan\theta^* = -\frac12$; $E'' > 0$ at a minimum.

\textbf{\color{brand}Step-by-step explanation.} $E'(\theta) = \sin\theta + \tfrac12\cos\theta = 0 \Rightarrow \tan\theta = -\tfrac12$. The minimum needs $E'' = \cos\theta - \tfrac12\sin\theta > 0$: with
$\theta^* = -\arctan\tfrac12 = -0.4636$, $\cos\theta^* = 0.894 > 0$ and $\sin\theta^* = -0.447$, so $E'' > 0$ and
$E(\theta^*) = -0.894 - 0.224 = -1.118 = E_0$. (The other root $\pi - 0.4636$ is the maximum $+1.118$.)
$E(0) = -1$, so the starting error is $-1 - (-1.118) = 0.118$. The negative angle mixes in $\ket1$ with a \emph{negative} amplitude because the
off-diagonal element $+\tfrac12$ is positive.

\textbf{\color{brand}Why the other options are incorrect.}
\begin{itemize}[leftmargin=2.4em, itemsep=1pt, topsep=2pt]
\item[\textbf{(A)}] The start $\ket0$ is not the ground state, so its error is not zero.
\item[\textbf{(B)}] $\arctan2$ inverts the ratio; $E(-1.107) = -0.447 - 0.447 = -0.894$.
\item[\textbf{(C)}] $+0.464$ has the wrong sign: there $E = -0.894 + 0.224 = -0.671$.
\end{itemize}
\textbf{\color{brand}Key takeaway.} Off-diagonal sign decides the sign of the optimal mixing angle.
\end{solbox}
\qsep

\needspace{13\baselineskip}
\begin{qitem}{16}
\qmeta{Single-correct MCQ}{Pauli decomposition ($2\times2$)}
The Hermitian matrix $M = \begin{pmatrix}3 & 1 - 2i\\ 1 + 2i & -1\end{pmatrix}$ written as $aI + bX + cY + dZ$ is
\par\smallskip\noindent\begin{tabularx}{\linewidth}{@{}*{4}{L}@{}}
\correctlabel{A}\,$I + X + 2Y + 2Z$ & \optlabel{B}\,$3I + X + 2Y - Z$ & \optlabel{C}\,$I + X - 2Y + 2Z$ & \optlabel{D}\,$2I + X + 2Y + Z$
\end{tabularx}
\end{qitem}
\begin{solbox}{A}
\textbf{\color{brand}Correct answer.} \textcolor{okgreen}{\textbf{(A)}} $I + X + 2Y + 2Z$

\textbf{\color{brand}Concepts required.} Pauli decomposition of a Hermitian matrix.

\textbf{\color{brand}Relevant formulae.} $M = \sum_P c_PP$, $c_P = \frac12\operatorname{tr}(PM)$; $M_{01} = b - ic$, $M_{00} = a + d$, $M_{11} = a - d$.

\textbf{\color{brand}Step-by-step explanation.} Use $c_P = \tfrac12\operatorname{tr}(PM)$ (Paulis are orthogonal: $\operatorname{tr}(PQ) = 2\delta_{PQ}$).
$a = \tfrac12(3 - 1) = 1$; $d = \tfrac12(3 + 1) = 2$; $b = \tfrac12(M_{01} + M_{10}) = \tfrac12\cdot2 = 1$; $c = \tfrac12\operatorname{tr}(YM) = \tfrac12(-iM_{10} + iM_{01}) = \tfrac12(-i(1 + 2i) + i(1 - 2i)) = \tfrac12(2 + 2) = 2$.
Check: $X + 2Y = \begin{psmallmatrix}0 & 1 - 2i\\ 1 + 2i & 0\end{psmallmatrix}$ \checkmark{} and $I + 2Z = \operatorname{diag}(3, -1)$ \checkmark{}. Real coefficients are guaranteed because $M$ is Hermitian --- the
reason the notebook rejects complex coefficients in \texttt{CUSTOM\_HAMILTONIAN}.

\textbf{\color{brand}Why the other options are incorrect.}
\begin{itemize}[leftmargin=2.4em, itemsep=1pt, topsep=2pt]
\item[\textbf{(B)}] $3I - Z$ gives diagonal $(2, 4)$.
\item[\textbf{(C)}] $-2Y$ gives $M_{01} = 1 + 2i$, the conjugate of the given entry.
\item[\textbf{(D)}] $2I + Z$ gives diagonal $(3, 1)$, not $(3, -1)$.
\end{itemize}
\textbf{\color{brand}Key takeaway.} Real coefficients for every Pauli string $\iff$ Hermitian matrix.
\end{solbox}
\qsep

\needspace{13\baselineskip}
\begin{qitem}{17}
\qmeta{Single-correct MCQ}{Pauli decomposition (diagonal)}
The diagonal two-qubit Hamiltonian $H = \operatorname{diag}(1, 2, 3, 4)$ (basis $\ket{00}, \ket{01}, \ket{10}, \ket{11}$) written as a sum of Pauli strings is
\par\smallskip\noindent\begin{tabularx}{\linewidth}{@{}*{2}{L}@{}}
\optlabel{A}\,$2.5\,II - 0.5\,ZI - 1.0\,IZ$ & \optlabel{B}\,$2.5\,II - 1.0\,ZI - 0.5\,IZ + 0.5\,ZZ$ \\[3pt]
\optlabel{C}\,$2.5\,II + 1.0\,ZI + 0.5\,IZ$ & \correctlabel{D}\,$2.5\,II - 1.0\,ZI - 0.5\,IZ$
\end{tabularx}
\end{qitem}
\begin{solbox}{D}
\textbf{\color{brand}Correct answer.} \textcolor{okgreen}{\textbf{(D)}} $2.5\,II - 1.0\,ZI - 0.5\,IZ$

\textbf{\color{brand}Concepts required.} Pauli decomposition of diagonal operators.

\textbf{\color{brand}Relevant formulae.} $c_P = 2^{-n}\operatorname{tr}(PH)$; $ZI = \operatorname{diag}(1,1,-1,-1)$, $IZ = \operatorname{diag}(1,-1,1,-1)$.

\textbf{\color{brand}Step-by-step explanation.} Coefficients $c_P = \tfrac14\operatorname{tr}(PH)$. On the basis states, $ZI = \operatorname{diag}(1, 1, -1, -1)$, $IZ = \operatorname{diag}(1, -1, 1, -1)$, $ZZ = \operatorname{diag}(1, -1, -1, 1)$.
$c_{II} = \frac{1+2+3+4}{4} = 2.5$; $c_{ZI} = \frac{1+2-3-4}{4} = -1$; $c_{IZ} = \frac{1-2+3-4}{4} = -0.5$; $c_{ZZ} = \frac{1-2-3+4}{4} = 0$.
Check $\ket{11}$: $2.5 + 1 + 0.5 = 4$ \checkmark{}. The first qubit (more significant) changes the energy by 2, so its $Z$ coefficient is $-1$.

\textbf{\color{brand}Why the other options are incorrect.}
\begin{itemize}[leftmargin=2.4em, itemsep=1pt, topsep=2pt]
\item[\textbf{(A)}] Swaps the two qubits' roles; the first qubit changes the energy by 2, the second by 1.
\item[\textbf{(B)}] $c_{ZZ} = \frac{1 - 2 - 3 + 4}{4} = 0$: the levels are additive, no interaction term.
\item[\textbf{(C)}] Positive $Z$ coefficients would make $\ket{00}$ the highest level; it is the lowest ($1$).
\end{itemize}
\textbf{\color{brand}Key takeaway.} Diagonal Hamiltonians need only $I$/$Z$ strings.
\end{solbox}
\qsep

\needspace{13\baselineskip}
\begin{qitem}{18}
\qmeta{Assertion--Reason}{RHF at dissociation}
\textbf{Assertion (A):} For H$_2$ at large bond length, the restricted Hartree--Fock (single closed-shell determinant) energy lies far above
the exact energy.\\[1mm]
\textbf{Reason (R):} The nuclear--nuclear repulsion $1/R$ tends to zero as $R \to \infty$.
\begin{itemize}[leftmargin=2.6em, labelsep=0.4em, itemsep=2pt, topsep=3pt]
\item[\optlabel{A}] Both Assertion (A) and Reason (R) are true, and R is the correct explanation of A
\item[\correctlabel{B}] Both A and R are true, but R is \emph{not} the correct explanation of A
\item[\optlabel{C}] A is true, but R is false
\item[\optlabel{D}] A is false, but R is true
\end{itemize}
\end{qitem}
\begin{solbox}{B}
\textbf{\color{brand}Correct answer.} \textcolor{okgreen}{\textbf{(B)}} Both A and R are true, but R is \emph{not} the correct explanation of A

\textbf{\color{brand}Concepts required.} Static correlation; RHF failure at dissociation.

\textbf{\color{brand}Relevant formulae.} $\sigma_g^2 \propto (1s_Q.1s_B + 1s_Q.21s_A) + (1s_A^2 + 1s_B^2)$: $50\,\%$ ionic.

\textbf{\color{brand}Step-by-step explanation.} \textbf{A is true:} the RHF determinant puts both electrons in the bonding orbital $\sigma_g = (1s_A + 1s_B)/\sqrt2$. Expanded in atomic
orbitals it contains ionic configurations (H$^+$H$^-$) with $50\,\%$ weight, which cost a lot of energy for separated atoms. The exact
state is $\approx \cos\theta\ket{\sigma_g^2} - \sin\theta\ket{\sigma_u^2}$ with $\theta \to \pi/4$: static correlation. (Same structure as the notebook's UCC
state $\cos\theta\ket{1100} - \sin\theta\ket{0011}$.)
\textbf{R is true} but is \emph{not} the explanation: $V_{NN}$ is identical in RHF and the exact calculation and cancels in the error. Answer: (B).

\textbf{\color{brand}Why the other options are incorrect.}
\begin{itemize}[leftmargin=2.4em, itemsep=1pt, topsep=2pt]
\item[\textbf{(A)}] needs A true, R true \emph{and} R explaining A; $V_{NN}$ is the same in both methods and cannot explain the error.
\item[\textbf{(C)}] needs A true and R false; R is true.
\item[\textbf{(D)}] needs A false and R true; A is true.
\end{itemize}
\textbf{\color{brand}Key takeaway.} Bond breaking needs more than one determinant --- VQE's double excitation provides it.
\end{solbox}
\qsep

\needspace{13\baselineskip}
\begin{qitem}{19}
\qmeta{Single-correct MCQ}{Correlation energy}
For the notebook's H$_2$ Hamiltonian, the Hartree--Fock energy is $E_{\mathrm{HF}} = \braket{1100|H|1100} = -1.116759$ Ha and the FCI energy is
$-1.137284$ Ha. The correlation energy $E_{\mathrm{corr}} = E_{\mathrm{FCI}} - E_{\mathrm{HF}}$ is about
\par\smallskip\noindent\begin{tabularx}{\linewidth}{@{}*{4}{L}@{}}
\correctlabel{A}\,$-20.5$ mHa & \optlabel{B}\,$+20.5$ mHa & \optlabel{C}\,$-2.05$ mHa & \optlabel{D}\,$-1137$ mHa
\end{tabularx}
\end{qitem}
\begin{solbox}{A}
\textbf{\color{brand}Correct answer.} \textcolor{okgreen}{\textbf{(A)}} $-20.5$ mHa

\textbf{\color{brand}Concepts required.} Correlation energy.

\textbf{\color{brand}Relevant formulae.} $E_{\mathrm{corr}} = E_{\mathrm{FCI}} - E_{\mathrm{HF}} \le 0$.

\textbf{\color{brand}Step-by-step explanation.} $E_{\mathrm{corr}} = -1.137284 - (-1.116759) = -0.020525$ Ha $= -20.5$ mHa. It is negative because FCI is variationally
\emph{lower} than HF (HF is one particular state in the same space). $20.5$ mHa $\approx 12.9$ kcal/mol, far above chemical accuracy:
HF alone is not good enough, and VQE has to recover this correlation energy. The UCC parameter does exactly that (see Q.59--Q.60).

\textbf{\color{brand}Why the other options are incorrect.}
\begin{itemize}[leftmargin=2.4em, itemsep=1pt, topsep=2pt]
\item[\textbf{(B)}] Positive sign would mean FCI is above HF, violating the variational principle.
\item[\textbf{(C)}] Off by a factor of 10: $-0.020525$ Ha is $-20.5$ mHa, not $-2.05$ mHa.
\item[\textbf{(D)}] $-1137$ mHa is the total FCI energy.
\end{itemize}
\textbf{\color{brand}Key takeaway.} H$_2$: $20.5$ mHa of correlation, about $13\times$ chemical accuracy.
\end{solbox}
\qsep

\needspace{13\baselineskip}
\begin{qitem}{20}
\qmeta{Single-correct MCQ}{Excited-state contamination}
A trial state of H$_2$ has $90\,\%$ weight on the ground state ($E_0 = -1.13728$ Ha) and $10\,\%$ on another eigenstate of the 4-qubit Hamiltonian with energy
$E_1 = -0.53821$ Ha. Its energy expectation and its error above $E_0$ are
\par\smallskip\noindent\begin{tabularx}{\linewidth}{@{}*{2}{L}@{}}
\optlabel{A}\,$-1.1971$ Ha and $-59.9$ mHa & \correctlabel{B}\,$-1.0774$ Ha and $59.9$ mHa \\[3pt]
\optlabel{C}\,$-0.8378$ Ha and $299.5$ mHa & \optlabel{D}\,$-1.1373$ Ha and $0$ mHa
\end{tabularx}
\end{qitem}
\begin{solbox}{B}
\textbf{\color{brand}Correct answer.} \textcolor{okgreen}{\textbf{(B)}} $-1.0774$ Ha and $59.9$ mHa

\textbf{\color{brand}Concepts required.} Energy of a state expanded in eigenstates.

\textbf{\color{brand}Relevant formulae.} $E = \sum|c_k|^2E_k$; error $= \sum_{k>0}|c_k|^2(E_k - E_0)$.

\textbf{\color{brand}Step-by-step explanation.} $E = \sum_k|c_k|^2E_k = 0.9(-1.13728) + 0.1(-0.53821) = -1.02355 - 0.05382 = -1.07738$ Ha.
Error $E - E_0 = 0.1(E_1 - E_0) = 0.1\times0.59907 = 0.0599$ Ha $= 59.9$ mHa. The error is \emph{linear} in the contaminating weight and is
always non-negative (variational principle). The energy is lower than the energies of all excited states but still above $E_0$.

\textbf{\color{brand}Why the other options are incorrect.}
\begin{itemize}[leftmargin=2.4em, itemsep=1pt, topsep=2pt]
\item[\textbf{(A)}] Below $E_0$: impossible for a normalised state.
\item[\textbf{(C)}] Uses $50\,\%$ weight instead of $10\,\%$.
\item[\textbf{(D)}] Ignores the $10\,\%$ excited weight.
\end{itemize}
\textbf{\color{brand}Key takeaway.} Energy errors are linear in contamination; state errors are larger than they look.
\end{solbox}
\qsep

\newpage

\begin{sectionbanner}SECTION B \enspace$\cdot$\enspace Quantum Computing Foundations \hfill Q.21--40\end{sectionbanner}

\needspace{13\baselineskip}
\begin{qitem}{21}
\qmeta{Single-correct MCQ}{$R_Y$ and the Born rule}
The state $R_Y(\pi/3)\ket0$ is measured in the computational basis. The probability of the outcome 1 is
\par\smallskip\noindent\begin{tabularx}{\linewidth}{@{}*{4}{L}@{}}
\optlabel{A}\,$\tfrac34$ & \optlabel{B}\,$\tfrac{\sqrt3}{2}$ & \optlabel{C}\,$\tfrac12$ & \correctlabel{D}\,$\tfrac14$
\end{tabularx}
\end{qitem}
\begin{solbox}{D}
\textbf{\color{brand}Correct answer.} \textcolor{okgreen}{\textbf{(D)}} $\tfrac14$

\textbf{\color{brand}Concepts required.} $R_Y$ rotation and the Born rule.

\textbf{\color{brand}Relevant formulae.} $R_Y(\theta)\ket0 = \cos\frac\theta2\ket0 + \sin\frac\theta2\ket1$; $P(1) = \sin^2\frac\theta2$.

\textbf{\color{brand}Step-by-step explanation.} $R_Y(\theta) = e^{-i\theta Y/2} = \begin{psmallmatrix}\cos\frac\theta2 & -\sin\frac\theta2\\ \sin\frac\theta2 & \cos\frac\theta2\end{psmallmatrix}$, so $R_Y(\theta)\ket0 = \cos\frac\theta2\ket0 + \sin\frac\theta2\ket1$
(real amplitudes). With $\theta = \pi/3$: $P(1) = \sin^2(\pi/6) = \tfrac14$, $P(0) = \cos^2(\pi/6) = \tfrac34$. The notebook's hardware-efficient
ansatz is built from exactly these \texttt{ry} rotations.

\textbf{\color{brand}Why the other options are incorrect.}
\begin{itemize}[leftmargin=2.4em, itemsep=1pt, topsep=2pt]
\item[\textbf{(A)}] $\frac34 = \cos^2(\pi/6)$ is $P(0)$, the probability of the \emph{other} outcome.
\item[\textbf{(B)}] $\frac{\sqrt3}{2} = \cos(\pi/6)$ is the \emph{amplitude} of $\ket0$; probabilities are squared moduli ($P(1) = \sin^2(\pi/6) = \frac14$).
\item[\textbf{(C)}] $\frac12$ needs $\theta = \pi/2$.
\end{itemize}
\textbf{\color{brand}Key takeaway.} Half angles in amplitudes: $R_Y(\theta)$ rotates the Bloch vector by $\theta$.
\end{solbox}
\qsep

\needspace{13\baselineskip}
\begin{qitem}{22}
\qmeta{Single-correct MCQ}{Pauli products}
The two-qubit operator $(X\otimes Y)(Y\otimes X)$ equals
\par\smallskip\noindent\begin{tabularx}{\linewidth}{@{}*{4}{L}@{}}
\correctlabel{A}\,$Z\otimes Z$ & \optlabel{B}\,$Y\otimes Y$ & \optlabel{C}\,$-Z\otimes Z$ & \optlabel{D}\,$i\,Z\otimes Z$
\end{tabularx}
\end{qitem}
\begin{solbox}{A}
\textbf{\color{brand}Correct answer.} \textcolor{okgreen}{\textbf{(A)}} $Z\otimes Z$

\textbf{\color{brand}Concepts required.} Pauli multiplication; tensor products.

\textbf{\color{brand}Relevant formulae.} $XY = iZ$, $YX = -iZ$; $(A\otimes B)(C\otimes D) = AC\otimes BD$.

\textbf{\color{brand}Step-by-step explanation.} Multiply factor-wise: $(X\otimes Y)(Y\otimes X) = (XY)\otimes(YX)$. Pauli table: $XY = iZ$, $YX = -iZ$. So the product is
$(iZ)\otimes(-iZ) = -i^2\,Z\otimes Z = Z\otimes Z$. Products of Pauli strings are again Pauli strings up to a phase $\pm1, \pm i$ --- the
algebra behind Hamiltonian decompositions such as the notebook's $\texttt{XXYY}$, $\texttt{XYYX}$, \dots\ terms.

\textbf{\color{brand}Why the other options are incorrect.}
\begin{itemize}[leftmargin=2.4em, itemsep=1pt, topsep=2pt]
\item[\textbf{(B)}] $Y\otimes Y$ ignores the products.
\item[\textbf{(C)}] $-Z\otimes Z$ uses $YX = +iZ$.
\item[\textbf{(D)}] Keeps a phase $i$; the two phases $i$ and $-i$ multiply to $1$.
\end{itemize}
\textbf{\color{brand}Key takeaway.} Track the phases $\pm i$ factor by factor.
\end{solbox}
\qsep

\needspace{13\baselineskip}
\begin{qitem}{23}
\qmeta{Matching lists}{Bloch vectors}
Match each single-qubit state (List-I) with its Bloch vector $(\braket X, \braket Y, \braket Z)$ (List-II).
\par\smallskip\noindent{\renewcommand{\arraystretch}{1.35}%
\begin{tabularx}{\linewidth}{@{}c>{\raggedright\arraybackslash}p{0.33\linewidth}@{\hspace{1.2em}}cL@{}}
\toprule
\multicolumn{2}{@{}l}{\textbf{List-I}} & \multicolumn{2}{l@{}}{\textbf{List-II}} \\ \midrule
\textbf{(P)} & $\ket0$ & \textbf{(1)} & $(0, 1, 0)$ \\
\textbf{(Q)} & $\ket- = \tfrac{1}{\sqrt2}(\ket0 - \ket1)$ & \textbf{(2)} & $(\tfrac{\sqrt3}{2}, 0, \tfrac12)$ \\
\textbf{(R)} & $\tfrac{1}{\sqrt2}(\ket0 + i\ket1)$ & \textbf{(3)} & $(0, 0, 1)$ \\
\textbf{(S)} & $R_Y(\pi/3)\ket0$ & \textbf{(4)} & $(-1, 0, 0)$ \\
\bottomrule
\end{tabularx}}
\par\smallskip\noindent\begin{tabularx}{\linewidth}{@{}*{2}{L}@{}}
\optlabel{A}\,P$\to$3, Q$\to$4, R$\to$2, S$\to$1 & \optlabel{B}\,P$\to$3, Q$\to$1, R$\to$4, S$\to$2 \\[3pt]
\optlabel{C}\,P$\to$4, Q$\to$3, R$\to$1, S$\to$2 & \correctlabel{D}\,P$\to$3, Q$\to$4, R$\to$1, S$\to$2
\end{tabularx}
\end{qitem}
\begin{solbox}{D}
\textbf{\color{brand}Correct answer.} \textcolor{okgreen}{\textbf{(D)}} P$\to$3, Q$\to$4, R$\to$1, S$\to$2

\textbf{\color{brand}Concepts required.} Bloch vector of a qubit state.

\textbf{\color{brand}Relevant formulae.} $(\braket X, \braket Y, \braket Z) = (2\mathrm{Re}\,a^*b, 2\mathrm{Im}\,a^*b, |a|^2 - |b|^2)$.

\textbf{\color{brand}Step-by-step explanation.} For $\ket\psi = a\ket0 + b\ket1$: $\braket Z = |a|^2 - |b|^2$, $\braket X = 2\,\mathrm{Re}(a^*b)$, $\braket Y = 2\,\mathrm{Im}(a^*b)$.
\begin{itemize}\itemsep0pt
\item P: $a = 1$, $b = 0$: $(0, 0, 1)$ (3).
\item Q: $a = -b = \frac{1}{\sqrt2}$: $\braket X = -1$: $(-1, 0, 0)$ (4).
\item R: $a^*b = \frac i2$: $\braket Y = 1$: $(0, 1, 0)$ (1).
\item S: $R_Y(\theta)\ket0$ has Bloch vector $(\sin\theta, 0, \cos\theta)$; $\theta = \pi/3$: $(\frac{\sqrt3}{2}, 0, \frac12)$ (2). $R_Y$ rotates in the $x$--$z$ plane:
real amplitudes, $\braket Y = 0$.
\end{itemize}

\textbf{\color{brand}Why the other options are incorrect.}
\begin{itemize}[leftmargin=2.4em, itemsep=1pt, topsep=2pt]
\item[\textbf{(A)}] R$\to$2, S$\to$1 gives $\ket{+i}$ a real Bloch vector; $R_Y$ never creates $\braket Y \neq 0$.
\item[\textbf{(B)}] Q$\to$1, R$\to$4 swaps $\ket-$ and $\ket{+i}$.
\item[\textbf{(C)}] P$\to$4 would put $\ket0$ on the $-x$ axis.
\end{itemize}
\textbf{\color{brand}Key takeaway.} $R_Y$ keeps states in the real $x$--$z$ plane.
\end{solbox}
\qsep

\needspace{13\baselineskip}
\begin{qitem}{24}
\qmeta{Matching lists}{Measurement basis changes}
To estimate $\braket P$ from computational-basis measurements, each Pauli string is first rotated into a string of $Z$'s and $I$'s. Match each
observable (List-I) with the gates applied before measuring (List-II).
\par\smallskip\noindent{\renewcommand{\arraystretch}{1.35}%
\begin{tabularx}{\linewidth}{@{}c>{\raggedright\arraybackslash}p{0.33\linewidth}@{\hspace{1.2em}}cL@{}}
\toprule
\multicolumn{2}{@{}l}{\textbf{List-I}} & \multicolumn{2}{l@{}}{\textbf{List-II}} \\ \midrule
\textbf{(P)} & $X$ & \textbf{(1)} & $H$ on qubit 0 only \\
\textbf{(Q)} & $Y$ & \textbf{(2)} & nothing \\
\textbf{(R)} & $Z$ & \textbf{(3)} & $H$ \\
\textbf{(S)} & $X_0Z_1$ & \textbf{(4)} & $S^\dagger$, then $H$ \\
\bottomrule
\end{tabularx}}
\par\smallskip\noindent\begin{tabularx}{\linewidth}{@{}*{2}{L}@{}}
\optlabel{A}\,P$\to$3, Q$\to$2, R$\to$4, S$\to$1 & \optlabel{B}\,P$\to$3, Q$\to$4, R$\to$1, S$\to$2 \\[3pt]
\correctlabel{C}\,P$\to$3, Q$\to$4, R$\to$2, S$\to$1 & \optlabel{D}\,P$\to$4, Q$\to$3, R$\to$2, S$\to$1
\end{tabularx}
\end{qitem}
\begin{solbox}{C}
\textbf{\color{brand}Correct answer.} \textcolor{okgreen}{\textbf{(C)}} P$\to$3, Q$\to$4, R$\to$2, S$\to$1

\textbf{\color{brand}Concepts required.} Measurement basis changes.

\textbf{\color{brand}Relevant formulae.} $HXH = Z$; $(HS^\dagger)Y(HS^\dagger)^\dagger = Z$.

\textbf{\color{brand}Step-by-step explanation.} If $UPU^\dagger = Z$, measuring $Z$ after applying $U$ gives the eigenvalues of $P$: $\braket{\psi|P|\psi} = \braket{U\psi|Z|U\psi}$.
\begin{itemize}\itemsep0pt
\item P: $HXH = Z$: apply $H$ (3).
\item Q: $(HS^\dagger)\,Y\,(HS^\dagger)^\dagger = Z$: apply $S^\dagger$ first, then $H$ (4) (checked numerically in \texttt{verify\_answers.py}).
\item R: $Z$ is already diagonal: nothing (2).
\item S: rotate only the $X$ factor: $H$ on qubit 0; then $\braket{X_0Z_1}$ is the average of $(-1)^{b_0\oplus b_1}$ (1).
\end{itemize}
The notebook's text describes the same rule (``$H$ for $X$, $S^\dagger H$ for $Y$''); \texttt{cudaq.observe} does it automatically with shots.

\textbf{\color{brand}Why the other options are incorrect.}
\begin{itemize}[leftmargin=2.4em, itemsep=1pt, topsep=2pt]
\item[\textbf{(A)}] Q$\to$2 means measuring $Y$ without a basis change, which measures $Z$.
\item[\textbf{(B)}] R$\to$1, S$\to$2: $Z$ needs nothing, but $X_0Z_1$ needs $H$ on qubit 0.
\item[\textbf{(D)}] P$\to$4, Q$\to$3 swaps the $X$ and $Y$ rotations.
\end{itemize}
\textbf{\color{brand}Key takeaway.} Rotate each Pauli factor to $Z$, then multiply the $\pm1$ outcomes.
\end{solbox}
\qsep

\needspace{13\baselineskip}
\begin{qitem}{25}
\qmeta{Single-correct MCQ}{Bell-state correlations}
For the Bell state $\ket{\Phi^+} = \tfrac{1}{\sqrt2}(\ket{00} + \ket{11})$, the expectation value $\braket{Y\otimes Y}$ is
\par\smallskip\noindent\begin{tabularx}{\linewidth}{@{}*{4}{L}@{}}
\optlabel{A}\,$+1$ & \optlabel{B}\,$-\tfrac12$ & \optlabel{C}\,$0$ & \correctlabel{D}\,$-1$
\end{tabularx}
\end{qitem}
\begin{solbox}{D}
\textbf{\color{brand}Correct answer.} \textcolor{okgreen}{\textbf{(D)}} $-1$

\textbf{\color{brand}Concepts required.} Action of $Y$; Bell-state correlations.

\textbf{\color{brand}Relevant formulae.} $Y\ket0 = i\ket1$, $Y\ket1 = -i\ket0$; $\ket{\Phi^+}$: $\braket{XX} = \braket{ZZ} = 1$, $\braket{YY} = -1$.

\textbf{\color{brand}Step-by-step explanation.} $Y\ket0 = i\ket1$, $Y\ket1 = -i\ket0$. So $(Y\otimes Y)\ket{00} = i\cdot i\ket{11} = -\ket{11}$ and $(Y\otimes Y)\ket{11} = (-i)(-i)\ket{00} = -\ket{00}$:
$(Y\otimes Y)\ket{\Phi^+} = -\ket{\Phi^+}$, an eigenstate with eigenvalue $-1$. (Consistency: $\braket{XX} = 1$, $\braket{ZZ} = 1$ and
$XX\cdot YY\cdot ZZ$-relation $(X\otimes X)(Y\otimes Y) = -(Z\otimes Z)$ gives $1\cdot\braket{YY} = -1$.)

\textbf{\color{brand}Why the other options are incorrect.}
\begin{itemize}[leftmargin=2.4em, itemsep=1pt, topsep=2pt]
\item[\textbf{(A)}] Forgets $i\cdot i = -1$.
\item[\textbf{(B)}] Expectation values of eigenstates are eigenvalues: $\pm1$.
\item[\textbf{(C)}] The state is an eigenstate of $Y\otimes Y$, so the value is $\pm1$, not $0$.
\end{itemize}
\textbf{\color{brand}Key takeaway.} $\ket{\Phi^+}$ is a joint eigenstate of $XX$, $YY$, $ZZ$ with eigenvalues $(1, -1, 1)$.
\end{solbox}
\qsep

\needspace{13\baselineskip}
\begin{qitem}{26}
\qmeta{Single-correct MCQ}{Qubit ordering in CUDA-Q}
The Hartree--Fock state of the notebook's H$_2$ example is $\ket{1100}$, written with qubit 0 first. In the vector returned by \texttt{cudaq.get\_state},
which lists amplitudes with qubit 0 as the \emph{least} significant bit, this basis state has the index
\par\smallskip\noindent\begin{tabularx}{\linewidth}{@{}*{4}{L}@{}}
\optlabel{A}\,$12$ & \correctlabel{B}\,$3$ & \optlabel{C}\,$0$ & \optlabel{D}\,$6$
\end{tabularx}
\end{qitem}
\begin{solbox}{B}
\textbf{\color{brand}Correct answer.} \textcolor{okgreen}{\textbf{(B)}} $3$

\textbf{\color{brand}Concepts required.} Qubit ordering conventions.

\textbf{\color{brand}Relevant formulae.} little-endian index $= \sum_kq_k2^k$; big-endian $= \sum_kq_k2^{n-1-k}$.

\textbf{\color{brand}Step-by-step explanation.} Qubit 0 first means the string $q_0q_1q_2q_3 = 1100$. With qubit 0 as the least significant bit the index is
$\sum_k q_k2^k = 1\cdot1 + 1\cdot2 + 0\cdot4 + 0\cdot8 = 3$. Read as a big-endian number the same string would be $12$. Mixing up the two
orderings is a classic source of wrong amplitudes; the repository's lecture scripts reorder CUDA-Q state vectors before comparing them
with NumPy for exactly this reason.

\textbf{\color{brand}Why the other options are incorrect.}
\begin{itemize}[leftmargin=2.4em, itemsep=1pt, topsep=2pt]
\item[\textbf{(A)}] $12$ reads the string big-endian (qubit 0 most significant).
\item[\textbf{(C)}] Index $0$ is the all-zero state $\ket{0000}$, the empty (vacuum) occupation, not Hartree--Fock.
\item[\textbf{(D)}] $6 = 0110_2$ has the wrong bits set.
\end{itemize}
\textbf{\color{brand}Key takeaway.} Always state the bit order; CUDA-Q state vectors are little-endian.
\end{solbox}
\qsep

\needspace{13\baselineskip}
\begin{qitem}{27}
\qmeta{Single-correct MCQ}{Entangling with CNOT}
Starting from $\ket{00}$, apply $R_Y(\pi/2)$ to qubit 0 and then a CNOT with control 0 and target 1. The resulting state is
\par\smallskip\noindent\begin{tabularx}{\linewidth}{@{}*{2}{L}@{}}
\correctlabel{A}\,$\tfrac{1}{\sqrt2}(\ket{00} + \ket{11})$ & \optlabel{B}\,$\tfrac{1}{\sqrt2}(\ket{00} + \ket{10})$ \\[3pt]
\optlabel{C}\,$\tfrac{1}{\sqrt2}(\ket{00} - \ket{11})$ & \optlabel{D}\,$\tfrac12(\ket{00} + \ket{01} + \ket{10} + \ket{11})$
\end{tabularx}
\end{qitem}
\begin{solbox}{A}
\textbf{\color{brand}Correct answer.} \textcolor{okgreen}{\textbf{(A)}} $\tfrac{1}{\sqrt2}(\ket{00} + \ket{11})$

\textbf{\color{brand}Concepts required.} Entangling with CNOT.

\textbf{\color{brand}Relevant formulae.} $\mathrm{CNOT}\ket{a, b} = \ket{a, a\oplus b}$.

\textbf{\color{brand}Step-by-step explanation.} $R_Y(\pi/2)\ket0 = \cos\frac\pi4\ket0 + \sin\frac\pi4\ket1 = \frac{1}{\sqrt2}(\ket0 + \ket1)$, so after the first gate $\frac{1}{\sqrt2}(\ket{00} + \ket{10})$.
The CNOT flips qubit 1 when qubit 0 is 1: $\ket{10} \to \ket{11}$. Result: the Bell state $\frac{1}{\sqrt2}(\ket{00} + \ket{11})$ --- entangled. This is
how the CNOT ladder of the hardware-efficient ansatz creates entanglement between neighbouring qubits.

\textbf{\color{brand}Why the other options are incorrect.}
\begin{itemize}[leftmargin=2.4em, itemsep=1pt, topsep=2pt]
\item[\textbf{(B)}] That is the state before the CNOT.
\item[\textbf{(C)}] $R_Y(\pi/2)$ gives a $+$ sign: $\sin\frac\pi4 > 0$.
\item[\textbf{(D)}] Needs a rotation on qubit 1 as well.
\end{itemize}
\textbf{\color{brand}Key takeaway.} $R_Y$ + CNOT ladder = the entangling backbone of the hardware-efficient ansatz.
\end{solbox}
\qsep

\needspace{13\baselineskip}
\begin{qitem}{28}
\qmeta{Statement I / II}{Normalisation; measuring $X$}
\textbf{Statement I:} Any circuit of gates applied to $\ket{0\dots0}$ prepares a normalised state, so the ansatz state $\ket{\psi(\boldsymbol\theta)}$ needs no
explicit normalisation.\\[1mm]
\textbf{Statement II:} Measuring in the computational basis directly reveals the eigenvalue of an $X_i$ term, without any basis change.
\begin{itemize}[leftmargin=2.6em, labelsep=0.4em, itemsep=2pt, topsep=3pt]
\item[\optlabel{A}] Both Statement I and Statement II are true
\item[\optlabel{B}] Both Statement I and Statement II are false
\item[\correctlabel{C}] Statement I is true, but Statement II is false
\item[\optlabel{D}] Statement I is false, but Statement II is true
\end{itemize}
\end{qitem}
\begin{solbox}{C}
\textbf{\color{brand}Correct answer.} \textcolor{okgreen}{\textbf{(C)}} Statement I is true, but Statement II is false

\textbf{\color{brand}Concepts required.} Unitarity preserves norm; measurement bases.

\textbf{\color{brand}Relevant formulae.} $\|U\ket\psi\| = \|\ket\psi\|$; measure $X$ via $H$.

\textbf{\color{brand}Step-by-step explanation.} \textbf{I is true:} gates are unitary, $\|U\ket\psi\| = \|\ket\psi\|$, so $\ket{\psi(\boldsymbol\theta)} = U(\boldsymbol\theta)\ket{0\dots0}$ has norm 1 for every $\boldsymbol\theta$. This is why
VQE can use $\braket{\psi|H|\psi}$ directly, without dividing by $\braket{\psi|\psi}$ as in a Rayleigh quotient.
\textbf{II is false:} a $Z$-basis measurement yields eigenvalues of $Z$; for $X_i$ a Hadamard on qubit $i$ is needed first (Q.24). Answer: (C).

\textbf{\color{brand}Why the other options are incorrect.}
\begin{itemize}[leftmargin=2.4em, itemsep=1pt, topsep=2pt]
\item[\textbf{(A)}] needs both statements true; II is false.
\item[\textbf{(B)}] needs both statements false; I is true.
\item[\textbf{(D)}] needs I false and II true; I is true, II false.
\end{itemize}
\textbf{\color{brand}Key takeaway.} Circuit states are automatically normalised; non-$Z$ terms need basis changes.
\end{solbox}
\qsep

\needspace{13\baselineskip}
\begin{qitem}{29}
\qmeta{Single-correct MCQ}{Commuting vs.\ qubit-wise commuting}
Consider the four exchange terms of the notebook's H$_2$ Hamiltonian, $\texttt{XXYY}$, $\texttt{XYYX}$, $\texttt{YXXY}$, $\texttt{YYXX}$. Which statement is
correct?
\begin{itemize}[leftmargin=2.6em, labelsep=0.4em, itemsep=2pt, topsep=3pt]
\item[\optlabel{A}] Every pair anticommutes
\item[\correctlabel{B}] All six pairs commute, but no pair is qubit-wise commuting
\item[\optlabel{C}] Only $\texttt{XXYY}$ and $\texttt{YYXX}$ commute
\item[\optlabel{D}] All six pairs are qubit-wise commuting, so one measurement setting suffices
\end{itemize}
\end{qitem}
\begin{solbox}{B}
\textbf{\color{brand}Correct answer.} \textcolor{okgreen}{\textbf{(B)}} All six pairs commute, but no pair is qubit-wise commuting

\textbf{\color{brand}Concepts required.} General versus qubit-wise commutation.

\textbf{\color{brand}Relevant formulae.} commute $\iff$ even number of anticommuting positions; QWC $\iff$ each position equal or $I$.

\textbf{\color{brand}Step-by-step explanation.} Two Pauli strings commute iff they differ by anticommuting single-qubit factors at an \emph{even} number of positions.
$\texttt{XXYY}$ vs $\texttt{XYYX}$: positions 1 and 3 differ ($X$/$Y$) --- two: commute. $\texttt{XXYY}$ vs $\texttt{YXXY}$: positions 0 and 2 --- commute.
$\texttt{XXYY}$ vs $\texttt{YYXX}$: all four positions --- commute. Likewise for the remaining three pairs.
\emph{Qubit-wise} commuting requires, at every position, equal letters or an $I$; every pair here has positions with $X$ against $Y$, so
no pair qualifies. Consequence: with simple qubit-wise grouping these 4 terms need 4 settings; with general commuting grouping (a
joint basis change containing CNOTs) they need only one (Q.54).

\textbf{\color{brand}Why the other options are incorrect.}
\begin{itemize}[leftmargin=2.4em, itemsep=1pt, topsep=2pt]
\item[\textbf{(A)}] Each pair differs at an even number of positions, so they commute.
\item[\textbf{(C)}] All six pairs commute, not only one.
\item[\textbf{(D)}] No pair is QWC: every pair has positions with $X$ against $Y$.
\end{itemize}
\textbf{\color{brand}Key takeaway.} Commuting $\ne$ qubit-wise commuting; the difference costs measurement settings.
\end{solbox}
\qsep

\needspace{13\baselineskip}
\begin{qitem}{30}
\qmeta{Single-correct MCQ}{\texttt{exp\_pauli} on eigenstates}
CUDA-Q's \texttt{exp\_pauli(theta, q, word)} applies $e^{i\theta P}$ for the Pauli string $P$ given by \texttt{word}. Applied with the word
$\texttt{"ZZ"}$ to the basis state $\ket{01}$, it gives
\par\smallskip\noindent\begin{tabularx}{\linewidth}{@{}*{2}{L}@{}}
\optlabel{A}\,$\cos\theta\ket{01} + i\sin\theta\ket{10}$ & \optlabel{B}\,$e^{+i\theta}\ket{01}$ \\[3pt]
\correctlabel{C}\,$e^{-i\theta}\ket{01}$ & \optlabel{D}\,$\ket{10}$
\end{tabularx}
\end{qitem}
\begin{solbox}{C}
\textbf{\color{brand}Correct answer.} \textcolor{okgreen}{\textbf{(C)}} $e^{-i\theta}\ket{01}$

\textbf{\color{brand}Concepts required.} Exponential of a Pauli string on its eigenstate.

\textbf{\color{brand}Relevant formulae.} $P\ket v = \lambda\ket v \Rightarrow e^{i\theta P}\ket v = e^{i\theta\lambda}\ket v$; $ZZ\ket{01} = -\ket{01}$.

\textbf{\color{brand}Step-by-step explanation.} $Z\otimes Z\ket{01} = (+1)(-1)\ket{01} = -\ket{01}$: an eigenstate with eigenvalue $-1$. On an eigenvector, $e^{i\theta P}$ acts as the number
$e^{i\theta\lambda} = e^{-i\theta}$. A diagonal Pauli string only adds phases; it never changes bit values. The sign convention ($e^{+i\theta P}$,
verified with CUDA-Q 0.16) matters for the notebook's UCC state $\cos\theta\ket{1100} - \sin\theta\ket{0011}$.

\textbf{\color{brand}Why the other options are incorrect.}
\begin{itemize}[leftmargin=2.4em, itemsep=1pt, topsep=2pt]
\item[\textbf{(A)}] A diagonal string cannot change bits.
\item[\textbf{(B)}] Uses eigenvalue $+1$; $\ket{01}$ has different bits.
\item[\textbf{(D)}] Bit flips need $X$/$Y$ factors.
\end{itemize}
\textbf{\color{brand}Key takeaway.} Diagonal Pauli exponentials only add phases.
\end{solbox}
\qsep

\needspace{13\baselineskip}
\begin{qitem}{31}
\qmeta{Single-correct MCQ}{$\texttt{XXXY}\ket{1100}$}
The four-qubit Pauli string $X_0X_1X_2Y_3$ (word $\texttt{"XXXY"}$) acting on $\ket{1100}$ gives
\par\smallskip\noindent\begin{tabularx}{\linewidth}{@{}*{4}{L}@{}}
\correctlabel{A}\,$i\ket{0011}$ & \optlabel{B}\,$\ket{0011}$ & \optlabel{C}\,$-i\ket{0011}$ & \optlabel{D}\,$i\ket{1100}$
\end{tabularx}
\end{qitem}
\begin{solbox}{A}
\textbf{\color{brand}Correct answer.} \textcolor{okgreen}{\textbf{(A)}} $i\ket{0011}$

\textbf{\color{brand}Concepts required.} Pauli strings on basis states.

\textbf{\color{brand}Relevant formulae.} $X\ket b = \ket{1-b}$; $Y\ket0 = i\ket1$, $Y\ket1 = -i\ket0$.

\textbf{\color{brand}Step-by-step explanation.} Qubit by qubit: $X\ket1 = \ket0$ (qubits 0, 1), $X\ket0 = \ket1$ (qubit 2), $Y\ket0 = i\ket1$ (qubit 3). Product: $i\ket{0011}$.
Hence $e^{i\theta P}\ket{1100} = \cos\theta\ket{1100} + i\sin\theta\,(i\ket{0011}) = \cos\theta\ket{1100} - \sin\theta\ket{0011}$: the notebook's
UCC state. The single Pauli string moves both electrons from orbital 0 (qubits 0, 1) to orbital 1 (qubits 2, 3) --- a double excitation.

\textbf{\color{brand}Why the other options are incorrect.}
\begin{itemize}[leftmargin=2.4em, itemsep=1pt, topsep=2pt]
\item[\textbf{(B)}] Drops the factor $i$ from $Y$.
\item[\textbf{(C)}] $-i$ comes from $Y\ket1$, but qubit 3 is in $\ket0$.
\item[\textbf{(D)}] $X$ flips bits; the state changes.
\end{itemize}
\textbf{\color{brand}Key takeaway.} $\texttt{XXXY}$ moves both electrons from orbital 0 to orbital 1.
\end{solbox}
\qsep

\needspace{13\baselineskip}
\begin{qitem}{32}
\qmeta{Assertion--Reason}{Real expectation values}
\textbf{Assertion (A):} For a Hermitian $H$, the expectation value $\braket{\psi|H|\psi}$ is a real number for every state $\ket\psi$.\\[1mm]
\textbf{Reason (R):} The identity $\braket{\psi|H|\psi} = \braket{\psi|H|\psi}^*$ holds only for normalised states.
\begin{itemize}[leftmargin=2.6em, labelsep=0.4em, itemsep=2pt, topsep=3pt]
\item[\optlabel{A}] Both Assertion (A) and Reason (R) are true, and R is the correct explanation of A
\item[\optlabel{B}] Both A and R are true, but R is \emph{not} the correct explanation of A
\item[\correctlabel{C}] A is true, but R is false
\item[\optlabel{D}] A is false, but R is true
\end{itemize}
\end{qitem}
\begin{solbox}{C}
\textbf{\color{brand}Correct answer.} \textcolor{okgreen}{\textbf{(C)}} A is true, but R is false

\textbf{\color{brand}Concepts required.} Hermiticity implies real expectation values.

\textbf{\color{brand}Relevant formulae.} $\braket{\psi|H|\psi}^* = \braket{\psi|H^\dagger|\psi} = \braket{\psi|H|\psi}$.

\textbf{\color{brand}Step-by-step explanation.} \textbf{A is true:} $\braket{\psi|H|\psi}^* = \braket{\psi|H^\dagger|\psi} = \braket{\psi|H|\psi}$, so it equals its own complex conjugate.
\textbf{R is false:} the identity uses only $H = H^\dagger$; it holds for every vector, normalised or not (the norm $\braket{\psi|\psi}$ is real anyway).
Real energies are why a VQE energy can be compared with $E_0$ at all. Answer: (C).

\textbf{\color{brand}Why the other options are incorrect.}
\begin{itemize}[leftmargin=2.4em, itemsep=1pt, topsep=2pt]
\item[\textbf{(A)}] needs A true, R true \emph{and} R explaining A; R is false.
\item[\textbf{(B)}] needs A and R both true but R not explaining A; R is false.
\item[\textbf{(D)}] needs A false and R true; A is true, R false.
\end{itemize}
\textbf{\color{brand}Key takeaway.} Real energies need only $H = H^\dagger$.
\end{solbox}
\qsep

\needspace{13\baselineskip}
\begin{qitem}{33}
\qmeta{Single-correct MCQ}{State-vector / density-matrix memory}
A complex fp64 amplitude needs 16 bytes. For LiH in STO-3G with Jordan--Wigner (12 qubits), the memory needed for the state vector and
for the density matrix ($2^n\times2^n$ amplitudes) is
\par\smallskip\noindent\begin{tabularx}{\linewidth}{@{}*{2}{L}@{}}
\optlabel{A}\,$192$ B and $2.3$ kB & \optlabel{B}\,$65.5$ kB and $131$ kB \\[3pt]
\optlabel{C}\,$268$ MB and $65.5$ kB & \correctlabel{D}\,$65.5$ kB and $268$ MB
\end{tabularx}
\end{qitem}
\begin{solbox}{D}
\textbf{\color{brand}Correct answer.} \textcolor{okgreen}{\textbf{(D)}} $65.5$ kB and $268$ MB

\textbf{\color{brand}Concepts required.} Memory of state vectors and density matrices.

\textbf{\color{brand}Relevant formulae.} state: $16\cdot2^n$ B; density matrix: $16\cdot4^n$ B.

\textbf{\color{brand}Step-by-step explanation.} State vector: $16\times2^{12} = 16\times4096 = 65\,536$ bytes $\approx 65.5$ kB. Density matrix: $16\times4^{12} = 16\times16\,777\,216 = 268\,435\,456$ bytes
$\approx 268$ MB. Density matrices (needed for exact noisy simulation, target \texttt{density-matrix-cpu}) square the cost; at the notebook's 15
qubits the density matrix would need $17.2$ GB versus $524$ kB for the state vector.

\textbf{\color{brand}Why the other options are incorrect.}
\begin{itemize}[leftmargin=2.4em, itemsep=1pt, topsep=2pt]
\item[\textbf{(A)}] $192 = 16\times12$ treats memory as linear in $n$.
\item[\textbf{(B)}] $131$ kB is twice the state vector, not $4^n$ scaling.
\item[\textbf{(C)}] Swaps the two values.
\end{itemize}
\textbf{\color{brand}Key takeaway.} Noisy (density-matrix) simulation squares the memory.
\end{solbox}
\qsep

\needspace{13\baselineskip}
\begin{qitem}{34}
\qmeta{Single-correct MCQ}{$\braket{Z_0Z_1}$ through a CNOT}
Two qubits start in $\ket{00}$; $R_Y(\theta_0)$ and $R_Y(\theta_1)$ are applied to qubits 0 and 1, followed by a CNOT (control 0, target 1). The
expectation value $\braket{Z_0Z_1}$ of the final state is
\par\smallskip\noindent\begin{tabularx}{\linewidth}{@{}*{4}{L}@{}}
\optlabel{A}\,$\cos(\theta_0 + \theta_1)$ & \correctlabel{B}\,$\cos\theta_1$ & \optlabel{C}\,$\cos\theta_0$ & \optlabel{D}\,$\cos\theta_0\cos\theta_1$
\end{tabularx}
\end{qitem}
\begin{solbox}{B}
\textbf{\color{brand}Correct answer.} \textcolor{okgreen}{\textbf{(B)}} $\cos\theta_1$

\textbf{\color{brand}Concepts required.} Heisenberg picture through a CNOT.

\textbf{\color{brand}Relevant formulae.} $\mathrm{CNOT}\,Z_0Z_1\,\mathrm{CNOT} = Z_1$; $\braket{Z}_{R_Y(\theta)\ket0} = \cos\theta$.

\textbf{\color{brand}Step-by-step explanation.} Heisenberg picture: $\braket{Z_0Z_1}_{\text{after}} = \braket{\mathrm{CNOT}\,Z_0Z_1\,\mathrm{CNOT}}_{\text{before}}$. CNOT conjugation maps $Z_1 \to Z_0Z_1$, hence
$Z_0Z_1 \to Z_0\cdot Z_0Z_1 = Z_1$. So $\braket{Z_0Z_1}_{\text{after}} = \braket{Z_1}_{\text{before}} = \cos\theta_1$ (Bloch vector of $R_Y(\theta_1)\ket0$).
Direct check with the bits: after the CNOT, $z_1' = z_0\oplus z_1$, so $z_0 \oplus z_1' = z_1$ and $(-1)^{z_0\oplus z_1'} = (-1)^{z_1}$. Before the CNOT the
answer would have been $\cos\theta_0\cos\theta_1$.

\textbf{\color{brand}Why the other options are incorrect.}
\begin{itemize}[leftmargin=2.4em, itemsep=1pt, topsep=2pt]
\item[\textbf{(A)}] Expectation values of product states multiply; they do not add angles.
\item[\textbf{(C)}] $Z_0Z_1$ maps to $Z_1$, not $Z_0$.
\item[\textbf{(D)}] $\cos\theta_0\cos\theta_1$ is the value \emph{before} the CNOT.
\end{itemize}
\textbf{\color{brand}Key takeaway.} Propagate the observable instead of the state for quick calculations.
\end{solbox}
\qsep

\needspace{13\baselineskip}
\begin{qitem}{35}
\qmeta{Single-correct MCQ}{Measuring $Y$}
Which unitary $U$ satisfies $UYU^\dagger = Z$, so that measuring $Z$ after applying $U$ measures $Y$?
\par\smallskip\noindent\begin{tabularx}{\linewidth}{@{}*{2}{L}@{}}
\optlabel{A}\,$U = S$ & \correctlabel{B}\,$U = HS^\dagger$ (first $S^\dagger$, then $H$) \\[3pt]
\optlabel{C}\,$U = SH$ (first $H$, then $S$) & \optlabel{D}\,$U = H$
\end{tabularx}
\end{qitem}
\begin{solbox}{B}
\textbf{\color{brand}Correct answer.} \textcolor{okgreen}{\textbf{(B)}} $U = HS^\dagger$ (first $S^\dagger$, then $H$)

\textbf{\color{brand}Concepts required.} Clifford conjugation of Paulis by $S$ and $H$.

\textbf{\color{brand}Relevant formulae.} $SXS^\dagger = Y$, $HXH = Z$, $HYH = -Y$.

\textbf{\color{brand}Step-by-step explanation.} With $S = \operatorname{diag}(1, i)$: $SXS^\dagger = \begin{psmallmatrix}1&0\\0&i\end{psmallmatrix}\begin{psmallmatrix}0&1\\1&0\end{psmallmatrix}\begin{psmallmatrix}1&0\\0&-i\end{psmallmatrix} = \begin{psmallmatrix}0&-i\\i&0\end{psmallmatrix} = Y$,
hence $S^\dagger YS = X$. Therefore
\[ (HS^\dagger)\,Y\,(HS^\dagger)^\dagger = H\,(S^\dagger YS)\,H = HXH = Z . \]
In circuit order: apply $S^\dagger$, then $H$, then measure in the computational basis. $H$ alone maps $Y \to -Y$; $S$ maps $Y \to SYS^\dagger = -X$;
neither diagonalises $Y$. (Verified numerically in \texttt{verify\_answers.py}.)

\textbf{\color{brand}Why the other options are incorrect.}
\begin{itemize}[leftmargin=2.4em, itemsep=1pt, topsep=2pt]
\item[\textbf{(A)}] $SYS^\dagger = -X$, not $Z$.
\item[\textbf{(C)}] $SH$ in that order maps $Y$ elsewhere ($SHYHS^\dagger = -SYS^\dagger = X$).
\item[\textbf{(D)}] $HYH = -Y$: still not diagonal.
\end{itemize}
\textbf{\color{brand}Key takeaway.} Measure $Y$: apply $S^\dagger$, then $H$.
\end{solbox}
\qsep

\needspace{20\baselineskip}
\begin{passagebox}{Paragraph for Questions 36 and 37}
A two-qubit transverse-field Ising Hamiltonian $H = -Z_0Z_1 - \tfrac12(X_0 + X_1)$ (the notebook's \texttt{"tfim"} model with $N = 2$,
$J = 1$, $h = 0.5$) is treated with a tiny hardware-efficient circuit: starting from $\ket{00}$, apply $R_Y(\theta_0)$ to qubit 0 and $R_Y(\theta_1)$
to qubit 1, then a CNOT with control 0 and target 1. The angles are $\theta_0 = \pi/2$, $\theta_1 = 0$.
\end{passagebox}
\needspace{13\baselineskip}
\begin{qitem}{36}
\qmeta{Paragraph based}{Two-qubit Ising: trial energy}
The energy $\braket{\psi|H|\psi}$ of the prepared state is
\par\smallskip\noindent\begin{tabularx}{\linewidth}{@{}*{4}{L}@{}}
\correctlabel{A}\,$-1$ & \optlabel{B}\,$0$ & \optlabel{C}\,$-\sqrt2 \approx -1.414$ & \optlabel{D}\,$-1.5$
\end{tabularx}
\end{qitem}
\begin{solbox}{A}
\textbf{\color{brand}Correct answer.} \textcolor{okgreen}{\textbf{(A)}} $-1$

\textbf{\color{brand}Concepts required.} Term-by-term energy of a Bell state.

\textbf{\color{brand}Relevant formulae.} $E = -\braket{Z_0Z_1} - \tfrac12(\braket{X_0} + \braket{X_1})$.

\textbf{\color{brand}Step-by-step explanation.} The circuit prepares the Bell state $\frac{1}{\sqrt2}(\ket{00} + \ket{11})$ (Q.27). Term by term: $\braket{Z_0Z_1} = 1$ (both components have equal bits);
$\braket{X_0} = 0$ because $X_0$ maps $\ket{00} \to \ket{10}$ and $\ket{11} \to \ket{01}$, states orthogonal to the Bell state; likewise $\braket{X_1} = 0$.
$E = -1\cdot1 - \tfrac12(0 + 0) = -1$. The state wins the full $ZZ$ energy but none of the field energy.

\textbf{\color{brand}Why the other options are incorrect.}
\begin{itemize}[leftmargin=2.4em, itemsep=1pt, topsep=2pt]
\item[\textbf{(B)}] $\braket{Z_0Z_1} = 1$ contributes $-1$.
\item[\textbf{(C)}] $-\sqrt2$ is the exact ground energy, not reached by this state.
\item[\textbf{(D)}] $-1.5$ assumes $\braket{X_0} = 1$ or similar; both are 0 for $\ket{\Phi^+}$.
\end{itemize}
\textbf{\color{brand}Key takeaway.} Entanglement alone is not enough: the ansatz must also capture the field terms.
\end{solbox}
\qsep

\needspace{13\baselineskip}
\begin{qitem}{37}
\qmeta{Paragraph based}{Two-qubit Ising: exact $E_0$}
The exact ground-state energy $E_0$ of this $H$, and the relative error of the prepared state, are
\par\smallskip\noindent\begin{tabularx}{\linewidth}{@{}*{2}{L}@{}}
\optlabel{A}\,$-1$; $0\,\%$ & \correctlabel{B}\,$-\sqrt2 \approx -1.414$; $29.3\,\%$ \\[3pt]
\optlabel{C}\,$-\sqrt2$; $41.4\,\%$ & \optlabel{D}\,$-1.5$; $33.3\,\%$
\end{tabularx}
\end{qitem}
\begin{solbox}{B}
\textbf{\color{brand}Correct answer.} \textcolor{okgreen}{\textbf{(B)}} $-\sqrt2 \approx -1.414$; $29.3\,\%$

\textbf{\color{brand}Concepts required.} Symmetry blocks; $2\times2$ diagonalisation; relative error.

\textbf{\color{brand}Relevant formulae.} $\lambda = \pm\sqrt{a^2 + b^2}$ for $\begin{psmallmatrix}-a&-b\\-b&a\end{psmallmatrix}$; rel.\ error $= (E - E_0)/|E_0|$.

\textbf{\color{brand}Step-by-step explanation.} $H$ commutes with $X_0X_1$, so it splits into two $2\times2$ blocks. The ground state lies in the block $X_0X_1 = +1$, spanned by
$\ket{\Phi^+} = \frac{1}{\sqrt2}(\ket{00} + \ket{11})$ and $\ket{\Psi^+} = \frac{1}{\sqrt2}(\ket{01} + \ket{10})$. Matrix elements: $\braket{\Phi^+|H|\Phi^+} = -1$ (from $-Z_0Z_1$),
$\braket{\Psi^+|H|\Psi^+} = +1$, and since $X_0\ket{\Psi^+} = X_1\ket{\Psi^+} = \ket{\Phi^+}$: $\braket{\Phi^+|H|\Psi^+} = -\tfrac12(1 + 1) = -1$. Eigenvalues of $\begin{psmallmatrix}-1&-1\\-1&1\end{psmallmatrix}$: $\pm\sqrt2$, so $E_0 = -\sqrt2$
(free-fermion formula and exact diagonalisation agree). Relative error $\frac{-1 - (-\sqrt2)}{\sqrt2} = \frac{0.414}{1.414} = 29.3\,\%$.
Choosing $\theta_1 \neq 0$ (or more layers) lets the ansatz mix in $\ket{\Psi^+}$ and lower the energy.

\textbf{\color{brand}Why the other options are incorrect.}
\begin{itemize}[leftmargin=2.4em, itemsep=1pt, topsep=2pt]
\item[\textbf{(A)}] $-1$ is the trial energy, not $E_0$.
\item[\textbf{(C)}] $41.4\,\%$ divides by $|E|$ instead of $|E_0|$.
\item[\textbf{(D)}] $-1.5$ ignores the off-diagonal coupling.
\end{itemize}
\textbf{\color{brand}Key takeaway.} Use symmetries to shrink an eigenproblem before solving it.
\end{solbox}
\qsep

\needspace{13\baselineskip}
\begin{qitem}{38}
\qmeta{Single-correct MCQ}{Pauli basis size}
Every Hermitian $2^n\times2^n$ matrix is a real linear combination of the $n$-qubit Pauli strings. For $n = 4$ the number of Pauli strings,
and the number actually needed for the notebook's H$_2$ Hamiltonian, are
\par\smallskip\noindent\begin{tabularx}{\linewidth}{@{}*{4}{L}@{}}
\optlabel{A}\,$256$ and $256$ & \optlabel{B}\,$64$ and $15$ & \optlabel{C}\,$16$ and $15$ & \correctlabel{D}\,$256$ and $15$
\end{tabularx}
\end{qitem}
\begin{solbox}{D}
\textbf{\color{brand}Correct answer.} \textcolor{okgreen}{\textbf{(D)}} $256$ and $15$

\textbf{\color{brand}Concepts required.} Pauli basis of operator space; sparsity of molecular Hamiltonians.

\textbf{\color{brand}Relevant formulae.} $\#\text{Pauli strings} = 4^n$; $c_P = 2^{-n}\operatorname{tr}(PH)$.

\textbf{\color{brand}Step-by-step explanation.} Each qubit carries one of $I, X, Y, Z$: $4^n$ strings; $n = 4$: $256$. They form an orthogonal basis of the $4^n$-dimensional space of
$2^n\times2^n$ matrices ($\operatorname{tr}(PQ) = 2^n\delta_{PQ}$), so $c_P = 2^{-n}\operatorname{tr}(PH)$. Molecular Hamiltonians are \emph{sparse} in this basis:
H$_2$ needs only 15 of the 256 (symmetries such as particle-number and spin conservation remove the rest).

\textbf{\color{brand}Why the other options are incorrect.}
\begin{itemize}[leftmargin=2.4em, itemsep=1pt, topsep=2pt]
\item[\textbf{(A)}] H$_2$ uses only 15 strings.
\item[\textbf{(B)}] $64 = 4^3$ would be the number of Pauli strings on 3 qubits, not 4.
\item[\textbf{(C)}] $16 = 2^4$ is the state dimension.
\end{itemize}
\textbf{\color{brand}Key takeaway.} $4^n$ possible terms, but physical Hamiltonians use few.
\end{solbox}
\qsep

\needspace{13\baselineskip}
\begin{qitem}{39}
\qmeta{Statement I / II}{Simultaneous measurement}
\textbf{Statement I:} Two Pauli strings either commute or anticommute.\\[1mm]
\textbf{Statement II:} A set of mutually commuting Pauli strings can be measured simultaneously after one common basis change.
\begin{itemize}[leftmargin=2.6em, labelsep=0.4em, itemsep=2pt, topsep=3pt]
\item[\correctlabel{A}] Both Statement I and Statement II are true
\item[\optlabel{B}] Both Statement I and Statement II are false
\item[\optlabel{C}] Statement I is true, but Statement II is false
\item[\optlabel{D}] Statement I is false, but Statement II is true
\end{itemize}
\end{qitem}
\begin{solbox}{A}
\textbf{\color{brand}Correct answer.} \textcolor{okgreen}{\textbf{(A)}} Both Statement I and Statement II are true

\textbf{\color{brand}Concepts required.} Pauli commutation; simultaneous measurement.

\textbf{\color{brand}Relevant formulae.} $PQ = \pm QP$; commuting observables share eigenbases.

\textbf{\color{brand}Step-by-step explanation.} \textbf{I is true:} position by position the factors commute or anticommute, so $PQ = (-1)^kQP$ with $k$ = number of anticommuting positions.
\textbf{II is true:} commuting Hermitian operators share an eigenbasis; a Clifford circuit rotating that basis to the computational basis lets
one set of shots estimate all of them (for qubit-wise commuting sets this circuit is just single-qubit rotations). This is the basis of
measurement grouping. Answer: (A).

\textbf{\color{brand}Why the other options are incorrect.}
\begin{itemize}[leftmargin=2.4em, itemsep=1pt, topsep=2pt]
\item[\textbf{(B)}] needs both statements false; both are true.
\item[\textbf{(C)}] needs I true and II false; II is true.
\item[\textbf{(D)}] needs I false and II true; I is true.
\end{itemize}
\textbf{\color{brand}Key takeaway.} Grouping commuting terms saves shots.
\end{solbox}
\qsep

\needspace{13\baselineskip}
\begin{qitem}{40}
\qmeta{Single-correct MCQ}{Energy variance}
For $H = -Z + \tfrac12X$ (one qubit), the energy variance $\mathrm{Var}(H) = \braket{H^2} - \braket H^2$ in the state $\ket0$, and in the exact ground state, are
\par\smallskip\noindent\begin{tabularx}{\linewidth}{@{}*{4}{L}@{}}
\optlabel{A}\,$0.25$ and $1.25$ & \optlabel{B}\,$1.25$ and $0$ & \correctlabel{C}\,$0.25$ and $0$ & \optlabel{D}\,$0$ and $0.25$
\end{tabularx}
\end{qitem}
\begin{solbox}{C}
\textbf{\color{brand}Correct answer.} \textcolor{okgreen}{\textbf{(C)}} $0.25$ and $0$

\textbf{\color{brand}Concepts required.} Energy variance; zero-variance principle.

\textbf{\color{brand}Relevant formulae.} $\mathrm{Var}(H) = \braket{H^2} - \braket H^2$; $(aZ + bX)^2 = (a^2 + b^2)I$.

\textbf{\color{brand}Step-by-step explanation.} $H^2 = Z^2 - \tfrac12(ZX + XZ) + \tfrac14X^2 = \tfrac54I$ (the cross terms cancel because $ZX = -XZ$). So $\braket{H^2} = 1.25$ in every state.
For $\ket0$: $\braket H = -1$, variance $1.25 - 1 = 0.25$. For the ground state: $\braket H = -\sqrt{5/4}$, $\braket H^2 = 1.25$, variance $0$.
\emph{Zero-variance principle}: exact eigenstates have zero energy variance, so the variance measures how far a VQE state is from an
eigenstate --- and how many shots an energy estimate needs (shot noise $\propto \sqrt{\mathrm{Var}/S}$).

\textbf{\color{brand}Why the other options are incorrect.}
\begin{itemize}[leftmargin=2.4em, itemsep=1pt, topsep=2pt]
\item[\textbf{(A)}] The exact eigenstate has zero variance.
\item[\textbf{(B)}] $1.25$ is $\braket{H^2}$, not the variance.
\item[\textbf{(D)}] Swaps the two states.
\end{itemize}
\textbf{\color{brand}Key takeaway.} Variance $\to 0$ signals an eigenstate and fewer shots needed.
\end{solbox}
\qsep

\newpage

\begin{sectionbanner}SECTION C \enspace$\cdot$\enspace VQE Fundamentals \hfill Q.41--60\end{sectionbanner}

\needspace{13\baselineskip}
\begin{qitem}{41}
\qmeta{Single-correct MCQ}{Energy below $E_0$ = bug}
A VQE run with \emph{exact} expectation values (no shots, no noise) reports $E(\boldsymbol\theta) = -15.10$ for the 15-qubit Ising chain whose exact ground
energy is $E_0 = -15.0825$. The most likely explanation is
\begin{itemize}[leftmargin=2.6em, labelsep=0.4em, itemsep=2pt, topsep=3pt]
\item[\optlabel{A}] the optimizer found a state better than the exact ground state
\item[\optlabel{B}] a deeper ansatz always produces energies below the exact value
\item[\correctlabel{C}] an error in the Hamiltonian, the reference energy or the state normalisation --- a correct calculation cannot go below $E_0$
\item[\optlabel{D}] shot noise pushed the estimate below $E_0$
\end{itemize}
\end{qitem}
\begin{solbox}{C}
\textbf{\color{brand}Correct answer.} \textcolor{okgreen}{\textbf{(C)}} an error in the Hamiltonian, the reference energy or the state normalisation --- a correct calculation cannot go below $E_0$

\textbf{\color{brand}Concepts required.} Variational bound as a correctness check.

\textbf{\color{brand}Relevant formulae.} $E(\boldsymbol\theta) = \braket{\psi(\boldsymbol\theta)|H|\psi(\boldsymbol\theta)} \ge E_0$ for normalised states.

\textbf{\color{brand}Step-by-step explanation.} Variational principle: for every normalised state $\braket{\psi|H|\psi} \ge E_0$ (Q.7). Exact expectation values of a normalised circuit state
(unitary gates, Q.28) therefore can never fall below $E_0$. A value below it signals a bug: a wrong Hamiltonian term or sign, a wrong
reference, or an expectation value of an unnormalised vector. With shots, single estimates \emph{can} fluctuate below $E_0$, but the
question excludes shots. This is why the notebook prints ``the VQE energy must stay $\ge E_0$''.

\textbf{\color{brand}Why the other options are incorrect.}
\begin{itemize}[leftmargin=2.4em, itemsep=1pt, topsep=2pt]
\item[\textbf{(A)}] No normalised state lies below $E_0$.
\item[\textbf{(B)}] More layers approach $E_0$ from above, never cross it.
\item[\textbf{(D)}] Shot noise is excluded (exact expectation values).
\end{itemize}
\textbf{\color{brand}Key takeaway.} An energy below the exact reference is a bug detector.
\end{solbox}
\qsep

\needspace{13\baselineskip}
\begin{qitem}{42}
\qmeta{Single-correct MCQ}{Rayleigh--Ritz bounds}
In the Rayleigh--Ritz method, $H$ is diagonalised inside an $m$-dimensional trial subspace, giving Ritz values $\varepsilon_1 \le \dots \le \varepsilon_m$. Which
statement is correct?
\begin{itemize}[leftmargin=2.6em, labelsep=0.4em, itemsep=2pt, topsep=3pt]
\item[\optlabel{A}] Only $\varepsilon_1$ is bounded; the other Ritz values can lie below the exact eigenvalues
\item[\optlabel{B}] The Ritz values are lower bounds to the exact eigenvalues
\item[\correctlabel{C}] $\varepsilon_k \ge E_{k-1}$ for every $k$: each Ritz value is an upper bound to the corresponding exact eigenvalue
\item[\optlabel{D}] The Ritz values are exact whenever $m \ge 2$
\end{itemize}
\end{qitem}
\begin{solbox}{C}
\textbf{\color{brand}Correct answer.} \textcolor{okgreen}{\textbf{(C)}} $\varepsilon_k \ge E_{k-1}$ for every $k$: each Ritz value is an upper bound to the corresponding exact eigenvalue

\textbf{\color{brand}Concepts required.} Rayleigh--Ritz / Hylleraas--Undheim--MacDonald theorem.

\textbf{\color{brand}Relevant formulae.} $\varepsilon_k \ge E_{k-1}$ (min--max theorem).

\textbf{\color{brand}Step-by-step explanation.} Write $P$ for the projector onto the trial space. By the Courant--Fischer min--max theorem the $k$-th eigenvalue of $PHP$ restricted to the
subspace is $\ge$ the $k$-th eigenvalue of $H$ (Hylleraas--Undheim--MacDonald theorem): $\varepsilon_k \ge E_{k-1}$ (labelling exact levels
$E_0 \le E_1 \le \dots$). For $k = 1$ this is the variational principle. Enlarging the subspace can only lower every Ritz value. The passage for
Q.59 and Q.60 is a Rayleigh--Ritz calculation with $m = 2$; there $\varepsilon_1$ happens to be exact because the ground state lies in the subspace.

\textbf{\color{brand}Why the other options are incorrect.}
\begin{itemize}[leftmargin=2.4em, itemsep=1pt, topsep=2pt]
\item[\textbf{(A)}] All Ritz values are upper bounds, not only the first.
\item[\textbf{(B)}] Upper, not lower, bounds.
\item[\textbf{(D)}] Exact only if the subspace contains the eigenvectors.
\end{itemize}
\textbf{\color{brand}Key takeaway.} Restricting the state space can only raise energies.
\end{solbox}
\qsep

\needspace{13\baselineskip}
\begin{qitem}{43}
\qmeta{Single-correct MCQ}{Hardware-efficient ansatz counts}
The notebook's hardware-efficient ansatz uses $L$ layers of [$R_Y$ on every qubit, CNOT ladder over neighbours] plus a final $R_Y$ layer. For
$N = 10$ qubits and $L = 3$, the numbers of parameters, CNOTs and gates are
\par\smallskip\noindent\begin{tabularx}{\linewidth}{@{}*{4}{L}@{}}
\optlabel{A}\,$40$, $30$ and $70$ & \optlabel{B}\,$30$, $27$ and $57$ & \optlabel{C}\,$40$, $9$ and $49$ & \correctlabel{D}\,$40$, $27$ and $67$
\end{tabularx}
\end{qitem}
\begin{solbox}{D}
\textbf{\color{brand}Correct answer.} \textcolor{okgreen}{\textbf{(D)}} $40$, $27$ and $67$

\textbf{\color{brand}Concepts required.} Hardware-efficient ansatz counts.

\textbf{\color{brand}Relevant formulae.} parameters $(L+1)N$; CNOTs $L(N-1)$; gates sum of both.

\textbf{\color{brand}Step-by-step explanation.} Parameters: one angle per $R_Y$: $(L + 1)N = 4\cdot10 = 40$. CNOTs: each ladder has $N - 1 = 9$ CNOTs, $L$ ladders: $27$.
Gates: $40 + 27 = 67$ (the notebook's estimate \texttt{num\_parameters + ANSATZ\_LAYERS * (n\_qubits - 1)}).
Default ($N = 15$, $L = 2$): $45$ parameters, $28$ CNOTs, $73$ gates --- exactly the recorded statistics \texttt{ry 45, cx 28}.

\textbf{\color{brand}Why the other options are incorrect.}
\begin{itemize}[leftmargin=2.4em, itemsep=1pt, topsep=2pt]
\item[\textbf{(A)}] Uses $N$ CNOTs per ladder (a closed ring); the notebook's ladder is open.
\item[\textbf{(B)}] Forgets the final rotation layer ($LN = 30$).
\item[\textbf{(C)}] $9$ CNOTs is a single ladder; with $L = 3$ layers there are three ladders, $27$ CNOTs.
\end{itemize}
\textbf{\color{brand}Key takeaway.} Default: 45 parameters, 28 CNOTs, 73 gates.
\end{solbox}
\qsep

\needspace{13\baselineskip}
\begin{qitem}{44}
\qmeta{Single-correct MCQ}{Hartree--Fock preparation}
The UCC kernel starts with \texttt{x(qubits[0]); x(qubits[1])}. In the Jordan--Wigner picture with spin-orbitals ordered (orbital 0 $\alpha$, orbital 0 $\beta$,
orbital 1 $\alpha$, orbital 1 $\beta$), this prepares
\begin{itemize}[leftmargin=2.6em, labelsep=0.4em, itemsep=2pt, topsep=3pt]
\item[\optlabel{A}] the state with both electrons in the antibonding orbital
\item[\optlabel{B}] an equal superposition of all 2-electron determinants
\item[\correctlabel{C}] the Hartree--Fock determinant $\ket{1100}$: both electrons in the lower (bonding) spatial orbital, one with each spin
\item[\optlabel{D}] a triplet state with both electrons spin-up
\end{itemize}
\end{qitem}
\begin{solbox}{C}
\textbf{\color{brand}Correct answer.} \textcolor{okgreen}{\textbf{(C)}} the Hartree--Fock determinant $\ket{1100}$: both electrons in the lower (bonding) spatial orbital, one with each spin

\textbf{\color{brand}Concepts required.} Jordan--Wigner occupation; Hartree--Fock reference.

\textbf{\color{brand}Relevant formulae.} $\ket{n_0n_1n_2n_3}$, $n_j = 1$ occupied; HF: lowest $N_e$ spin-orbitals filled.

\textbf{\color{brand}Step-by-step explanation.} Under Jordan--Wigner, qubit $j$ in $\ket1$ means spin-orbital $j$ is occupied. $X$ on qubits 0 and 1 turns $\ket{0000}$ into $\ket{1100}$: spin-orbitals
0 ($\sigma_g\alpha$) and 1 ($\sigma_g\beta$) occupied. This closed-shell determinant is the restricted Hartree--Fock state; its energy
$\braket{1100|H|1100} = -1.116759$ Ha is the HF energy (Q.48). VQE starts there ($\theta = 0$) and the double excitation adds the correlation.

\textbf{\color{brand}Why the other options are incorrect.}
\begin{itemize}[leftmargin=2.4em, itemsep=1pt, topsep=2pt]
\item[\textbf{(A)}] Antibonding occupation is $\ket{0011}$.
\item[\textbf{(B)}] Two $X$ gates create one basis state, not a superposition.
\item[\textbf{(D)}] Spin-orbitals 0 and 1 are $\alpha$ and $\beta$ of the same orbital: a singlet-type closed shell.
\end{itemize}
\textbf{\color{brand}Key takeaway.} Initial state matters: HF is the natural VQE starting point for molecules.
\end{solbox}
\qsep

\needspace{13\baselineskip}
\begin{qitem}{45}
\qmeta{Single-correct MCQ}{Two-state energy formula}
Let $\ket u$, $\ket v$ be orthonormal, $H_{uu} = a$, $H_{vv} = c$, $H_{uv} = H_{vu} = b$ (real). For the trial state $\ket{\psi(\theta)} = \cos\theta\ket u - \sin\theta\ket v$, the energy is
\par\smallskip\noindent\begin{tabularx}{\linewidth}{@{}*{2}{L}@{}}
\correctlabel{A}\,$a\cos^2\theta + c\sin^2\theta - b\sin2\theta$ & \optlabel{B}\,$a\cos^2\theta - c\sin^2\theta$ \\[3pt]
\optlabel{C}\,$a\cos\theta + c\sin\theta - b$ & \optlabel{D}\,$a\cos^2\theta + c\sin^2\theta + b\sin2\theta$
\end{tabularx}
\end{qitem}
\begin{solbox}{A}
\textbf{\color{brand}Correct answer.} \textcolor{okgreen}{\textbf{(A)}} $a\cos^2\theta + c\sin^2\theta - b\sin2\theta$

\textbf{\color{brand}Concepts required.} Energy of a two-state superposition.

\textbf{\color{brand}Relevant formulae.} $E = a\cos^2\theta + c\sin^2\theta - b\sin2\theta$; $\tan2\theta^* = \frac{2b}{c - a}$.

\textbf{\color{brand}Step-by-step explanation.} $E = \braket{\psi|H|\psi} = \cos^2\theta\,H_{uu} + \sin^2\theta\,H_{vv} - \sin\theta\cos\theta(H_{uv} + H_{vu}) = a\cos^2\theta + c\sin^2\theta - 2b\sin\theta\cos\theta$,
and $2\sin\theta\cos\theta = \sin2\theta$. The minus sign comes from the $-\sin\theta$ amplitude. Minimising: $E'(\theta) = (c - a)\sin2\theta - 2b\cos2\theta = 0$,
i.e.\ $\tan2\theta^* = \frac{2b}{c - a}$, and the minimum equals the lower eigenvalue of $\begin{psmallmatrix}a&b\\b&c\end{psmallmatrix}$.

\textbf{\color{brand}Why the other options are incorrect.}
\begin{itemize}[leftmargin=2.4em, itemsep=1pt, topsep=2pt]
\item[\textbf{(B)}] Drops the cross term and has a wrong sign.
\item[\textbf{(C)}] Amplitudes, not probabilities: energies need squares.
\item[\textbf{(D)}] $+b\sin2\theta$ is for the state $\cos\theta\ket u + \sin\theta\ket v$.
\end{itemize}
\textbf{\color{brand}Key takeaway.} Two-level problems reduce VQE to one angle.
\end{solbox}
\qsep

\needspace{13\baselineskip}
\begin{qitem}{46}
\qmeta{Matching lists}{VQE $\leftrightarrow$ code}
Match each VQE ingredient (List-I) with the code in \texttt{VQE.ipynb} that implements it (List-II).
\par\smallskip\noindent{\renewcommand{\arraystretch}{1.35}%
\begin{tabularx}{\linewidth}{@{}c>{\raggedright\arraybackslash}p{0.33\linewidth}@{\hspace{1.2em}}cL@{}}
\toprule
\multicolumn{2}{@{}l}{\textbf{List-I}} & \multicolumn{2}{l@{}}{\textbf{List-II}} \\ \midrule
\textbf{(P)} & Hartree--Fock reference state & \textbf{(1)} & \texttt{minimize(energy, initial\_theta, method="COBYLA", ...)} \\
\textbf{(Q)} & energy $E(\boldsymbol\theta) = \braket{\psi(\boldsymbol\theta)|H|\psi(\boldsymbol\theta)}$ & \textbf{(2)} & \texttt{coefficient * cudaq.SpinOperator.from\_word(word)}, summed \\
\textbf{(R)} & classical parameter update & \textbf{(3)} & \texttt{x(qubits[0]); x(qubits[1])} \\
\textbf{(S)} & Hamiltonian as a sum of Pauli strings & \textbf{(4)} & \texttt{cudaq.observe(kernel, hamiltonian, *args).expectation()} \\
\bottomrule
\end{tabularx}}
\par\smallskip\noindent\begin{tabularx}{\linewidth}{@{}*{2}{L}@{}}
\correctlabel{A}\,P$\to$3, Q$\to$4, R$\to$1, S$\to$2 & \optlabel{B}\,P$\to$3, Q$\to$4, R$\to$2, S$\to$1 \\[3pt]
\optlabel{C}\,P$\to$4, Q$\to$3, R$\to$1, S$\to$2 & \optlabel{D}\,P$\to$3, Q$\to$1, R$\to$4, S$\to$2
\end{tabularx}
\end{qitem}
\begin{solbox}{A}
\textbf{\color{brand}Correct answer.} \textcolor{okgreen}{\textbf{(A)}} P$\to$3, Q$\to$4, R$\to$1, S$\to$2

\textbf{\color{brand}Concepts required.} Mapping VQE to the notebook's code.

\textbf{\color{brand}Relevant formulae.} $E(\boldsymbol\theta) \leftarrow$ \texttt{observe}; $\boldsymbol\theta \leftarrow$ \texttt{minimize}; $H \leftarrow$ \texttt{SpinOperator.from\_word}.

\textbf{\color{brand}Step-by-step explanation.} \begin{itemize}\itemsep0pt
\item P: two $X$ gates occupy spin-orbitals 0 and 1 (3).
\item Q: \texttt{cudaq.observe} evaluates $\sum_kc_k\braket{P_k}$ for the ansatz kernel; \texttt{.expectation()} returns the number (4).
\item R: SciPy's \texttt{minimize} proposes new angles from the returned energies (1).
\item S: \texttt{build\_spin\_operator} turns each \texttt{(coefficient, word)} pair into a \texttt{SpinOperator} and adds them (2).
\end{itemize}

\textbf{\color{brand}Why the other options are incorrect.}
\begin{itemize}[leftmargin=2.4em, itemsep=1pt, topsep=2pt]
\item[\textbf{(B)}] R$\to$2, S$\to$1 calls the Hamiltonian builder the optimizer.
\item[\textbf{(C)}] P$\to$4: the HF state is prepared by gates, not computed by \texttt{observe}.
\item[\textbf{(D)}] Q$\to$1, R$\to$4 swaps the quantum and the classical step.
\end{itemize}
\textbf{\color{brand}Key takeaway.} Quantum: prepare and measure; classical: update angles.
\end{solbox}
\qsep

\needspace{13\baselineskip}
\begin{qitem}{47}
\qmeta{Single-correct MCQ}{Energy from Pauli averages}
$H = 0.5\,Z_0 - 0.3\,Z_1 + 0.2\,X_0X_1$. A trial state has $\braket{Z_0} = 0.6$, $\braket{Z_1} = -0.2$ and $\braket{X_0X_1} = 0.5$. Its energy is
\par\smallskip\noindent\begin{tabularx}{\linewidth}{@{}*{4}{L}@{}}
\optlabel{A}\,$0.34$ & \optlabel{B}\,$1.3$ & \correctlabel{C}\,$0.46$ & \optlabel{D}\,$0.26$
\end{tabularx}
\end{qitem}
\begin{solbox}{C}
\textbf{\color{brand}Correct answer.} \textcolor{okgreen}{\textbf{(C)}} $0.46$

\textbf{\color{brand}Concepts required.} Linearity of expectation values.

\textbf{\color{brand}Relevant formulae.} $E = \sum_kc_k\braket{P_k}$.

\textbf{\color{brand}Step-by-step explanation.} Linearity: $E = 0.5(0.6) - 0.3(-0.2) + 0.2(0.5) = 0.30 + 0.06 + 0.10 = 0.46$. Each $\braket{P_k}$ is estimated separately (or in commuting groups) and the
energy is the coefficient-weighted sum --- exactly what \texttt{cudaq.observe} does term by term.

\textbf{\color{brand}Why the other options are incorrect.}
\begin{itemize}[leftmargin=2.4em, itemsep=1pt, topsep=2pt]
\item[\textbf{(A)}] $0.34$ uses $-0.3\times0.2$ (sign error).
\item[\textbf{(B)}] $1.3$ sums $|c_k|$ only.
\item[\textbf{(D)}] $0.26$ drops the $XX$ term.
\end{itemize}
\textbf{\color{brand}Key takeaway.} Multiply each measured Pauli average by its coefficient and add.
\end{solbox}
\qsep

\needspace{13\baselineskip}
\begin{qitem}{48}
\qmeta{Single-correct MCQ}{HF energy from $Z$ terms}
Only the identity and the $Z$-type terms of the notebook's H$_2$ Hamiltonian contribute to $\braket{1100|H|1100}$. With coefficients
$c_{IIII} = -0.097066$, $c_{ZIII} = c_{IZII} = 0.171413$, $c_{IIZI} = c_{IIIZ} = -0.223432$, $c_{ZZII} = 0.168689$, $c_{ZIZI} = c_{IZIZ} = 0.120625$,
$c_{ZIIZ} = c_{IZZI} = 0.165928$, $c_{IIZZ} = 0.174413$, the Hartree--Fock energy is
\par\smallskip\noindent\begin{tabularx}{\linewidth}{@{}*{4}{L}@{}}
\optlabel{A}\,$-1.137284$ Ha & \optlabel{B}\,$-0.097066$ Ha & \optlabel{C}\,$+0.462618$ Ha & \correctlabel{D}\,$-1.116759$ Ha
\end{tabularx}
\end{qitem}
\begin{solbox}{D}
\textbf{\color{brand}Correct answer.} \textcolor{okgreen}{\textbf{(D)}} $-1.116759$ Ha

\textbf{\color{brand}Concepts required.} Diagonal elements from $Z$-type terms.

\textbf{\color{brand}Relevant formulae.} $Z_j\ket{n} = (1 - 2n_j)\ket n$; off-diagonal strings give 0 on the diagonal.

\textbf{\color{brand}Step-by-step explanation.} On $\ket{1100}$: $Z_0 = Z_1 = -1$ (occupied), $Z_2 = Z_3 = +1$. The $XY$-type terms flip bits and have zero diagonal element.
\[ E_{\mathrm{HF}} = -0.097066 - 2(0.171413) - 2(0.223432) + 0.168689 - 2(0.120625) - 2(0.165928) + 0.174413 . \]
$= -0.097066 - 0.342826 - 0.446863 + 0.168689 - 0.241250 - 0.331856 + 0.174413 = -1.116759$ Ha.
(Signs: $Z_0Z_1 = (+1)$, $Z_2Z_3 = (+1)$, mixed pairs $(-1)$.) $-1.137284$ is FCI; $+0.462618$ is $\braket{0011|H|0011}$.

\textbf{\color{brand}Why the other options are incorrect.}
\begin{itemize}[leftmargin=2.4em, itemsep=1pt, topsep=2pt]
\item[\textbf{(A)}] $-1.137284$ is FCI, below HF.
\item[\textbf{(B)}] Only the identity term.
\item[\textbf{(C)}] $+0.462618$ is $\braket{0011|H|0011}$.
\end{itemize}
\textbf{\color{brand}Key takeaway.} HF energy = expectation of $H$ in one determinant.
\end{solbox}
\qsep

\needspace{13\baselineskip}
\begin{qitem}{49}
\qmeta{Assertion--Reason}{Exactness of the UCC ansatz}
\textbf{Assertion (A):} The notebook's one-parameter UCC ansatz can reach the exact (FCI) ground energy of its H$_2$ Hamiltonian.\\[1mm]
\textbf{Reason (R):} The exact ground state of that Hamiltonian lies in $\mathrm{span}\{\ket{1100}, \ket{0011}\}$, and $\cos\theta\ket{1100} - \sin\theta\ket{0011}$ covers every
real normalised vector of that span (up to sign).
\begin{itemize}[leftmargin=2.6em, labelsep=0.4em, itemsep=2pt, topsep=3pt]
\item[\correctlabel{A}] Both Assertion (A) and Reason (R) are true, and R is the correct explanation of A
\item[\optlabel{B}] Both A and R are true, but R is \emph{not} the correct explanation of A
\item[\optlabel{C}] A is true, but R is false
\item[\optlabel{D}] A is false, but R is true
\end{itemize}
\end{qitem}
\begin{solbox}{A}
\textbf{\color{brand}Correct answer.} \textcolor{okgreen}{\textbf{(A)}} Both Assertion (A) and Reason (R) are true, and R is the correct explanation of A

\textbf{\color{brand}Concepts required.} Expressiveness of an ansatz; exact subspace.

\textbf{\color{brand}Relevant formulae.} $\ket{\psi_0} = 0.9937\ket{1100} - 0.1125\ket{0011}$; $\theta^* = 0.112783$.

\textbf{\color{brand}Step-by-step explanation.} \textbf{R is true:} exact diagonalisation of the $16\times16$ matrix gives a ground state with weights $0.98733$ on $\ket{1100}$ and $0.01267$ on $\ket{0011}$
and nothing else (particle-number, spin and spatial symmetry forbid single excitations and the other determinants). Any real unit vector
$(\cos\theta, -\sin\theta)$ in that plane is reached by some $\theta$.
\textbf{A follows:} at $\theta^* = 0.112783$ the energy equals $-1.137283834$ Ha to $10^{-15}$ (NumPy and CUDA-Q). R explains A. Answer: (A).

\textbf{\color{brand}Why the other options are incorrect.}
\begin{itemize}[leftmargin=2.4em, itemsep=1pt, topsep=2pt]
\item[\textbf{(B)}] needs A and R both true but R not explaining A; R is exactly why A holds.
\item[\textbf{(C)}] needs A true and R false; R is true.
\item[\textbf{(D)}] needs A false and R true; A is true.
\end{itemize}
\textbf{\color{brand}Key takeaway.} An ansatz is exact if it spans the subspace containing the ground state.
\end{solbox}
\qsep

\needspace{13\baselineskip}
\begin{qitem}{50}
\qmeta{Single-correct MCQ}{Energy of $\ket{0\dots0}$}
For the default Ising chain ($N = 15$, $J = 1$, $h = 0.5$, $H = -J\sum Z_iZ_{i+1} - h\sum X_i$), the energy of $\ket{00\dots0}$ --- close to where the
hardware-efficient ansatz starts, since its initial angles are drawn from $\mathcal N(0, 0.1^2)$ --- is
\par\smallskip\noindent\begin{tabularx}{\linewidth}{@{}*{4}{L}@{}}
\optlabel{A}\,$-15$ & \optlabel{B}\,$-21.5$ & \optlabel{C}\,$-7.5$ & \correctlabel{D}\,$-14$
\end{tabularx}
\end{qitem}
\begin{solbox}{D}
\textbf{\color{brand}Correct answer.} \textcolor{okgreen}{\textbf{(D)}} $-14$

\textbf{\color{brand}Concepts required.} Energy of a product state for the Ising chain.

\textbf{\color{brand}Relevant formulae.} $\braket{0\dots0|Z_iZ_{i+1}|0\dots0} = 1$, $\braket{X_i} = 0$.

\textbf{\color{brand}Step-by-step explanation.} On $\ket{0\dots0}$ every $Z_iZ_{i+1} = +1$: $14$ bonds give $-14$. $\braket{X_i} = 0$ for a computational basis state. So $E = -14$, compared with
$E_0 = -15.0825$: the field term contributes about $-1.08$ in the true ground state through small admixtures of flipped spins.
Small random angles start the optimizer close to this classical ordered state (not exactly at it, to break symmetries).

\textbf{\color{brand}Why the other options are incorrect.}
\begin{itemize}[leftmargin=2.4em, itemsep=1pt, topsep=2pt]
\item[\textbf{(A)}] $-15$ counts 15 bonds; an open chain has 14.
\item[\textbf{(B)}] $-21.5$ adds $-h$ per site, but $\braket{X_i} = 0$ on a basis state.
\item[\textbf{(C)}] $-7.5$ uses only the field.
\end{itemize}
\textbf{\color{brand}Key takeaway.} The ordered product state already captures most of the energy.
\end{solbox}
\qsep

\needspace{13\baselineskip}
\begin{qitem}{51}
\qmeta{Single-correct MCQ}{Gradient cost}
Each of the 45 angles of the default hardware-efficient ansatz enters exactly one $R_Y$ gate. The number of circuit evaluations for one exact
parameter-shift gradient, and for one forward finite-difference gradient (L-BFGS-B), are
\par\smallskip\noindent\begin{tabularx}{\linewidth}{@{}*{4}{L}@{}}
\optlabel{A}\,$90$ and $90$ & \correctlabel{B}\,$90$ and $46$ & \optlabel{C}\,$2$ and $46$ & \optlabel{D}\,$45$ and $46$
\end{tabularx}
\end{qitem}
\begin{solbox}{B}
\textbf{\color{brand}Correct answer.} \textcolor{okgreen}{\textbf{(B)}} $90$ and $46$

\textbf{\color{brand}Concepts required.} Gradient cost: parameter shift versus finite differences.

\textbf{\color{brand}Relevant formulae.} shift: $2n_p$ runs (one gate per parameter); forward differences: $n_p + 1$.

\textbf{\color{brand}Step-by-step explanation.} Parameter shift for a gate $e^{-i\theta Y/2}$: $\partial_\theta E = \tfrac12[E(\theta + \tfrac\pi2) - E(\theta - \tfrac\pi2)]$ --- 2 evaluations per parameter, exact: $2\times45 = 90$.
Forward differences: $E(\boldsymbol\theta)$ once plus one shifted evaluation per parameter: $45 + 1 = 46$ (the notebook's table: ``N\_params + 1
runs per gradient''), approximate and fragile under shot noise. COBYLA (default) uses no gradients at all.

\textbf{\color{brand}Why the other options are incorrect.}
\begin{itemize}[leftmargin=2.4em, itemsep=1pt, topsep=2pt]
\item[\textbf{(A)}] Finite differences need one extra run per parameter plus the base point.
\item[\textbf{(C)}] Two runs suffice only for one parameter.
\item[\textbf{(D)}] Parameter shift needs two runs per parameter.
\end{itemize}
\textbf{\color{brand}Key takeaway.} Gradients cost $O(n_p)$ circuit runs.
\end{solbox}
\qsep

\needspace{13\baselineskip}
\begin{qitem}{52}
\qmeta{Statement I / II}{COBYLA; best-of-history}
\textbf{Statement I:} COBYLA needs no gradient evaluations.\\[1mm]
\textbf{Statement II:} The notebook reports the energy of the optimizer's final point \texttt{opt.x} as the VQE energy.
\begin{itemize}[leftmargin=2.6em, labelsep=0.4em, itemsep=2pt, topsep=3pt]
\item[\optlabel{A}] Both Statement I and Statement II are true
\item[\optlabel{B}] Both Statement I and Statement II are false
\item[\correctlabel{C}] Statement I is true, but Statement II is false
\item[\optlabel{D}] Statement I is false, but Statement II is true
\end{itemize}
\end{qitem}
\begin{solbox}{C}
\textbf{\color{brand}Correct answer.} \textcolor{okgreen}{\textbf{(C)}} Statement I is true, but Statement II is false

\textbf{\color{brand}Concepts required.} Optimizer properties; best-of-history bookkeeping.

\textbf{\color{brand}Relevant formulae.} \texttt{best\_index = argmin(energy\_history)}.

\textbf{\color{brand}Step-by-step explanation.} \textbf{I is true:} COBYLA builds linear models from function values inside a trust region (initial radius \texttt{rhobeg = 0.3} in the notebook).
\textbf{II is false:} the code stores every evaluation and uses \texttt{best\_index = int(np.argmin(energy\_history))}; the reported energy and the sampled
state belong to the best evaluated angles, which need not be \texttt{opt.x}. Answer: (C).

\textbf{\color{brand}Why the other options are incorrect.}
\begin{itemize}[leftmargin=2.4em, itemsep=1pt, topsep=2pt]
\item[\textbf{(A)}] needs both statements true; II is false.
\item[\textbf{(B)}] needs both statements false; I is true.
\item[\textbf{(D)}] needs I false and II true; I true, II false.
\end{itemize}
\textbf{\color{brand}Key takeaway.} Report the best evaluated point, not the optimizer's last guess.
\end{solbox}
\qsep

\needspace{13\baselineskip}
\begin{qitem}{53}
\qmeta{Single-correct MCQ}{Chemical accuracy check}
Four VQE runs on the notebook's H$_2$ Hamiltonian (exact expectation values; FCI $= -1.137284$ Ha) report the energies below. Which one is a valid,
chemically accurate result?
\par\smallskip\noindent\begin{tabularx}{\linewidth}{@{}*{4}{L}@{}}
\optlabel{A}\,$-1.11676$ Ha & \correctlabel{B}\,$-1.13600$ Ha & \optlabel{C}\,$-1.13500$ Ha & \optlabel{D}\,$-1.13900$ Ha
\end{tabularx}
\end{qitem}
\begin{solbox}{B}
\textbf{\color{brand}Correct answer.} \textcolor{okgreen}{\textbf{(B)}} $-1.13600$ Ha

\textbf{\color{brand}Concepts required.} Validity and chemical accuracy.

\textbf{\color{brand}Relevant formulae.} $0 \le E - E_{\mathrm{FCI}} < 1.6$ mHa.

\textbf{\color{brand}Step-by-step explanation.} Errors $E - E_{\mathrm{FCI}}$: $-1.13600$: $+1.28$ mHa (above FCI, below $1.6$ mHa: chemically accurate, valid).
$-1.13500$: $+2.28$ mHa (valid but not chemically accurate). $-1.11676$: $+20.5$ mHa, the Hartree--Fock energy ($\theta = 0$).
$-1.13900$: $-1.72$ mHa, \emph{below} FCI, impossible with exact expectation values (Q.41).

\textbf{\color{brand}Why the other options are incorrect.}
\begin{itemize}[leftmargin=2.4em, itemsep=1pt, topsep=2pt]
\item[\textbf{(A)}] That is the HF energy: $20.5$ mHa error.
\item[\textbf{(C)}] $2.28$ mHa error: above threshold.
\item[\textbf{(D)}] Below FCI: violates the variational principle.
\end{itemize}
\textbf{\color{brand}Key takeaway.} Valid means $\ge E_0$; accurate means within 1.6 mHa.
\end{solbox}
\qsep

\needspace{13\baselineskip}
\begin{qitem}{54}
\qmeta{Single-correct MCQ}{Measurement grouping}
The notebook's H$_2$ Hamiltonian has the identity, 10 $Z$-type terms and 4 terms $\texttt{XXYY}$, $\texttt{XYYX}$, $\texttt{YXXY}$, $\texttt{YYXX}$. With
\emph{qubit-wise commuting} grouping, the minimum number of measurement settings is ..., and with general commuting grouping ...
\par\smallskip\noindent\begin{tabularx}{\linewidth}{@{}*{4}{L}@{}}
\correctlabel{A}\,$5$; $2$ & \optlabel{B}\,$15$; $5$ & \optlabel{C}\,$2$; $1$ & \optlabel{D}\,$4$; $1$
\end{tabularx}
\end{qitem}
\begin{solbox}{A}
\textbf{\color{brand}Correct answer.} \textcolor{okgreen}{\textbf{(A)}} $5$; $2$

\textbf{\color{brand}Concepts required.} Measurement grouping.

\textbf{\color{brand}Relevant formulae.} settings = number of groups (identity needs none).

\textbf{\color{brand}Step-by-step explanation.} The identity needs no measurement. All $Z$-type terms are qubit-wise commuting: one setting (plain $Z$ measurement). No two of the four
exchange terms are qubit-wise commuting (Q.29), so each needs its own setting: $1 + 4 = 5$. Under general commutation the four exchange
terms commute pairwise and form one group (measured with an entangling basis change), so $1 + 1 = 2$. (Whether the $Z$-group and the
exchange group could merge: $\texttt{ZIII}$ anticommutes with $\texttt{XXYY}$, so no.) Grouping directly reduces the shot cost of VQE.

\textbf{\color{brand}Why the other options are incorrect.}
\begin{itemize}[leftmargin=2.4em, itemsep=1pt, topsep=2pt]
\item[\textbf{(B)}] $15$ measures every term separately including the identity.
\item[\textbf{(C)}] $2$ is the general-commuting count; QWC needs 5.
\item[\textbf{(D)}] Ignores the $Z$ group or merges incompatible groups.
\end{itemize}
\textbf{\color{brand}Key takeaway.} Better grouping $\Rightarrow$ fewer settings $\Rightarrow$ fewer shots.
\end{solbox}
\qsep

\needspace{13\baselineskip}
\begin{qitem}{55}
\qmeta{Single-correct MCQ}{Recorded relative error}
In the previously recorded default run, $E_{\mathrm{VQE}} = -15.0533898284$ and the exact $E_0 = -15.0825179883$. The relative error is
\par\smallskip\noindent\begin{tabularx}{\linewidth}{@{}*{4}{L}@{}}
\optlabel{A}\,$1.93\,\%$ & \correctlabel{B}\,$0.193\,\%$ & \optlabel{C}\,$0.0291\,\%$ & \optlabel{D}\,$-0.193\,\%$
\end{tabularx}
\end{qitem}
\begin{solbox}{B}
\textbf{\color{brand}Correct answer.} \textcolor{okgreen}{\textbf{(B)}} $0.193\,\%$

\textbf{\color{brand}Concepts required.} Absolute and relative error.

\textbf{\color{brand}Relevant formulae.} rel.\ error $= (E_{\mathrm{VQE}} - E_0)/|E_0|$.

\textbf{\color{brand}Step-by-step explanation.} Absolute error $E_{\mathrm{VQE}} - E_0 = -15.0533898 + 15.0825180 = 0.0291282$ (positive, as the variational principle requires). Relative:
$0.0291282/15.0825180 = 0.00193 = 0.193\,\%$ --- the notebook prints ``2.913e-02 (relative 0.193\%)''. Note that TFIM energies are dimensionless
(in units of $J$); chemical accuracy does not apply to this model.

\textbf{\color{brand}Why the other options are incorrect.}
\begin{itemize}[leftmargin=2.4em, itemsep=1pt, topsep=2pt]
\item[\textbf{(A)}] $1.93\,\%$ is ten times too large: $0.0291/15.08 = 0.00193 = 0.193\,\%$.
\item[\textbf{(C)}] $0.0291$ is the absolute error read as a percentage.
\item[\textbf{(D)}] Negative would mean below $E_0$.
\end{itemize}
\textbf{\color{brand}Key takeaway.} Relative errors can hide what matters; compare with a baseline (Q.79).
\end{solbox}
\qsep

\needspace{13\baselineskip}
\begin{qitem}{56}
\qmeta{Single-correct MCQ}{Real ansatz for real $H$}
The notebook's hardware-efficient ansatz uses only $R_Y$ rotations and CNOTs, so it produces states with \emph{real} amplitudes. Why is this
adequate for its Ising and H$_2$ Hamiltonians?
\begin{itemize}[leftmargin=2.6em, labelsep=0.4em, itemsep=2pt, topsep=3pt]
\item[\optlabel{A}] Complex amplitudes cannot be measured, so they are never needed
\item[\correctlabel{B}] Both Hamiltonians are real symmetric matrices, so a ground state can always be chosen real
\item[\optlabel{C}] Real amplitudes always give lower energies than complex ones
\item[\optlabel{D}] CNOT gates generate all complex phases automatically
\end{itemize}
\end{qitem}
\begin{solbox}{B}
\textbf{\color{brand}Correct answer.} \textcolor{okgreen}{\textbf{(B)}} Both Hamiltonians are real symmetric matrices, so a ground state can always be chosen real

\textbf{\color{brand}Concepts required.} Real Hamiltonians and real ansätze.

\textbf{\color{brand}Relevant formulae.} $H = H^{\mathsf T}$ real $\Rightarrow$ $\mathrm{Re}\,v$ and $\mathrm{Im}\,v$ are eigenvectors.

\textbf{\color{brand}Step-by-step explanation.} $R_Y(\theta)$ and CNOT have real matrices, so they map real vectors to real vectors. The Ising Hamiltonian ($Z$ and $X$ only) is real; the H$_2$
Hamiltonian is real too ($\texttt{XXYY}$-type strings contain an even number of $Y$'s, each $Y$ contributes $\pm i$, so the products are real).
For a real symmetric $H$ the real and imaginary parts of an eigenvector are eigenvectors themselves, so a real ground state exists. For
Hamiltonians with genuinely complex matrix elements (e.g.\ magnetic fields) a real ansatz would fail.

\textbf{\color{brand}Why the other options are incorrect.}
\begin{itemize}[leftmargin=2.4em, itemsep=1pt, topsep=2pt]
\item[\textbf{(A)}] Complex phases do affect energies through interference.
\item[\textbf{(C)}] Real and complex ansätze give the same minimum for real $H$; neither is always lower.
\item[\textbf{(D)}] CNOT is a real permutation matrix.
\end{itemize}
\textbf{\color{brand}Key takeaway.} Match the ansatz's number field to the Hamiltonian's.
\end{solbox}
\qsep

\needspace{13\baselineskip}
\begin{qitem}{57}
\qmeta{Single-correct MCQ}{Symmetry breaking}
The recorded optimised state of the Ising chain gives $\ket{000\dots0}$ with probability $0.70$ but $\ket{111\dots1}$ (not in the top 8) with nearly zero
probability, although $H$ is symmetric under flipping all spins. The best explanation is
\begin{itemize}[leftmargin=2.6em, labelsep=0.4em, itemsep=2pt, topsep=3pt]
\item[\optlabel{A}] $\ket{111\dots1}$ has a higher energy than $\ket{000\dots0}$
\item[\optlabel{B}] the exact ground state contains only $\ket{000\dots0}$
\item[\optlabel{C}] the simulator has a bias towards $\ket0$
\item[\correctlabel{D}] for $h < J$ the two ordered configurations are nearly degenerate; VQE settled on one of them (symmetry breaking), which costs only an energy exponentially small in $N$
\end{itemize}
\end{qitem}
\begin{solbox}{D}
\textbf{\color{brand}Correct answer.} \textcolor{okgreen}{\textbf{(D)}} for $h < J$ the two ordered configurations are nearly degenerate; VQE settled on one of them (symmetry breaking), which costs only an energy exponentially small in $N$

\textbf{\color{brand}Concepts required.} Spontaneous symmetry breaking in finite VQE.

\textbf{\color{brand}Relevant formulae.} $[H, \prod_iX_i] = 0$; splitting $\propto (h/J)^N$ for $h < J$.

\textbf{\color{brand}Step-by-step explanation.} $\prod_iX_i$ commutes with $H$, so the exact ground state is symmetric: $P(0\dots0) = P(1\dots1)$. For $h < J$ it is close to the ``cat'' state
$\frac{1}{\sqrt2}(\ket{\Uparrow} + \ket{\Downarrow})$ (dressed by spin flips), whose symmetric and antisymmetric combinations are split only by an energy that
decreases exponentially with $N$. A shallow ansatz started near $\ket{0\dots0}$ cannot build the cat state cheaply and settles near one branch, as the
notebook's text explains; the energy cost is tiny compared with the recorded error $0.029$. Both configurations have the same energy.

\textbf{\color{brand}Why the other options are incorrect.}
\begin{itemize}[leftmargin=2.4em, itemsep=1pt, topsep=2pt]
\item[\textbf{(A)}] Both configurations have energy $-14$.
\item[\textbf{(B)}] The exact ground state is symmetric between the two configurations.
\item[\textbf{(C)}] No such bias: the state is set by the optimised angles.
\end{itemize}
\textbf{\color{brand}Key takeaway.} A broken-symmetry VQE state can still have an excellent energy.
\end{solbox}
\qsep

\needspace{13\baselineskip}
\begin{qitem}{58}
\qmeta{Matching lists}{Optimizers}
Match each optimizer offered by the notebook (List-I) with its working principle (List-II).
\par\smallskip\noindent{\renewcommand{\arraystretch}{1.35}%
\begin{tabularx}{\linewidth}{@{}c>{\raggedright\arraybackslash}p{0.33\linewidth}@{\hspace{1.2em}}cL@{}}
\toprule
\multicolumn{2}{@{}l}{\textbf{List-I}} & \multicolumn{2}{l@{}}{\textbf{List-II}} \\ \midrule
\textbf{(P)} & COBYLA & \textbf{(1)} & successive line searches along a set of conjugate directions, derivative-free \\
\textbf{(Q)} & Nelder--Mead & \textbf{(2)} & quasi-Newton method with a limited-memory Hessian estimate, using gradients (finite differences here) and box bounds \\
\textbf{(R)} & Powell & \textbf{(3)} & linear models of the objective inside a shrinking trust region, derivative-free \\
\textbf{(S)} & L-BFGS-B & \textbf{(4)} & a simplex of $n+1$ points updated by reflection, expansion and contraction \\
\bottomrule
\end{tabularx}}
\par\smallskip\noindent\begin{tabularx}{\linewidth}{@{}*{2}{L}@{}}
\correctlabel{A}\,P$\to$3, Q$\to$4, R$\to$1, S$\to$2 & \optlabel{B}\,P$\to$3, Q$\to$1, R$\to$4, S$\to$2 \\[3pt]
\optlabel{C}\,P$\to$4, Q$\to$3, R$\to$1, S$\to$2 & \optlabel{D}\,P$\to$3, Q$\to$4, R$\to$2, S$\to$1
\end{tabularx}
\end{qitem}
\begin{solbox}{A}
\textbf{\color{brand}Correct answer.} \textcolor{okgreen}{\textbf{(A)}} P$\to$3, Q$\to$4, R$\to$1, S$\to$2

\textbf{\color{brand}Concepts required.} Classical optimizers.

\textbf{\color{brand}Relevant formulae.} COBYLA: trust region; NM: simplex; Powell: line searches; L-BFGS-B: quasi-Newton.

\textbf{\color{brand}Step-by-step explanation.} \begin{itemize}\itemsep0pt
\item P: COBYLA = Constrained Optimization BY Linear Approximation, trust region (3) --- default, \texttt{maxiter = 300}, \texttt{rhobeg = 0.3}.
\item Q: Nelder--Mead simplex (4) --- \texttt{maxfev} budget.
\item R: Powell's conjugate-direction method (1).
\item S: L-BFGS-B (2) --- the only one needing gradients; without analytic gradients SciPy uses $n + 1$ evaluations per gradient.
\end{itemize}

\textbf{\color{brand}Why the other options are incorrect.}
\begin{itemize}[leftmargin=2.4em, itemsep=1pt, topsep=2pt]
\item[\textbf{(B)}] Q$\to$1, R$\to$4 swaps Nelder--Mead and Powell.
\item[\textbf{(C)}] P$\to$4, Q$\to$3 swaps COBYLA and Nelder--Mead.
\item[\textbf{(D)}] R$\to$2, S$\to$1 calls Powell gradient-based.
\end{itemize}
\textbf{\color{brand}Key takeaway.} Derivative-free methods dominate VQE with noisy energies.
\end{solbox}
\qsep

\needspace{20\baselineskip}
\begin{passagebox}{Paragraph for Questions 59 and 60}
In the notebook's H$_2$ example the UCC state $\ket{\psi(\theta)} = \cos\theta\ket{1100} - \sin\theta\ket{0011}$ only involves two basis states. Within
this two-dimensional space the Hamiltonian is the real $2\times2$ matrix
\[ \begin{pmatrix} H_{11} & H_{12}\\ H_{12} & H_{22}\end{pmatrix} = \begin{pmatrix} -1.116759 & 0.181210\\ 0.181210 & 0.462618\end{pmatrix}\ \text{Ha}, \]
with $H_{11} = \braket{1100|H|1100}$, $H_{22} = \braket{0011|H|0011}$, $H_{12} = \braket{1100|H|0011}$ (computed from the notebook's 15 Pauli terms), so that
$E(\theta) = H_{11}\cos^2\theta + H_{22}\sin^2\theta - H_{12}\sin2\theta$.
\end{passagebox}
\needspace{13\baselineskip}
\begin{qitem}{59}
\qmeta{Paragraph based}{H$_2$ two-level reduction: $E_0$}
The lowest eigenvalue of the $2\times2$ matrix --- the best energy the UCC ansatz can reach --- is
\par\smallskip\noindent\begin{tabularx}{\linewidth}{@{}*{4}{L}@{}}
\optlabel{A}\,$-1.116759$ Ha & \optlabel{B}\,$+0.483143$ Ha & \optlabel{C}\,$-0.654141$ Ha & \correctlabel{D}\,$-1.137284$ Ha
\end{tabularx}
\end{qitem}
\begin{solbox}{D}
\textbf{\color{brand}Correct answer.} \textcolor{okgreen}{\textbf{(D)}} $-1.137284$ Ha

\textbf{\color{brand}Concepts required.} Eigenvalues of a $2\times2$ symmetric matrix.

\textbf{\color{brand}Relevant formulae.} $\lambda_\pm = \frac{a+c}{2} \pm \sqrt{\left(\frac{c-a}{2}\right)^2 + b^2}$.

\textbf{\color{brand}Step-by-step explanation.} For $\begin{psmallmatrix}a&b\\b&c\end{psmallmatrix}$: $\lambda_\pm = \frac{a+c}{2} \pm \sqrt{\left(\frac{c-a}{2}\right)^2 + b^2}$.
$\frac{a+c}{2} = \frac{-1.116759 + 0.462618}{2} = -0.327071$; $\frac{c-a}{2} = 0.789689$; $\sqrt{0.789689^2 + 0.181210^2} = \sqrt{0.623608 + 0.032837} = 0.810213$.
$\lambda_- = -0.327071 - 0.810213 = -1.137284$ Ha $= E_{\mathrm{FCI}}$; $\lambda_+ = +0.483143$ (an excited state of the full Hamiltonian, too).
The coupling $H_{12}$ lowers the energy below $H_{11} = E_{\mathrm{HF}}$ by $20.5$ mHa, the correlation energy.

\textbf{\color{brand}Why the other options are incorrect.}
\begin{itemize}[leftmargin=2.4em, itemsep=1pt, topsep=2pt]
\item[\textbf{(A)}] That is $H_{11} = E_{\mathrm{HF}}$.
\item[\textbf{(B)}] The upper eigenvalue.
\item[\textbf{(C)}] $-0.654141 = H_{11} + H_{22}$ is the trace (sum of both eigenvalues), not the lower eigenvalue.
\end{itemize}
\textbf{\color{brand}Key takeaway.} Two-level reduction reproduces FCI exactly for H$_2$.
\end{solbox}
\qsep

\needspace{13\baselineskip}
\begin{qitem}{60}
\qmeta{Paragraph based}{H$_2$: optimal angle, weights}
The optimal angle $\theta^*$ and the probability of measuring $\ket{0011}$ in the optimal state are
\par\smallskip\noindent\begin{tabularx}{\linewidth}{@{}*{2}{L}@{}}
\optlabel{A}\,$\theta^* \approx 0.2295$ rad and $0.0518$ & \correctlabel{B}\,$\theta^* \approx 0.1128$ rad and $0.0127$ \\[3pt]
\optlabel{C}\,$\theta^* \approx -0.1128$ rad and $0.0127$ & \optlabel{D}\,$\theta^* \approx 0.1128$ rad and $0.1125$
\end{tabularx}
\end{qitem}
\begin{solbox}{B}
\textbf{\color{brand}Correct answer.} \textcolor{okgreen}{\textbf{(B)}} $\theta^* \approx 0.1128$ rad and $0.0127$

\textbf{\color{brand}Concepts required.} Optimal mixing angle and weights.

\textbf{\color{brand}Relevant formulae.} $\tan2\theta^* = \frac{2H_{12}}{H_{22} - H_{11}}$; $P(\ket{0011}) = \sin^2\theta^*$.

\textbf{\color{brand}Step-by-step explanation.} $E'(\theta) = (H_{22} - H_{11})\sin2\theta - 2H_{12}\cos2\theta = 0 \Rightarrow \tan2\theta^* = \frac{2H_{12}}{H_{22} - H_{11}} = \frac{0.362420}{1.579377} = 0.229471$,
$2\theta^* = 0.225566$, $\theta^* = 0.112783$ rad (the minimum, $E'' > 0$). Probability of $\ket{0011}$: $\sin^2\theta^* = 0.0127$ (amplitude $-0.1125$).
So the exact H$_2$ ground state is $98.7\,\%$ Hartree--Fock and $1.3\,\%$ doubly excited; this small admixture is worth $20.5$ mHa.
Negative $\theta$ would raise the energy; $0.2295$ is $\tan2\theta^*$, not $\theta^*$.

\textbf{\color{brand}Why the other options are incorrect.}
\begin{itemize}[leftmargin=2.4em, itemsep=1pt, topsep=2pt]
\item[\textbf{(A)}] Confuses $\tan2\theta^*$ with $\theta^*$.
\item[\textbf{(C)}] Negative $\theta$ raises the energy.
\item[\textbf{(D)}] $0.1125$ is the amplitude's modulus, not the probability.
\end{itemize}
\textbf{\color{brand}Key takeaway.} $1.3\,\%$ admixture $\to$ $20.5$ mHa correlation energy.
\end{solbox}
\qsep

\newpage

\begin{sectionbanner}SECTION D \enspace$\cdot$\enspace Advanced VQE and Implementation \hfill Q.61--80\end{sectionbanner}

\needspace{13\baselineskip}
\begin{qitem}{61}
\qmeta{Single-correct MCQ}{Jordan--Wigner $a_j^\dagger$}
With the Jordan--Wigner transformation $a_j^\dagger = \big(\prod_{k<j}Z_k\big)\,\tfrac12(X_j - iY_j)$ (qubit $\ket1$ = occupied), the creation operator $a_2^\dagger$ on
4 spin-orbitals becomes
\par\smallskip\noindent\begin{tabularx}{\linewidth}{@{}*{2}{L}@{}}
\optlabel{A}\,$\tfrac12\,Z_0Z_1(X_2 + iY_2)$ & \optlabel{B}\,$\tfrac12(X_2 - iY_2)$ \\[3pt]
\optlabel{C}\,$\tfrac12\,Z_3(X_2 - iY_2)$ & \correctlabel{D}\,$\tfrac12\,Z_0Z_1(X_2 - iY_2)$
\end{tabularx}
\end{qitem}
\begin{solbox}{D}
\textbf{\color{brand}Correct answer.} \textcolor{okgreen}{\textbf{(D)}} $\tfrac12\,Z_0Z_1(X_2 - iY_2)$

\textbf{\color{brand}Concepts required.} Jordan--Wigner transformation.

\textbf{\color{brand}Relevant formulae.} $a_j^\dagger = Z_0\cdots Z_{j-1}\tfrac12(X_j - iY_j)$.

\textbf{\color{brand}Step-by-step explanation.} $\tfrac12(X - iY) = \begin{psmallmatrix}0&0\\1&0\end{psmallmatrix} = \ket1\bra0$ fills an empty orbital. The string $Z_0Z_1$ counts the parity of the orbitals \emph{before}
$j = 2$ and supplies the sign $(-1)^{n_0 + n_1}$ that makes operators on different orbitals \emph{anti}commute. Result: $\tfrac12 Z_0Z_1X_2 - \tfrac i2Z_0Z_1Y_2$
(OpenFermion gives the same: \texttt{ZZXI}: $0.5$, \texttt{ZZYI}: $-0.5i$). Without the string (option 2) $\{a_0, a_2^\dagger\}$ would not vanish;
$\tfrac12(X + iY)$ is the annihilation operator.

\textbf{\color{brand}Why the other options are incorrect.}
\begin{itemize}[leftmargin=2.4em, itemsep=1pt, topsep=2pt]
\item[\textbf{(A)}] $\frac12(X + iY)$ annihilates.
\item[\textbf{(B)}] Without the $Z$ string fermionic anticommutation fails.
\item[\textbf{(C)}] The string runs over \emph{earlier} qubits, not later ones.
\end{itemize}
\textbf{\color{brand}Key takeaway.} JW = local occupation + non-local parity string.
\end{solbox}
\qsep

\needspace{13\baselineskip}
\begin{qitem}{62}
\qmeta{Single-correct MCQ}{JW hopping term}
Under Jordan--Wigner, the hopping term $a_0^\dagger a_1 + a_1^\dagger a_0$ between neighbouring spin-orbitals 0 and 1 becomes
\par\smallskip\noindent\begin{tabularx}{\linewidth}{@{}*{2}{L}@{}}
\optlabel{A}\,$X_0X_1$ & \correctlabel{B}\,$\tfrac12(X_0X_1 + Y_0Y_1)$ \\[3pt]
\optlabel{C}\,$\tfrac12(Z_0 + Z_1)$ & \optlabel{D}\,$\tfrac12(X_0X_1 - Y_0Y_1)$
\end{tabularx}
\end{qitem}
\begin{solbox}{B}
\textbf{\color{brand}Correct answer.} \textcolor{okgreen}{\textbf{(B)}} $\tfrac12(X_0X_1 + Y_0Y_1)$

\textbf{\color{brand}Concepts required.} Hopping terms under Jordan--Wigner.

\textbf{\color{brand}Relevant formulae.} $a_j^\dagger a_{j+1} + \text{h.c.} = \tfrac12(X_jX_{j+1} + Y_jY_{j+1})$.

\textbf{\color{brand}Step-by-step explanation.} $a_0^\dagger = \tfrac12(X_0 - iY_0)$ and $a_1 = Z_0\,\tfrac12(X_1 + iY_1)$ (the Jordan--Wigner string of orbital 1 is $Z_0$). Since
$(X - iY)Z = XZ - iYZ = -iY + X$, we get $(X_0 - iY_0)Z_0 = X_0 - iY_0$, so
\[ a_0^\dagger a_1 = \tfrac14(X_0 - iY_0)(X_1 + iY_1) = \tfrac14\big(X_0X_1 + Y_0Y_1 + iX_0Y_1 - iY_0X_1\big). \]
The Hermitian conjugate $a_1^\dagger a_0$ has the opposite imaginary parts, so the sum is $\tfrac12(X_0X_1 + Y_0Y_1)$. It swaps $\ket{01} \leftrightarrow \ket{10}$ (one electron
hops) and annihilates $\ket{00}$, $\ket{11}$. Verified with OpenFermion (\texttt{XX}: $0.5$, \texttt{YY}: $0.5$). $\tfrac12(X_0X_1 - Y_0Y_1)$ would create or destroy \emph{pairs}.

\textbf{\color{brand}Why the other options are incorrect.}
\begin{itemize}[leftmargin=2.4em, itemsep=1pt, topsep=2pt]
\item[\textbf{(A)}] $X_0X_1$ alone also flips $\ket{00} \leftrightarrow \ket{11}$, changing electron number.
\item[\textbf{(C)}] $Z$ terms are number operators, not hopping.
\item[\textbf{(D)}] $XX - YY$ is the pairing term $a^\dagger a^\dagger + aa$.
\end{itemize}
\textbf{\color{brand}Key takeaway.} Neighbouring hops are short; distant hops carry $Z$ strings.
\end{solbox}
\qsep

\needspace{13\baselineskip}
\begin{qitem}{63}
\qmeta{Single-correct MCQ}{JW vs.\ BK Pauli weight}
The Pauli weight (number of non-identity factors) of the image of $a_7^\dagger$ on 8 spin-orbitals is ... under Jordan--Wigner and ... under Bravyi--Kitaev.
\par\smallskip\noindent\begin{tabularx}{\linewidth}{@{}*{4}{L}@{}}
\correctlabel{A}\,$8$; $4$ & \optlabel{B}\,$8$; $8$ & \optlabel{C}\,$1$; $4$ & \optlabel{D}\,$4$; $8$
\end{tabularx}
\end{qitem}
\begin{solbox}{A}
\textbf{\color{brand}Correct answer.} \textcolor{okgreen}{\textbf{(A)}} $8$; $4$

\textbf{\color{brand}Concepts required.} Pauli weight of encodings.

\textbf{\color{brand}Relevant formulae.} JW: weight $j + 1$ for $a_j^\dagger$; BK: $O(\log_2N)$.

\textbf{\color{brand}Step-by-step explanation.} Jordan--Wigner stores occupations locally but the parity of all previous orbitals non-locally: $a_7^\dagger = Z_0\cdots Z_6\,\tfrac12(X_7 - iY_7)$, weight $7 + 1 = 8$ --- $O(N)$.
Bravyi--Kitaev stores partial sums of occupations in a binary-tree pattern so that both occupation and parity information are spread over only
$O(\log N)$ qubits; for $N = 8$ the strings of $a_7^\dagger$ have weight $\log_28 + 1 = 4$ (OpenFermion: maximum weight 4). Lower weight means
shorter circuits for $e^{i\theta P}$ on hardware. The notebook uses neither mapping explicitly; its H$_2$ coefficients are already Jordan--Wigner.

\textbf{\color{brand}Why the other options are incorrect.}
\begin{itemize}[leftmargin=2.4em, itemsep=1pt, topsep=2pt]
\item[\textbf{(B)}] BK reduces the weight.
\item[\textbf{(C)}] Weight 1 ignores the parity string.
\item[\textbf{(D)}] Reversed: Jordan--Wigner has the long string (weight 8), Bravyi--Kitaev the short one (weight 4).
\end{itemize}
\textbf{\color{brand}Key takeaway.} BK trades simplicity for logarithmic locality.
\end{solbox}
\qsep

\needspace{13\baselineskip}
\begin{qitem}{64}
\qmeta{Matching lists}{Fermion encodings}
Match each fermion-to-qubit encoding (List-I) with its characteristic property (List-II).
\par\smallskip\noindent{\renewcommand{\arraystretch}{1.35}%
\begin{tabularx}{\linewidth}{@{}c>{\raggedright\arraybackslash}p{0.33\linewidth}@{\hspace{1.2em}}cL@{}}
\toprule
\multicolumn{2}{@{}l}{\textbf{List-I}} & \multicolumn{2}{l@{}}{\textbf{List-II}} \\ \midrule
\textbf{(P)} & Jordan--Wigner & \textbf{(1)} & qubit $j$ stores the parity $n_0 \oplus\dots\oplus n_j$; total parity and $S_z$-parity sit on single qubits, so two qubits can be removed \\
\textbf{(Q)} & Bravyi--Kitaev & \textbf{(2)} & symmetries ($\mathbb Z_2$ stabilisers commuting with $H$) are used to remove qubits from any encoding \\
\textbf{(R)} & parity mapping & \textbf{(3)} & qubit $j$ stores the occupation $n_j$; strings of length $O(N)$ \\
\textbf{(S)} & qubit tapering & \textbf{(4)} & qubits store partial sums of occupations on a binary tree; strings of length $O(\log N)$ \\
\bottomrule
\end{tabularx}}
\par\smallskip\noindent\begin{tabularx}{\linewidth}{@{}*{2}{L}@{}}
\optlabel{A}\,P$\to$4, Q$\to$3, R$\to$1, S$\to$2 & \correctlabel{B}\,P$\to$3, Q$\to$4, R$\to$1, S$\to$2 \\[3pt]
\optlabel{C}\,P$\to$3, Q$\to$1, R$\to$4, S$\to$2 & \optlabel{D}\,P$\to$3, Q$\to$4, R$\to$2, S$\to$1
\end{tabularx}
\end{qitem}
\begin{solbox}{B}
\textbf{\color{brand}Correct answer.} \textcolor{okgreen}{\textbf{(B)}} P$\to$3, Q$\to$4, R$\to$1, S$\to$2

\textbf{\color{brand}Concepts required.} Fermion encodings and symmetry reductions.

\textbf{\color{brand}Relevant formulae.} JW: $q_j = n_j$; parity: $q_j = \bigoplus_{k\le j}n_k$; BK: tree partial sums.

\textbf{\color{brand}Step-by-step explanation.} \begin{itemize}\itemsep0pt
\item P: Jordan--Wigner: occupation basis, $Z$-strings (3).
\item Q: Bravyi--Kitaev: binary-tree partial sums, logarithmic weight (4).
\item R: parity mapping: cumulative parities; with interleaved/blocked spin ordering the electron-number parity and spin parity become single-qubit
$Z$'s, which allows the well-known reduction of H$_2$ from 4 to 2 qubits (1).
\item S: tapering (Bravyi et al.\ 2017) is a general symmetry reduction applied after any encoding (2).
\end{itemize}

\textbf{\color{brand}Why the other options are incorrect.}
\begin{itemize}[leftmargin=2.4em, itemsep=1pt, topsep=2pt]
\item[\textbf{(A)}] P$\to$4, Q$\to$3 swaps JW and BK.
\item[\textbf{(C)}] Q$\to$1, R$\to$4 swaps BK and parity.
\item[\textbf{(D)}] R$\to$2, S$\to$1 swaps parity mapping and tapering.
\end{itemize}
\textbf{\color{brand}Key takeaway.} Choice of encoding changes circuit cost, not physics.
\end{solbox}
\qsep

\needspace{13\baselineskip}
\begin{qitem}{65}
\qmeta{Single-correct MCQ}{Bravyi--Kitaev encoding}
For 4 modes, the Bravyi--Kitaev qubit values are $b = Bn \bmod 2$ with $B = \begin{psmallmatrix}1&0&0&0\\1&1&0&0\\0&0&1&0\\1&1&1&1\end{psmallmatrix}$ and $n$ the occupation vector.
The Hartree--Fock occupation $n = (1, 1, 0, 0)$ is encoded as
\par\smallskip\noindent\begin{tabularx}{\linewidth}{@{}*{4}{L}@{}}
\optlabel{A}\,$\ket{0100}$ & \optlabel{B}\,$\ket{1100}$ & \optlabel{C}\,$\ket{1101}$ & \correctlabel{D}\,$\ket{1000}$
\end{tabularx}
\end{qitem}
\begin{solbox}{D}
\textbf{\color{brand}Correct answer.} \textcolor{okgreen}{\textbf{(D)}} $\ket{1000}$

\textbf{\color{brand}Concepts required.} Bravyi--Kitaev encoding of occupations.

\textbf{\color{brand}Relevant formulae.} $b = Bn \bmod 2$.

\textbf{\color{brand}Step-by-step explanation.} $b_0 = n_0 = 1$; $b_1 = n_0 + n_1 = 2 \equiv 0$; $b_2 = n_2 = 0$; $b_3 = n_0 + n_1 + n_2 + n_3 = 2 \equiv 0$. So $\ket{1000}$.
Check with OpenFermion: the Bravyi--Kitaev Hamiltonian of H$_2$ has $\braket{1000|H_{\mathrm{BK}}|1000} = E_{\mathrm{HF}}$, while $\braket{1100|H_{\mathrm{BK}}|1100}$ is a different
energy. The same physical state has different bitstrings in different encodings --- an ansatz must start from the encoding-correct reference.

\textbf{\color{brand}Why the other options are incorrect.}
\begin{itemize}[leftmargin=2.4em, itemsep=1pt, topsep=2pt]
\item[\textbf{(A)}] $\ket{0100}$ has $b_0 = 0$, but the first BK qubit equals $n_0 = 1$.
\item[\textbf{(B)}] That is the JW bitstring.
\item[\textbf{(C)}] $b_3 = 1 + 1 + 0 + 0 = 0 \pmod 2$, not 1.
\end{itemize}
\textbf{\color{brand}Key takeaway.} Reference states must be re-encoded for each mapping.
\end{solbox}
\qsep

\needspace{13\baselineskip}
\begin{qitem}{66}
\qmeta{Single-correct MCQ}{UCCSD excitation count}
Counting excitations in spin-orbitals without spin restrictions, LiH in STO-3G (12 spin-orbitals, 4 electrons) has how many single and double
excitations from the Hartree--Fock determinant?
\par\smallskip\noindent\begin{tabularx}{\linewidth}{@{}*{2}{L}@{}}
\correctlabel{A}\,$32$ singles and $168$ doubles & \optlabel{B}\,$32$ singles and $1024$ doubles \\[3pt]
\optlabel{C}\,$8$ singles and $28$ doubles & \optlabel{D}\,$48$ singles and $168$ doubles
\end{tabularx}
\end{qitem}
\begin{solbox}{A}
\textbf{\color{brand}Correct answer.} \textcolor{okgreen}{\textbf{(A)}} $32$ singles and $168$ doubles

\textbf{\color{brand}Concepts required.} Counting UCCSD excitations.

\textbf{\color{brand}Relevant formulae.} singles $n_on_v$; doubles $\binom{n_o}{2}\binom{n_v}{2}$.

\textbf{\color{brand}Step-by-step explanation.} Occupied spin-orbitals $n_{\mathrm{occ}} = 4$, virtual $n_{\mathrm{virt}} = 12 - 4 = 8$. Singles: choose one occupied and one virtual: $4\times8 = 32$.
Doubles: choose two occupied and two virtual: $\binom42\binom82 = 6\times28 = 168$. UCCSD has one parameter per excitation (fewer if spin and
spatial symmetry are imposed). For H$_2$: $\binom22\binom22 = 1$ double --- the notebook's single parameter.

\textbf{\color{brand}Why the other options are incorrect.}
\begin{itemize}[leftmargin=2.4em, itemsep=1pt, topsep=2pt]
\item[\textbf{(B)}] $1024 = 32^2$ squares the singles.
\item[\textbf{(C)}] Spatial-orbital counts without spin.
\item[\textbf{(D)}] $48 = 4\times12$ includes occupied targets.
\end{itemize}
\textbf{\color{brand}Key takeaway.} UCCSD parameters grow like $n_o^2n_v^2$.
\end{solbox}
\qsep

\needspace{13\baselineskip}
\begin{qitem}{67}
\qmeta{Single-correct MCQ}{JW double excitation}
Under Jordan--Wigner, the anti-Hermitian double-excitation generator $a_3^\dagger a_2^\dagger a_1a_0 - a_0^\dagger a_1^\dagger a_2a_3$ becomes a sum of how many Pauli strings,
and what does the notebook's UCC kernel use?
\begin{itemize}[leftmargin=2.6em, labelsep=0.4em, itemsep=2pt, topsep=3pt]
\item[\optlabel{A}] 4 strings with coefficients $\pm\tfrac12$; the notebook keeps $\texttt{XXYY}$
\item[\optlabel{B}] 1 string; the notebook uses all of it
\item[\correctlabel{C}] 8 strings with coefficients $\pm i/8$; the notebook keeps the single string $\texttt{XXXY}$
\item[\optlabel{D}] 16 strings; the notebook keeps 8
\end{itemize}
\end{qitem}
\begin{solbox}{C}
\textbf{\color{brand}Correct answer.} \textcolor{okgreen}{\textbf{(C)}} 8 strings with coefficients $\pm i/8$; the notebook keeps the single string $\texttt{XXXY}$

\textbf{\color{brand}Concepts required.} JW form of excitation operators.

\textbf{\color{brand}Relevant formulae.} $a^\dagger a^\dagger aa - \text{h.c.} = \frac i8\sum_{8\text{ strings}}\pm P$.

\textbf{\color{brand}Step-by-step explanation.} Each of the four operators contributes $\tfrac12(X \mp iY)$, giving $2^4 = 16$ products; the Hermitian conjugate cancels half of them, leaving 8 strings
with an odd number of $Y$'s and coefficients $\pm i/8$ (OpenFermion: $\texttt{XXXY}, \texttt{XXYX}, \texttt{XYXX}, \texttt{YXXX}, \texttt{XYYY}, \dots$). Full UCC
would exponentiate all 8 (they commute, so $e^{\theta(T - T^\dagger)}$ factorises exactly). The notebook keeps only $e^{i\theta\,\texttt{XXXY}}$: starting from $\ket{1100}$ it
already rotates within $\mathrm{span}\{\ket{1100}, \ket{0011}\}$ (Q.31), which contains the exact ground state (Q.49), so nothing is lost for H$_2$.

\textbf{\color{brand}Why the other options are incorrect.}
\begin{itemize}[leftmargin=2.4em, itemsep=1pt, topsep=2pt]
\item[\textbf{(A)}] Coefficients are $\pm i/8$ and the strings have an odd number of $Y$'s; $\texttt{XXYY}$ is a Hamiltonian term.
\item[\textbf{(B)}] Four ladder operators give 16 products, 8 survive.
\item[\textbf{(D)}] Half of the 16 cancel against the conjugate.
\end{itemize}
\textbf{\color{brand}Key takeaway.} Problem-specific shortcuts are fine when they keep the exact subspace.
\end{solbox}
\qsep

\needspace{13\baselineskip}
\begin{qitem}{68}
\qmeta{Matching lists}{Obstacles and remedies}
Match each practical VQE difficulty (List-I) with a standard remedy (List-II).
\par\smallskip\noindent{\renewcommand{\arraystretch}{1.35}%
\begin{tabularx}{\linewidth}{@{}c>{\raggedright\arraybackslash}p{0.33\linewidth}@{\hspace{1.2em}}cL@{}}
\toprule
\multicolumn{2}{@{}l}{\textbf{List-I}} & \multicolumn{2}{l@{}}{\textbf{List-II}} \\ \midrule
\textbf{(P)} & barren plateaus (vanishing gradients for deep random circuits) & \textbf{(1)} & error mitigation, e.g.\ zero-noise extrapolation \\
\textbf{(Q)} & shot noise in energy estimates & \textbf{(2)} & chemically motivated ansatz, local cost functions, careful (e.g.\ Hartree--Fock) initialisation \\
\textbf{(R)} & gate errors on hardware & \textbf{(3)} & several restarts / good initial parameters \\
\textbf{(S)} & local minima of $E(\boldsymbol\theta)$ & \textbf{(4)} & measurement grouping and more shots for terms with large $|c_k|$ \\
\bottomrule
\end{tabularx}}
\par\smallskip\noindent\begin{tabularx}{\linewidth}{@{}*{2}{L}@{}}
\optlabel{A}\,P$\to$4, Q$\to$2, R$\to$1, S$\to$3 & \correctlabel{B}\,P$\to$2, Q$\to$4, R$\to$1, S$\to$3 \\[3pt]
\optlabel{C}\,P$\to$2, Q$\to$1, R$\to$4, S$\to$3 & \optlabel{D}\,P$\to$2, Q$\to$4, R$\to$3, S$\to$1
\end{tabularx}
\end{qitem}
\begin{solbox}{B}
\textbf{\color{brand}Correct answer.} \textcolor{okgreen}{\textbf{(B)}} P$\to$2, Q$\to$4, R$\to$1, S$\to$3

\textbf{\color{brand}Concepts required.} Practical VQE obstacles.

\textbf{\color{brand}Relevant formulae.} $\mathrm{Var}[\partial_\theta E] \sim 2^{-n}$ (barren plateau); SE $\propto 1/\sqrt S$.

\textbf{\color{brand}Step-by-step explanation.} \begin{itemize}\itemsep0pt
\item P: barren plateaus (McClean et al.\ 2018): gradient variance decays exponentially with $n$ for random deep ansätze; structured ansätze, local cost
functions and good initialisation help (2).
\item Q: shot noise $\propto \sqrt{\mathrm{Var}/S}$: group commuting terms, allocate shots by $|c_k|$ (4).
\item R: gate noise biases the energy upwards; mitigation such as zero-noise extrapolation (Q.72) (1).
\item S: local minima: multiple starts, Hartree--Fock or INTERP-like initialisation (3).
\end{itemize}

\textbf{\color{brand}Why the other options are incorrect.}
\begin{itemize}[leftmargin=2.4em, itemsep=1pt, topsep=2pt]
\item[\textbf{(A)}] P$\to$4, Q$\to$2 swaps barren plateaus and shot noise.
\item[\textbf{(C)}] Q$\to$1, R$\to$4 swaps shot noise and gate errors.
\item[\textbf{(D)}] R$\to$3, S$\to$1 swaps gate errors and local minima.
\end{itemize}
\textbf{\color{brand}Key takeaway.} Each obstacle has its own remedy; none is free.
\end{solbox}
\qsep

\needspace{13\baselineskip}
\begin{qitem}{69}
\qmeta{Single-correct MCQ}{Shots for chemical accuracy}
An energy estimator has a per-shot standard deviation of $0.1$ Ha. How many shots are needed for a standard error of $1.6$ mHa (chemical accuracy)?
\par\smallskip\noindent\begin{tabularx}{\linewidth}{@{}*{4}{L}@{}}
\optlabel{A}\,about $1.6\times10^{5}$ & \optlabel{B}\,about $63$ & \correctlabel{C}\,about $3.9\times10^3$ & \optlabel{D}\,about $2.6\times10^{-4}$
\end{tabularx}
\end{qitem}
\begin{solbox}{C}
\textbf{\color{brand}Correct answer.} \textcolor{okgreen}{\textbf{(C)}} about $3.9\times10^3$

\textbf{\color{brand}Concepts required.} Shot count for a target precision.

\textbf{\color{brand}Relevant formulae.} $S = (\sigma/\varepsilon)^2$.

\textbf{\color{brand}Step-by-step explanation.} $\mathrm{SE} = \sigma/\sqrt S \le 0.0016 \Rightarrow S \ge (0.1/0.0016)^2 = 62.5^2 = 3906$. $63$ forgets to square. Precision costs quadratically: halving
the error needs four times the shots, and every optimizer step needs such an estimate --- the main practical cost of VQE on hardware.

\textbf{\color{brand}Why the other options are incorrect.}
\begin{itemize}[leftmargin=2.4em, itemsep=1pt, topsep=2pt]
\item[\textbf{(A)}] $1.6\times10^5$ uses $\varepsilon = 0.25$ mHa.
\item[\textbf{(B)}] $63 = \sigma/\varepsilon$, not squared.
\item[\textbf{(D)}] $2.6\times10^{-4} = (\varepsilon/\sigma)^2$ inverts the ratio; a shot count must be large when $\varepsilon \ll \sigma$.
\end{itemize}
\textbf{\color{brand}Key takeaway.} Ten times more precision costs a hundred times more shots.
\end{solbox}
\qsep

\needspace{13\baselineskip}
\begin{qitem}{70}
\qmeta{Single-correct MCQ}{Fully depolarised energy}
Under complete depolarisation the state becomes $\rho = I/2^n$. For the notebook's 4-qubit H$_2$ Hamiltonian the measured energy would then be
\par\smallskip\noindent\begin{tabularx}{\linewidth}{@{}*{2}{L}@{}}
\correctlabel{A}\,$-0.0971$ Ha, the coefficient of $\texttt{IIII}$ & \optlabel{B}\,$-1.1168$ Ha \\[3pt]
\optlabel{C}\,$-1.1373$ Ha & \optlabel{D}\,$0$ Ha
\end{tabularx}
\end{qitem}
\begin{solbox}{A}
\textbf{\color{brand}Correct answer.} \textcolor{okgreen}{\textbf{(A)}} $-0.0971$ Ha, the coefficient of $\texttt{IIII}$

\textbf{\color{brand}Concepts required.} Maximally mixed state energy.

\textbf{\color{brand}Relevant formulae.} $\operatorname{tr}(H)/2^n = c_{I\dots I}$.

\textbf{\color{brand}Step-by-step explanation.} $E = \operatorname{tr}(\rho H) = 2^{-n}\operatorname{tr}H = \sum_kc_k\,2^{-n}\operatorname{tr}P_k = c_{IIII}$, since every non-identity Pauli string is traceless. Here $-0.0970663$ Ha.
Noise thus pulls the energy \emph{up} towards $c_{IIII}$, far above $E_{\mathrm{FCI}} = -1.1373$; partial depolarisation interpolates between the ideal
value and this limit.

\textbf{\color{brand}Why the other options are incorrect.}
\begin{itemize}[leftmargin=2.4em, itemsep=1pt, topsep=2pt]
\item[\textbf{(B)}] HF needs the pure state $\ket{1100}$.
\item[\textbf{(C)}] FCI needs the ideal ground state.
\item[\textbf{(D)}] The identity coefficient is not zero.
\end{itemize}
\textbf{\color{brand}Key takeaway.} Noise drags energies towards the identity coefficient.
\end{solbox}
\qsep

\needspace{13\baselineskip}
\begin{qitem}{71}
\qmeta{Assertion--Reason}{Variational bound vs.\ shot noise}
\textbf{Assertion (A):} With a finite number of shots, an estimated VQE energy can never fall below $E_0$.\\[1mm]
\textbf{Reason (R):} For every density matrix $\rho$ (pure or mixed, noisy or not), $\operatorname{tr}(\rho H) \ge E_0$.
\begin{itemize}[leftmargin=2.6em, labelsep=0.4em, itemsep=2pt, topsep=3pt]
\item[\optlabel{A}] Both Assertion (A) and Reason (R) are true, and R is the correct explanation of A
\item[\optlabel{B}] Both A and R are true, but R is \emph{not} the correct explanation of A
\item[\optlabel{C}] A is true, but R is false
\item[\correctlabel{D}] A is false, but R is true
\end{itemize}
\end{qitem}
\begin{solbox}{D}
\textbf{\color{brand}Correct answer.} \textcolor{okgreen}{\textbf{(D)}} A is false, but R is true

\textbf{\color{brand}Concepts required.} Variational principle for mixed states versus sampling noise.

\textbf{\color{brand}Relevant formulae.} $\operatorname{tr}(\rho H) \ge E_0$; estimator SE $\sigma/\sqrt S$.

\textbf{\color{brand}Step-by-step explanation.} \textbf{R is true:} $\rho = \sum_jp_j\ket{\psi_j}\bra{\psi_j}$ with $p_j \ge 0$, $\sum p_j = 1$, so $\operatorname{tr}(\rho H) = \sum_jp_j\braket{\psi_j|H|\psi_j} \ge E_0$: the variational principle survives noise.
\textbf{A is false:} a finite-shot \emph{estimate} fluctuates around $\operatorname{tr}(\rho H)$ with standard error $\sigma/\sqrt S$, and can fall below $E_0$ when the true value is
close to it. The notebook warns about this (``only shot noise can push single evaluations slightly below''). Answer: (D).

\textbf{\color{brand}Why the other options are incorrect.}
\begin{itemize}[leftmargin=2.4em, itemsep=1pt, topsep=2pt]
\item[\textbf{(A)}] needs A true, R true \emph{and} R explaining A; A is false.
\item[\textbf{(B)}] needs A and R both true but R not explaining A; A is false.
\item[\textbf{(C)}] needs A true and R false; A is false and R true.
\end{itemize}
\textbf{\color{brand}Key takeaway.} The true expectation obeys the bound; finite-shot estimates may not.
\end{solbox}
\qsep

\needspace{13\baselineskip}
\begin{qitem}{72}
\qmeta{Single-correct MCQ}{Zero-noise extrapolation}
Zero-noise extrapolation: an energy measured at noise scale factors $\lambda = 1$ and $\lambda = 3$ gives $-1.10$ Ha and $-1.00$ Ha. A linear extrapolation to
$\lambda = 0$ gives
\par\smallskip\noindent\begin{tabularx}{\linewidth}{@{}*{4}{L}@{}}
\optlabel{A}\,$-1.10$ Ha & \correctlabel{B}\,$-1.15$ Ha & \optlabel{C}\,$-1.05$ Ha & \optlabel{D}\,$-1.20$ Ha
\end{tabularx}
\end{qitem}
\begin{solbox}{B}
\textbf{\color{brand}Correct answer.} \textcolor{okgreen}{\textbf{(B)}} $-1.15$ Ha

\textbf{\color{brand}Concepts required.} Zero-noise extrapolation.

\textbf{\color{brand}Relevant formulae.} $E(0) = E(\lambda_1) - \lambda_1\frac{E(\lambda_2) - E(\lambda_1)}{\lambda_2 - \lambda_1}$.

\textbf{\color{brand}Step-by-step explanation.} Slope: $\frac{-1.00 - (-1.10)}{3 - 1} = 0.05$ Ha per unit $\lambda$. $E(0) = E(1) - 0.05 = -1.15$ Ha. Extrapolation removes the bias of noise but amplifies
statistical errors, and the extrapolated value carries \emph{no} variational guarantee: it can land below $E_0$.

\textbf{\color{brand}Why the other options are incorrect.}
\begin{itemize}[leftmargin=2.4em, itemsep=1pt, topsep=2pt]
\item[\textbf{(A)}] Value at $\lambda = 1$, not extrapolated.
\item[\textbf{(C)}] Wrong direction: noise raises the energy.
\item[\textbf{(D)}] Extrapolates by the full difference twice.
\end{itemize}
\textbf{\color{brand}Key takeaway.} Mitigated energies lose the variational guarantee.
\end{solbox}
\qsep

\needspace{13\baselineskip}
\begin{qitem}{73}
\qmeta{Single-correct MCQ}{GPU request without GPU}
\texttt{VQE.ipynb} is run with \texttt{BACKEND = "gpu"} on a machine where \texttt{cudaq.num\_available\_gpus()} returns 0. What happens?
\begin{itemize}[leftmargin=2.6em, labelsep=0.4em, itemsep=2pt, topsep=3pt]
\item[\optlabel{A}] The notebook runs on the GPU emulator \texttt{tensornet} instead
\item[\optlabel{B}] CUDA-Q silently falls back to \texttt{qpp-cpu}
\item[\correctlabel{C}] \texttt{resolve\_backend} raises a \texttt{RuntimeError} explaining that no NVIDIA GPU is visible; no GPU target is selected
\item[\optlabel{D}] The run starts and crashes later with an out-of-memory error
\end{itemize}
\end{qitem}
\begin{solbox}{C}
\textbf{\color{brand}Correct answer.} \textcolor{okgreen}{\textbf{(C)}} \texttt{resolve\_backend} raises a \texttt{RuntimeError} explaining that no NVIDIA GPU is visible; no GPU target is selected

\textbf{\color{brand}Concepts required.} Backend safety checks in the shared utilities.

\textbf{\color{brand}Relevant formulae.} \texttt{if backend in ("gpu", "multi-gpu") and num\_gpus == 0: raise RuntimeError}.

\textbf{\color{brand}Step-by-step explanation.} The shared utilities check \texttt{if backend in ("gpu", "multi-gpu") and num\_gpus == 0} and raise \texttt{RuntimeError("BACKEND='gpu' was requested but CUDA-Q
sees no NVIDIA GPU ...")}. GPU targets are never selected without a GPU because selecting \texttt{tensornet} on such a machine aborted the
Python process (README, known limitations). Only \texttt{BACKEND = "auto"} falls back to the CPU. GPU paths are implemented but were not executed.

\textbf{\color{brand}Why the other options are incorrect.}
\begin{itemize}[leftmargin=2.4em, itemsep=1pt, topsep=2pt]
\item[\textbf{(A)}] No GPU target is ever selected without a GPU.
\item[\textbf{(B)}] Only \texttt{"auto"} falls back to CPU.
\item[\textbf{(D)}] The check happens before any simulation.
\end{itemize}
\textbf{\color{brand}Key takeaway.} Explicit backend requests fail loudly, as they should.
\end{solbox}
\qsep

\needspace{13\baselineskip}
\begin{qitem}{74}
\qmeta{Statement I / II}{Exact references}
\textbf{Statement I:} For the Ising chain, the notebook computes its reference energy by sparse diagonalisation for every $N$.\\[1mm]
\textbf{Statement II:} For \texttt{"H2"} and \texttt{"custom"} Hamiltonians, an exact reference is computed only up to \texttt{REFERENCE\_MAX\_QUBITS = 18} qubits.
\begin{itemize}[leftmargin=2.6em, labelsep=0.4em, itemsep=2pt, topsep=3pt]
\item[\optlabel{A}] Both Statement I and Statement II are true
\item[\optlabel{B}] Both Statement I and Statement II are false
\item[\optlabel{C}] Statement I is true, but Statement II is false
\item[\correctlabel{D}] Statement I is false, but Statement II is true
\end{itemize}
\end{qitem}
\begin{solbox}{D}
\textbf{\color{brand}Correct answer.} \textcolor{okgreen}{\textbf{(D)}} Statement I is false, but Statement II is true

\textbf{\color{brand}Concepts required.} Exact references in the notebook.

\textbf{\color{brand}Relevant formulae.} TFIM: free fermions (any $N$); others: \texttt{eigsh} for $N \le 18$.

\textbf{\color{brand}Step-by-step explanation.} \textbf{I is false:} the Ising reference uses the free-fermion formula \texttt{tfim\_exact\_energy} for any $N$; sparse diagonalisation is only an
\emph{additional cross-check} for $N \le 16$ (it printed $-15.0825179883$, identical).
\textbf{II is true:} \texttt{elif n\_qubits <= REFERENCE\_MAX\_QUBITS: reference\_energy = exact\_ground\_energy\_sparse(...)}; above that, ``not computed''. Answer: (D).

\textbf{\color{brand}Why the other options are incorrect.}
\begin{itemize}[leftmargin=2.4em, itemsep=1pt, topsep=2pt]
\item[\textbf{(A)}] needs both statements true; I is false.
\item[\textbf{(B)}] needs both statements false; II is true.
\item[\textbf{(C)}] needs I true and II false; I is false and II true.
\end{itemize}
\textbf{\color{brand}Key takeaway.} Know where your reference comes from and where it stops.
\end{solbox}
\qsep

\needspace{13\baselineskip}
\begin{qitem}{75}
\qmeta{Single-correct MCQ}{Free-fermion matrix size}
The notebook's exact Ising reference maps the chain to free fermions and computes $E_0 = -\tfrac12\sum|\lambda(iA)|$ for a real antisymmetric matrix $A$. For
$N = 100$ qubits, $A$ has size
\par\smallskip\noindent\begin{tabularx}{\linewidth}{@{}*{4}{L}@{}}
\optlabel{A}\,$4\times4$ & \correctlabel{B}\,$200\times200$ & \optlabel{C}\,$100\times100$ & \optlabel{D}\,$2^{100}\times2^{100}$
\end{tabularx}
\end{qitem}
\begin{solbox}{B}
\textbf{\color{brand}Correct answer.} \textcolor{okgreen}{\textbf{(B)}} $200\times200$

\textbf{\color{brand}Concepts required.} Free-fermion solution of the Ising chain.

\textbf{\color{brand}Relevant formulae.} $2N$ Majorana modes; $E_0 = -\frac12\sum|\lambda(iA)|$.

\textbf{\color{brand}Step-by-step explanation.} Each spin becomes two Majorana operators ($a_j$, $b_j$), so $A$ is $2N\times2N$: $200\times200$, diagonalised in milliseconds. The Hamiltonian matrix
would be $2^{100}\times2^{100}$ --- impossible to store. Exactly solvable models give VQE benchmarks at any size; molecules generally do not.

\textbf{\color{brand}Why the other options are incorrect.}
\begin{itemize}[leftmargin=2.4em, itemsep=1pt, topsep=2pt]
\item[\textbf{(A)}] Constant size only for $N = 2$.
\item[\textbf{(C)}] Forgets the two Majoranas per site.
\item[\textbf{(D)}] That is the Hilbert-space dimension.
\end{itemize}
\textbf{\color{brand}Key takeaway.} Integrable models give exact large-$N$ benchmarks.
\end{solbox}
\qsep

\needspace{13\baselineskip}
\begin{qitem}{76}
\qmeta{Single-correct MCQ}{Recorded timing}
In the recorded run the quantum execution time was $69.10$ s for $300$ energy evaluations, while the notebook's pre-run estimate was $\sim102.9$ s.
The average time per evaluation, and the ratio estimate/measured cell time ($70.1$ s), are about
\par\smallskip\noindent\begin{tabularx}{\linewidth}{@{}*{2}{L}@{}}
\optlabel{A}\,$0.0032$ s and $1.5$ & \optlabel{B}\,$69$ s and $1.5$ \\[3pt]
\optlabel{C}\,$0.23$ s and $0.68$ & \correctlabel{D}\,$0.23$ s and $1.5$
\end{tabularx}
\end{qitem}
\begin{solbox}{D}
\textbf{\color{brand}Correct answer.} \textcolor{okgreen}{\textbf{(D)}} $0.23$ s and $1.5$

\textbf{\color{brand}Concepts required.} Timing analysis of a recorded run.

\textbf{\color{brand}Relevant formulae.} $t_{\mathrm{eval}} = t_{\mathrm{exec}}/N_{\mathrm{eval}}$.

\textbf{\color{brand}Step-by-step explanation.} $69.10/300 = 0.230$ s per \texttt{cudaq.observe} call (73-gate circuit, 29 terms, 15 qubits). $102.9/70.1 = 1.47$: the estimate was about $50\,\%$ high,
within the stated $\pm3\times$ (it assumed an extra state pass for the energy, ``x2 rough allowance''). $0.0032$ s is the time per \emph{gate}
($0.230/73$), not per evaluation; $0.68$ is the inverse ratio.

\textbf{\color{brand}Why the other options are incorrect.}
\begin{itemize}[leftmargin=2.4em, itemsep=1pt, topsep=2pt]
\item[\textbf{(A)}] $0.0032$ s is per gate.
\item[\textbf{(B)}] $69$ s is the total execution time of all 300 evaluations, not the time of one.
\item[\textbf{(C)}] $0.68$ is measured/estimate.
\end{itemize}
\textbf{\color{brand}Key takeaway.} Estimates are rough; measured timings are the benchmark.
\end{solbox}
\qsep

\needspace{13\baselineskip}
\begin{qitem}{77}
\qmeta{Single-correct MCQ}{Frozen core}
For LiH (STO-3G, 12 spin-orbitals under Jordan--Wigner), freezing the doubly occupied Li $1s$ core orbital removes how many qubits?
\par\smallskip\noindent\begin{tabularx}{\linewidth}{@{}*{2}{L}@{}}
\correctlabel{A}\,$2$, leaving $10$ qubits & \optlabel{B}\,$0$; frozen orbitals still need qubits \\[3pt]
\optlabel{C}\,$4$, leaving $8$ qubits & \optlabel{D}\,$1$, leaving $11$ qubits
\end{tabularx}
\end{qitem}
\begin{solbox}{A}
\textbf{\color{brand}Correct answer.} \textcolor{okgreen}{\textbf{(A)}} $2$, leaving $10$ qubits

\textbf{\color{brand}Concepts required.} Frozen core approximation.

\textbf{\color{brand}Relevant formulae.} qubits removed $= 2\times$ frozen spatial orbitals.

\textbf{\color{brand}Step-by-step explanation.} One spatial orbital = two spin-orbitals ($\alpha$, $\beta$) = two qubits under Jordan--Wigner. Freezing it (always occupied, absorbed into the constant and
one-body terms) leaves $12 - 2 = 10$ qubits and 2 active electrons. Further symmetry tapering and removing weakly occupied virtual orbitals reduce the
count further; every such reduction is an approximation choice except tapering, which is exact.

\textbf{\color{brand}Why the other options are incorrect.}
\begin{itemize}[leftmargin=2.4em, itemsep=1pt, topsep=2pt]
\item[\textbf{(B)}] Frozen orbitals are removed from the active space.
\item[\textbf{(C)}] Only one orbital is frozen.
\item[\textbf{(D)}] Spin doubles the count.
\end{itemize}
\textbf{\color{brand}Key takeaway.} Active spaces trade accuracy for qubits.
\end{solbox}
\qsep

\needspace{13\baselineskip}
\begin{qitem}{78}
\qmeta{Single-correct MCQ}{LiH Pauli terms}
Generated with PySCF and OpenFermion (STO-3G, Jordan--Wigner, bond length 1.595 \AA), the LiH qubit Hamiltonian has how many Pauli terms, compared with
15 for H$_2$?
\par\smallskip\noindent\begin{tabularx}{\linewidth}{@{}*{4}{L}@{}}
\optlabel{A}\,$4096$ & \optlabel{B}\,$144$ & \correctlabel{C}\,$631$ & \optlabel{D}\,$30$
\end{tabularx}
\end{qitem}
\begin{solbox}{C}
\textbf{\color{brand}Correct answer.} \textcolor{okgreen}{\textbf{(C)}} $631$

\textbf{\color{brand}Concepts required.} Growth of molecular Pauli sums.

\textbf{\color{brand}Relevant formulae.} $O(N^4)$ terms; LiH/STO-3G/JW: 631.

\textbf{\color{brand}Step-by-step explanation.} The number of terms grows like $N^4$ (Q.12): from $N = 4$ to $N = 12$ spin-orbitals the naive factor is $3^4 = 81$; symmetries remove many, and the
generated Hamiltonian has 631 Pauli terms (checked in \texttt{verify\_answers.py}). $4096 = 2^{12}$ is the dimension of the state, not a term count.
More terms mean more measurement settings and shots per energy evaluation, the dominant cost of molecular VQE.

\textbf{\color{brand}Why the other options are incorrect.}
\begin{itemize}[leftmargin=2.4em, itemsep=1pt, topsep=2pt]
\item[\textbf{(A)}] $2^{12}$ is the state dimension.
\item[\textbf{(B)}] $144 = 12^2$ counts only one-body terms.
\item[\textbf{(D)}] Far too few for 12 qubits with all two-electron integrals.
\end{itemize}
\textbf{\color{brand}Key takeaway.} Measurement cost, not qubit count, dominates molecular VQE.
\end{solbox}
\qsep

\needspace{20\baselineskip}
\begin{passagebox}{Paragraph for Questions 79 and 80}
\emph{Previously recorded default run of \texttt{VQE.ipynb}} (CUDA-Q 0.16.0, target \texttt{qpp-cpu}, 4 vCPU): Ising chain with $N = 15$, $J = 1$,
$h = 0.5$ (29 Pauli terms); hardware-efficient ansatz with 2 layers (45 parameters, 73 gates); COBYLA stopped with ``the objective function has
been evaluated MAXFUN times'' after 300 evaluations; $E_{\mathrm{VQE}} = -15.0533898284$, exact $E_0 = -15.0825179883$; quantum execution time
$69.10$ s; the optimised state gave $\ket{000\dots0}$ with probability $0.7005$. The product state $\ket{000\dots0}$ itself has energy $-14$.
\end{passagebox}
\needspace{13\baselineskip}
\begin{qitem}{79}
\qmeta{Paragraph based}{Recorded run: recovered fraction}
Measured from the product state $\ket{000\dots0}$ ($E = -14$), what fraction of the remaining energy gap to $E_0$ did VQE recover?
\par\smallskip\noindent\begin{tabularx}{\linewidth}{@{}*{4}{L}@{}}
\correctlabel{A}\,$97.3\,\%$ & \optlabel{B}\,$99.8\,\%$ & \optlabel{C}\,$70.1\,\%$ & \optlabel{D}\,$2.7\,\%$
\end{tabularx}
\end{qitem}
\begin{solbox}{A}
\textbf{\color{brand}Correct answer.} \textcolor{okgreen}{\textbf{(A)}} $97.3\,\%$

\textbf{\color{brand}Concepts required.} Baseline-normalised accuracy.

\textbf{\color{brand}Relevant formulae.} fraction $= \frac{E_{\mathrm{base}} - E_{\mathrm{VQE}}}{E_{\mathrm{base}} - E_0}$.

\textbf{\color{brand}Step-by-step explanation.} Gap from the product state to $E_0$: $-14 - (-15.0825180) = 1.0825180$. Recovered: $-14 - (-15.0533898) = 1.0533898$.
Fraction $1.0533898/1.0825180 = 0.973$. The relative error $0.193\,\%$ (Q.55) looks tiny because most of the energy ($-14$) comes from the trivial
ordered state; measured against this baseline the ansatz still misses $2.7\,\%$ of the quantum-fluctuation energy. $99.8\,\%$ is $1 - 0.193\,\%$.

\textbf{\color{brand}Why the other options are incorrect.}
\begin{itemize}[leftmargin=2.4em, itemsep=1pt, topsep=2pt]
\item[\textbf{(B)}] $1 - 0.193\,\%$ ignores the baseline.
\item[\textbf{(C)}] $0.70$ is the probability of $\ket{0\dots0}$.
\item[\textbf{(D)}] $2.7\,\%$ is the missing fraction.
\end{itemize}
\textbf{\color{brand}Key takeaway.} Measure progress against a trivial reference.
\end{solbox}
\qsep

\needspace{13\baselineskip}
\begin{qitem}{80}
\qmeta{Paragraph based}{Recorded run: next step}
Which change is most likely to reduce the remaining error of this run?
\begin{itemize}[leftmargin=2.6em, labelsep=0.4em, itemsep=2pt, topsep=3pt]
\item[\optlabel{A}] More final sampling shots
\item[\optlabel{B}] Switching from the free-fermion reference to sparse diagonalisation
\item[\correctlabel{C}] More evaluations (the budget was exhausted) and, if needed, more ansatz layers
\item[\optlabel{D}] Reducing \texttt{MAX\_EVALUATIONS} to stop earlier
\end{itemize}
\end{qitem}
\begin{solbox}{C}
\textbf{\color{brand}Correct answer.} \textcolor{okgreen}{\textbf{(C)}} More evaluations (the budget was exhausted) and, if needed, more ansatz layers

\textbf{\color{brand}Concepts required.} Diagnosing a VQE run.

\textbf{\color{brand}Relevant formulae.} budget exhausted $\Rightarrow$ not converged; expressiveness grows with $L$.

\textbf{\color{brand}Step-by-step explanation.} COBYLA stopped because of \texttt{MAXFUN}, not because of a convergence criterion, so the optimiser had not finished; a 2-layer ansatz may also lack
expressiveness (the notebook suggests more layers and evaluations, e.g.\ 4 layers and 800 evaluations for the harder critical point). The final
sampling only displays bitstrings; it does not change the energy. Both references are exact and agree to $10^{-10}$. Fewer evaluations can only
make the error larger or equal.

\textbf{\color{brand}Why the other options are incorrect.}
\begin{itemize}[leftmargin=2.4em, itemsep=1pt, topsep=2pt]
\item[\textbf{(A)}] Sampling shots do not change the optimised energy.
\item[\textbf{(B)}] Both references are exact and agree.
\item[\textbf{(D)}] Fewer evaluations cannot help.
\end{itemize}
\textbf{\color{brand}Key takeaway.} Read the optimizer's stopping message before trusting the energy.
\end{solbox}
\qsep

\end{document}